\chapter{Algebraic number theory}\label{ch:ant}

The paper lives over $\Q$, but its arguments constantly leave $\Q$: the coefficients of a congruent eigenform lie in a ring of integers $\mathcal O$ with a prime
$\lambda\mid\ell$; the Galois group acts through decomposition and inertia groups and Frobenius elements; the quadratic fields $\Q(i)$, $\Q(\sqrt{-3})$,
$\Q(\sqrt{q^*})$ and $\Q(\sqrt{\ell^*})$ supply the characters $\chi_{-4}$, $\chi_{-3}$, $\chi_{q^*}$ and $\chi_{\ell^*}$; and the rings $\Z_3[\zeta_3]$ and
$\Z_\ell\otimes\Z[\zeta_3]$ carry the action of inertia at the additive primes of types $\kod{IV}$ and $\kod{II}$. This chapter supplies that algebraic number
theory, from rings of integers and unique factorisation of ideals to the splitting of primes, the Frobenius element, cyclotomic fields, the dictionary
between Dirichlet characters and characters of $G_\Q$, the Chebotarev density theorem and complex multiplication. It uses \Cref{ch:groups} for group actions
and profinite groups, \Cref{ch:rings} for rings, finite fields, Kronecker symbols, quadratic reciprocity and Dirichlet characters, \Cref{ch:modules} for modules
over principal ideal domains and discrete valuation rings, \Cref{ch:reps} for representations, and \Cref{ch:galois} for Galois theory; local fields are in
\Cref{ch:local}.

\section*{Core results at a glance}

Throughout, $K$, $L$ are number fields, $\mathcal O_K$ the ring of integers, $\mathfrak p$, $\mathfrak P$, $\lambda$ nonzero primes, $N\mathfrak p=\#\mathcal O_K/\mathfrak p$,
$D$ the discriminant of a quadratic field and $\chi_D$ its character; $\pi=1-\zeta_3$; $\Frob$ is the arithmetic Frobenius and $\Frob^{-1}$ the geometric one.

\begin{enumerate}[label=\textbf{\thechapter.\arabic*}]
\item \textbf{Conventions.} $\Qbar\subset\C$, $\bar\Z$ the algebraic integers, $\zeta_n=e^{2\pi i/n}$; $D$ is a quadratic discriminant and $D_{\mathfrak P}$ a decomposition
group; $\chi_D$ is indexed by the fundamental discriminant ($\chi_{-4}$, $\chi_{-3}$, $\chi_8$, $\chi_{q^*}$); $\chi_\ell$ is cyclotomic; the geometric Frobenius is written
$\Frob_p^{-1}$. \Cref{conv:ant:conventions}.
\item \textbf{Number fields} (definition): $[K:\Q]=n<\infty$, $K=\Q(\theta)$, $n$ embeddings into $\C$, $n=r_1+2r_2$. \Cref{def:ant:numberfield}.
\item \textbf{Algebraic integers.} $\alpha$ is integral iff its minimal polynomial is in $\Z[X]$ iff $\Z[\alpha]$ is a finitely generated $\Z$-module iff $\alpha$ lies in
a subring finitely generated as a $\Z$-module; $\mathcal O_K$ is a ring, $\mathcal O_K\cap\Q=\Z$, $\Frac\mathcal O_K=K$, and $\mathcal O_K$ is integrally closed.
\emph{(proved here)} \Cref{prop:ant:integral}.
\item \textbf{Trace and norm} (definition): trace and determinant of $x\mapsto\alpha x$. \Cref{def:ant:tracenorm}.
\item \textbf{Properties of trace and norm.} $\Tr(\alpha)=\sum\sigma_i(\alpha)$, $N(\alpha)=\prod\sigma_i(\alpha)$; integers on $\mathcal O_K$; $\alpha\in\mathcal O_K^\times$ iff
$N(\alpha)=\pm1$; transitivity in towers. \emph{(proved here)} \Cref{prop:ant:tracenorm}.
\item \textbf{Integral bases and the discriminant.} The trace form is nondegenerate; $\mathcal O_K\simeq\Z^n$; $d_K=\det(\Tr(\beta_i\beta_j))=\det(\sigma_k\beta_j)^2$ is
independent of the basis; $\disc(\alpha_1,\dots,\alpha_n)=[\mathcal O_K:\sum\Z\alpha_i]^2d_K$; $\disc(1,\theta,\dots,\theta^{n-1})=\disc(g)=(-1)^{n(n-1)/2}N(g'(\theta))$.
\emph{(proved here)} \Cref{thm:ant:integralbasis}.
\item \textbf{Quadratic fields.} For $d$ squarefree, $\mathcal O=\Z[\sqrt d]$, $d_K=4d$ if $d\equiv2,3\pmod4$, and $\mathcal O=\Z[\frac{1+\sqrt d}2]$, $d_K=d$ if
$d\equiv1\pmod4$; quadratic fields correspond to fundamental discriminants; $\Z[i]$ ($D=-4$), $\Z[\zeta_3]$ ($D=-3$). \emph{(proved here)}
\Cref{prop:ant:quadraticintegers}.
\item \textbf{Units.} (i) $\mathcal O_K^\times\simeq\mu(K)\times\Z^{r_1+r_2-1}$ \emph{(stated; reference)}; (ii) for $K$ imaginary quadratic $\mathcal O_K^\times=\mu(K)$, which is
$\{\pm1\}$ except $\mu_4$ for $\Q(i)$ and $\mu_6$ for $\Q(\sqrt{-3})$; (iii) $\Z[\sqrt2]^\times=\pm(1+\sqrt2)^\Z$ \emph{((ii), (iii) proved here)}. \Cref{thm:ant:units}.
\item \textbf{Dedekind domains} (definition): Noetherian, integrally closed, not a field, nonzero primes maximal. \Cref{def:ant:dedekind}.
\item \textbf{$\mathcal O_K$ is a Dedekind domain}; $\mathfrak p\cap\Z=p\Z$ and $\mathcal O_K/\mathfrak p$ is a finite field; PIDs (not fields) are Dedekind; a Noetherian domain
(not a field) is Dedekind iff all its localisations at nonzero primes are discrete valuation rings. \emph{(proved here)} \Cref{thm:ant:dedekind}.
\item \textbf{Inverse of a prime.} Every nonzero ideal contains a product of primes; $\mathfrak p^{-1}=\{x:x\mathfrak p\subseteq R\}$ satisfies $\mathfrak a\mathfrak p^{-1}\ne\mathfrak a$
and $\mathfrak p\mathfrak p^{-1}=R$. \emph{(proved here)} \Cref{lem:ant:primeproducts}.
\item \textbf{Unique factorisation of ideals.} In a Dedekind domain every nonzero proper ideal is uniquely a product of primes; fractional ideals form a group,
free on the primes; to divide is to contain; $v_{\mathfrak p}$ is a discrete valuation. \emph{(proved here)} \Cref{thm:ant:idealfactorisation}.
\item \textbf{Norm of an ideal.} $N\mathfrak a=\#\mathcal O_K/\mathfrak a$ is multiplicative, $N(\alpha)=\abs{N_{K/\Q}\alpha}$, $N\mathfrak p=p^f$ with $\mathcal O_K/\mathfrak p\simeq\F_{p^f}$,
$\mathfrak p^k/\mathfrak p^{k+1}\simeq\mathcal O_K/\mathfrak p$; ideals of prime norm are prime; $\disc$ of a basis of $\mathfrak a$ is $(N\mathfrak a)^2d_K$. \emph{(proved here)}
\Cref{prop:ant:idealnorm}.
\item \textbf{Chinese remainder theorem, two generators, localisation.} $R/\mathfrak a\simeq\prod R/\mathfrak p_i^{a_i}$; every ideal is $(\alpha,\beta)$ with $\alpha$
arbitrary; finitely many primes implies PID; UFD iff PID; $R_{\mathfrak p}$ is a discrete valuation ring with residue field $R/\mathfrak p$ and valuation $v_{\mathfrak p}$.
\emph{(proved here)} \Cref{prop:ant:crt}.
\item \textbf{Class group.} $\Cl(K)=J_K/P_K$; $h_K=1$ iff PID iff UFD \emph{(proved here)}; Minkowski: $\Cl(K)$ is finite, each class contains an ideal of norm
$\le(4/\pi)^{r_2}(n!/n^n)\abs{d_K}^{1/2}$, and $\abs{d_K}>1$ for $K\ne\Q$ \emph{(stated; reference)}. Examples: $h=1$ for $D=-3,-4,-7,5,8,-11$, $h=2$ for $D=-20$.
\Cref{thm:ant:classgroup}.
\item \textbf{$\sum e_if_i=n$.} $\mathfrak p\mathcal O_L=\prod\mathfrak P_i^{e_i}$ with residue degrees $f_i$ satisfies $\sum e_if_i=[L:K]$; ramified, inert, split; $e$ and $f$
are multiplicative in towers. \emph{(proved here)} \Cref{thm:ant:efg}.
\item \textbf{Dedekind-Kummer.} If $p\nmid[\mathcal O_K:\Z[\theta]]$ and $g\equiv\prod\bar g_i^{e_i}\pmod p$, then $p\mathcal O_K=\prod(p,g_i(\theta))^{e_i}$ with
$f_i=\deg\bar g_i$ \emph{(proved here)}; relative version \emph{(stated; reference)}. \Cref{thm:ant:dedekindkummer}.
\item \textbf{Dedekind's discriminant theorem.} $p$ ramifies in $K$ iff $p\mid d_K$ \emph{(proved here)}; relative version \emph{(stated; reference)}.
\Cref{thm:ant:discram}.
\item \textbf{Decomposition and inertia groups.} For $L/K$ Galois, $G$ is transitive on the primes above $\mathfrak p$; $efg=n$; $D_{\mathfrak P}$ is the stabiliser,
$\#D_{\mathfrak P}=ef$, $I_{\mathfrak P}\trianglelefteq D_{\mathfrak P}$, both conjugate along the orbit. \emph{(proved here)} \Cref{prop:ant:galoissplitting}.
\item \textbf{Residue field extension.} $D_{\mathfrak P}\to\Gal(\kappa(\mathfrak P)/\kappa(\mathfrak p))$ is surjective with kernel $I_{\mathfrak P}$, so $\#I_{\mathfrak P}=e$ and
$D/I$ is cyclic of order $f$; behaviour in towers; split completely iff $D=1$, unramified iff $I=1$. \emph{(proved here)} \Cref{thm:ant:residuegalois}.
\item \textbf{Arithmetic and geometric Frobenius} (convention): $\Frob_p$ is arithmetic, $\chi_\ell(\Frob_p)=p$; $L_p(V,s)=\det(1-\Frob_p^{-1}p^{-s}\mid V^{I_p})^{-1}$; the two
conventions agree when the eigenvalues are inversion-stable, and differ at a multiplicative prime, where only the geometric one gives $1-q^{-1-s}$.
\Cref{conv:ant:frobenius}, \Cref{ex:ant:multiplicative}.
\item \textbf{The Frobenius element.} For $\mathfrak p$ unramified there is a unique $\Frob_{\mathfrak P}$ with $\Frob_{\mathfrak P}(x)\equiv x^{N\mathfrak p}\pmod{\mathfrak P}$; it
generates $D_{\mathfrak P}$; its conjugacy class is $\Frob_{\mathfrak p}$; it is $1$ iff $\mathfrak p$ splits completely; functoriality; at ramified primes only a coset of
inertia. \emph{(proved here)} \Cref{prop:ant:frobenius}.
\item \textbf{Frobenius and subfields.} For $M=L^H$, the primes of $M$ above $p$ correspond to the $\langle\Frob_{\mathfrak P}\rangle$-orbits on $G/H$, with residue degrees
the orbit lengths. \emph{(proved here)} \Cref{prop:ant:frobsubfield}.
\item \textbf{Dedekind's theorem on cycle types.} If $p\nmid a\disc(f)$, $p$ is unramified in the splitting field of $f$ and $\Frob_{\mathfrak P}$ permutes the roots with
cycle type the degrees of the irreducible factors of $f\bmod p$. \emph{(proved here)} \Cref{cor:ant:cycletype}.
\item \textbf{Frobenius elements in $G_\Q$.} Primes of $\bar\Z$ above $p$; $D_p$, $I_p$ closed, conjugate, with $1\to I_p\to D_p\to\Gal(\Fbar_p/\F_p)\to1$ exact and surjective
projections to finite levels; $\Frob_p$ defined up to $I_p$ and conjugation; unramified representations and $\rho(\Frob_p)$ \emph{(proved here)}; $D_p\simeq G_{\Q_p}$
\emph{(stated; \Cref{ch:local})}. \Cref{prop:ant:frobGQ}.
\item \textbf{Tame inertia.} If $\#I_{\mathfrak P}$ is prime to $p$, $I_{\mathfrak P}$ is cyclic and $F\tau F^{-1}=\tau^{N\mathfrak p}$ for every Frobenius $F$ and
$\tau\in I_{\mathfrak P}$. \emph{(proof sketched here; proved in \Cref{ch:local})} \Cref{prop:ant:tameconj}.
\item \textbf{Frobenius on a tame image of inertia.} For $\rho\colon G_\Q\to\GL_2(\Z_\ell)$, $\ell$ odd, $\ell\ne q$, with $\rho(I_q)$ finite of order prime to $q$:
$\rho(I_q)$ is cyclic and $\rho(F)\rho(\tau)\rho(F)^{-1}=\rho(\tau)^q$; so Frobenius centralises or inverts an element of order $3$ according as $q\equiv1$ or $2\pmod3$,
and one of order $4$ according as $q\equiv1$ or $3\pmod4$. \emph{(proved here, given \Cref{prop:ant:tameconj})} \Cref{cor:ant:frobsigma}.
\item \textbf{Splitting in quadratic fields.} $p$ ramifies, splits or is inert in $\Q(\sqrt d)$ according as $\bigl(\frac Dp\bigr)=0,1,-1$, with explicit primes,
including $p=2$. \emph{(proved here)} \Cref{thm:ant:quadsplit}.
\item \textbf{The character of a quadratic field.} $\chi_D(\Frob_p)=\bigl(\frac Dp\bigr)$ for $p\nmid D$; $\chi_D$ is ramified exactly at $p\mid D$; $\chi_{-4}(\Frob_p)=(-1)^{(p-1)/2}$,
$\chi_{-3}(\Frob_p)=1$ iff $p\equiv1\pmod3$, $\chi_{q^*}(\Frob_p)=\bigl(\frac pq\bigr)$ for $p\ne q$ including $p=2$. \emph{(proved here)} \Cref{prop:ant:quadchar}.
\item \textbf{Ramified only at $q$.} $\Q(\sqrt{q^*})$ is the unique quadratic field ramified only at the odd prime $q$; fundamental discriminants are products of
prime discriminants; $\chi_S=\prod_{q\in S}\chi_{q^*}$ is the character of $\Q(\sqrt{\prod q^*})$, of conductor $\prod q$, ramified exactly at $S$. \emph{(proved here)}
\Cref{prop:ant:ramonlyat}.
\item \textbf{Gaussian and Eisenstein primes.} $\Z[i]$, $\Z[\zeta_3]$ are PIDs; their primes are $1+i$, $a\pm bi$ ($p=a^2+b^2\equiv1\bmod 4$), $p\equiv3\pmod4$; resp.
$1-\zeta_3$, $a+b\zeta_3^{\pm1}$ ($p=a^2-ab+b^2\equiv1\bmod 3$), $p\equiv2\pmod3$. \emph{(proved here)} \Cref{prop:ant:gausseisenstein}.
\item \textbf{The prime above $3$.} $\sqrt{-3}=1+2\zeta_3=\zeta_3\pi$, $\pi^2=-3\zeta_3$, $3=-\zeta_3^2\pi^2$; $\Z[\zeta_3]/(\pi)=\F_3$; on $\Z[\zeta_3]/(\pi^2)$ multiplication by
$\zeta_3=1-\pi$ is a nontrivial unipotent with fixed line $(\pi)/(\pi^2)$; at $2$ in $\Z[i]$, $(2)=(1+i)^2$ and $i$ acts unipotently on $\Z[i]/(2)$. \emph{(proved here)}
\Cref{prop:ant:above3}.
\item \textbf{Quadratic algebras over $\Z_\ell$.} $\Z_\ell[X]/(g)$ is an unramified discrete valuation ring if $\bar g$ is irreducible, a ramified one with uniformiser
$x-c$ if $g(X+c)$ is Eisenstein, and $\Z_\ell\times\Z_\ell$ if $\bar g$ has distinct roots; lattices $T$ with $T\otimes\Q_\ell\simeq R\otimes\Q_\ell$ are free of rank one.
\emph{(proved here)} \Cref{lem:ant:quadraticZl}.
\item \textbf{$\Z_3[\zeta_3]$.} $\Z_3[X]/(X^2+X+1)$ is a discrete valuation ring with uniformiser $\pi=1-\zeta_3$ and residue field $\F_3$, $3=\pi^2\cdot$unit, and it is
$\Z_3\otimes\Z[\zeta_3]$ \emph{(proved here)}; it is the ring of integers of $\Q_3(\sqrt{-3})$ \emph{(stated; \Cref{ch:local})}. \Cref{thm:ant:z3zeta3dvr}.
\item \textbf{The lattice of Lemma lem:IV.} If $\sigma\in\Aut_{\Z_3}(\Z_3^2)$ satisfies $\sigma^2+\sigma+1=0$, then $T=\Z_3^2$ is free of rank one over $\Z_3[\sigma]\simeq\Z_3[\zeta_3]$,
$T/3T\simeq\Z_3[\zeta_3]/(\pi^2)$ with $\sigma$ a nontrivial unipotent fixing the line $\pi T/3T$, and $T^\sigma=0$. \emph{(proved here)} \Cref{cor:ant:z3zeta3modules}.
\item \textbf{$\Z_\ell[\zeta_3]$ and $\Z_\ell[i]$.} $\Z_\ell[\zeta_3]:=\Z_\ell[X]/(X^2+X+1)$ is ramified for $\ell=3$, $\simeq\Z_\ell\times\Z_\ell$ (not a domain) for $\ell\equiv1\pmod3$,
an unramified discrete valuation ring for $\ell\equiv2\pmod3$; likewise $\Z_\ell[i]$ according to $\ell\bmod4$; $\zeta_3\mapsto\zeta_3^2$ negates $\sqrt{-3}$, a unit iff
$\ell\ne3$. \emph{(proved here)} \Cref{prop:ant:zlzeta3}.
\item \textbf{Lattices with an automorphism of order $3$, $4$, $6$.} For $\ell$ odd and $g\in\Aut(\Z_\ell^2)$ with characteristic polynomial $\Phi_m$, $m\in\{3,4,6\}$
(e.g. $g\in\SL_2(\Z_\ell)$ of order $m$): $\Z_\ell[g]\simeq\Z_\ell[X]/(\Phi_m)$, the lattice is free of rank one over it, and the commutant of $g$ is $\Z_\ell[g]$.
\emph{(proved here)} \Cref{prop:ant:latticefree}.
\item \textbf{The commutant and the invariant line.} For $\sigma$ with characteristic polynomial $X^2+X+1$ on a plane over a field $F$ of characteristic $0$:
the commutant is $F[\sigma]\simeq F\otimes\Q(\zeta_3)$, its trace-zero part is $F(2\sigma+1)$ with $(2\sigma+1)^2=-3$, and $A\sigma A^{-1}=\sigma^q$, $q\ne3$ prime,
makes $A^{\pm1}$ act on it by $\chi_{-3}(q)$; for $\tau$ with characteristic polynomial $X^2+1$ (for example of order $4$ in $\SL_2$) and $q$ an odd prime, the line $F\tau$ and $\chi_{-4}(q)$; integrally, for $\ell$ odd,
$\End_{\langle\sigma\rangle}(T)=\Z_\ell[\sigma]\simeq\Z_\ell\otimes\Z[\zeta_3]$.
\emph{(proved here)} \Cref{prop:ant:commutant}.
\item \textbf{Cyclotomic integers.} $\mathcal O=\Z[\zeta_n]$; for $n=\ell^a$, $\ell=u(1-\zeta)^{\varphi(\ell^a)}$ is totally ramified; the discriminant formula; $p$ ramifies in
$\Q(\zeta_n)$ iff $p\mid n$ ($n\not\equiv2\bmod4$). \emph{((i), (ii) proved here for prime powers, (iii) for odd primes, (iv) in general; the rest stated with reference)} \Cref{thm:ant:cyclotomicintegers}.
\item \textbf{Splitting in cyclotomic fields.} For $n=p^am$: $e=\varphi(p^a)$, $f=$ order of $p$ modulo $m$, $g=\varphi(m)/f$; inertia $(\Z/p^a)^\times\times1$, decomposition
$(\Z/p^a)^\times\times\langle p\rangle$; $\Frob_p=\sigma_p$ for $p\nmid n$. \emph{(proved here)} \Cref{thm:ant:cyclotomicsplitting}.
\item \textbf{$\Q(\sqrt{\ell^*})$.} The quadratic subfield of $\Q(\zeta_\ell)$, ramified only at $\ell$; $\chi_{\ell^*}(\Frob_p)=\bigl(\frac p\ell\bigr)$ for all $p\ne\ell$; imaginary iff
$\ell\equiv3\pmod4$. \emph{(proved here)} \Cref{prop:ant:sqrtlstar}.
\item \textbf{Quadratic reciprocity}, with both supplements, from \Cref{prop:ant:quadchar} and \Cref{prop:ant:sqrtlstar}. \emph{(proved here)} \Cref{cor:ant:qr}.
\item \textbf{Kronecker-Weber.} Every abelian $K/\Q$ lies in some $\Q(\zeta_n)$; conductor; $G_\Q^{\mathrm{ab}}\simeq\hat\Z^\times$. \emph{(stated; reference)}
\Cref{thm:ant:kroneckerweber}.
\item \textbf{Artin map over $\Q$.} For $K\subseteq\Q(\zeta_n)$, $\Frob_p(K/\Q)=\sigma_p|_K$ depends only on $p\bmod n$ \emph{(proved here)}; idelic form and its sign
conventions \emph{(stated; reference)}. \Cref{thm:ant:artinQ}.
\item \textbf{Dirichlet characters and characters of $G_\Q$.} $\chi\mapsto\psi_\chi$, $\psi_\chi(\Frob_p)=\chi(p)$, is a bijection from primitive Dirichlet characters to
continuous characters of $G_\Q$; ramification and conductor match; $\chi_\ell$ and $\bar\chi_\ell=\psi_{\omega_\ell}$ \emph{(proved here, given Kronecker-Weber and Chebotarev)};
twisting newforms \emph{(stated; \Cref{ch:hecke})}. \Cref{thm:ant:dirichletgalois}.
\item \textbf{Dedekind zeta and densities.} $\zeta_K$ converges for $\mathrm{Re}\,s>1$ \emph{(proved here)}; simple pole at $1$ \emph{(stated; reference)}; Dirichlet
density; primes of degree one have density one \emph{(proved here, given the pole)}. \Cref{prop:ant:dedekindzeta}.
\item \textbf{Chebotarev.} For $L/K$ Galois and a class $C$, the primes with $\Frob_{\mathfrak p}=C$ have density $\#C/\#G$. \emph{(stated; reference)}
\Cref{thm:ant:chebotarev}.
\item \textbf{Consequences of Chebotarev.} Every element is a Frobenius; split primes have density $1/[L:K]$; Bauer's theorem; Dirichlet's theorem; inert primes
of $\Q(\sqrt{\ell^*})$ have density $\frac12$. \emph{(proved here, given \Cref{thm:ant:chebotarev})} \Cref{cor:ant:chebotarev}.
\item \textbf{Detecting a quadratic twist.} For $\rho$ semisimple, unramified outside $S$, in characteristic $\ne2$: $\rho\simeq\rho\otimes\chi_K$ iff $\tr\rho(\Frob_p)=0$ for
the primes $p\notin S$ inert in $K$ (outside density zero). \emph{(proved here, given \Cref{ch:galois})} \Cref{prop:ant:twistdetect}.
\item \textbf{Complex multiplication.} For $E/\Q$ with CM by $K$: class number one, $13$ values of $j$ \emph{(stated; reference)}; $\rho\simeq\rho\otimes\chi_K$, $a_p=0$ at good
inert primes, and primes ramified in $K$ are bad \emph{(proof sketched here)}; conversely $a_p=0$ at almost all inert primes forces CM by $K$, and without CM
$\{p:a_p=0\}$ has density $0$ \emph{(stated; reference)}; so CM by $\Q(\sqrt{\ell^*})$ forces $\ell\mid N$. \Cref{thm:ant:cm}.
\item \textbf{Coefficient fields.} The image of a ring finite over $\Z$ in $\C$ is an order in a number field; $K_f$ is a number field and $a_n(f)\in\mathcal O_{K_f}$
\emph{(proved here)}; conjugates are newforms, $K_f$ totally real \emph{(stated; \Cref{ch:hecke})}. \Cref{prop:ant:coefffield}.
\item \textbf{Primes above $\ell$.} Reduction modulo $\lambda$; congruences of integers modulo $\lambda$ and $\ell$; enlarging $\mathcal O$ and $\lambda$; lying over for
$\mathcal O[\alpha]$; the root of $X^2-aX+p$ congruent to a given root modulo a prime above $\lambda$. \emph{(proved here)} \Cref{prop:ant:primesabove}.
\item \textbf{Maximal ideals from congruences.} $\ker(R\to\mathcal O\to\mathcal O/\lambda)$ is maximal of residue characteristic $\ell$, the same for congruent
homomorphisms, and an ideal $J$ inside it with $R/J\simeq\Z/r$ has $\ell\mid r$ \emph{(proved here)}; $R\otimes\Z_\ell=\prod R_{\mathfrak m}$ \emph{(stated;
\Cref{ch:modules})}. \Cref{prop:ant:congruencemaximal}.
\item \textbf{$\mathcal O_K\otimes\Z_\ell=\prod_{\lambda\mid\ell}\mathcal O_{K,\lambda}$}, each a complete discrete valuation ring of rank $e_\lambda f_\lambda$ over $\Z_\ell$.
\emph{(stated; \Cref{ch:local})} \Cref{prop:ant:otimesZl}.
\end{enumerate}

\section*{Where the paper uses this chapter}

\begin{itemize}
\item \textbf{Lemma lem:IV} (types $\kod{IV}$, $\kod{IV}^*$ at $q\ne3$): $\Z_3[\zeta_3]$ a discrete valuation ring with uniformiser $\pi$, \Cref{thm:ant:z3zeta3dvr}; $T_3E$ free of
rank one over it and $E[3]\simeq\Z_3[\zeta_3]/(\pi^2)$ with a unipotent action, \Cref{cor:ant:z3zeta3modules}, \Cref{prop:ant:above3}; Frobenius conjugating $\sigma$ to
$\sigma^q=\sigma^{\chi_{-3}(q)}$, \Cref{cor:ant:frobsigma}, \Cref{prop:ant:quadchar}; the line of $\ad V$ and its sign, \Cref{prop:ant:commutant}.
\item \textbf{Lemma lem:euler} (the omitted factors at tame primes) \textbf{and the table of Section 1}: the tame relation, \Cref{prop:ant:tameconj},
\Cref{cor:ant:frobsigma}; $\End_{I_q}(T)=\Z_\ell[\sigma]\simeq\Z_\ell\otimes\Z[\zeta_3]$, not a field for $\ell\equiv1\pmod3$, \Cref{prop:ant:zlzeta3}, \Cref{prop:ant:latticefree},
\Cref{prop:ant:commutant}; $\chi_{-3}(q)$, $\chi_{-4}(q)$, \Cref{prop:ant:quadchar}, \Cref{ex:ant:characters}; the geometric Frobenius in $L_q(A_f,s)$ and the row
$\kod{I}_n$, \Cref{conv:ant:frobenius}, \Cref{ex:ant:multiplicative}.
\item \textbf{Section 1.1} (the local terms): $I_p$, $\Frob_p$, $V^{I_p}$, \Cref{prop:ant:frobGQ}; $L_p(A_f,s)=\det(1-\Frob_p^{-1}p^{-s}\mid A_f^{I_p})^{-1}$,
\Cref{conv:ant:frobenius}. \textbf{Lemma lem:fpr}: $L_p(A_f,1)=L_p(A_f(1),0)$, \Cref{conv:ant:frobenius}. \textbf{Lemma lem:tamfactor}: the arithmetic
Frobenius, with $\Frob^{-1}t\Frob=t^{q^{-1}}$ on tame inertia, \Cref{conv:ant:frobenius}, \Cref{prop:ant:tameconj}.
\item \textbf{Section 1.4, Lemma lem:twist, Theorems thm:twist and thm:quant}: $\chi_{q^*}$, the character of conductor $q$ of $\Q(\sqrt{q^*})$, ramified only at $q$,
\Cref{prop:ant:ramonlyat}, \Cref{prop:ant:quadchar}, \Cref{cor:ant:qr}; $\chi_S$ and its conductor $\prod_{q\in S}q$ (the paper's $D$), \Cref{prop:ant:ramonlyat}(iii); characters as Dirichlet and as
Galois characters, \Cref{thm:ant:dirichletgalois}, \Cref{ex:ant:characters}.
\item \textbf{Theorem thm:main} (the hypothesis on $\Gal(\Qbar/\Q(\sqrt{\ell^*}))$): $\Q(\sqrt{\ell^*})\subseteq\Q(\zeta_\ell)$ and its splitting law, \Cref{prop:ant:sqrtlstar};
the condition $\bar\rho\simeq\bar\rho\otimes\chi_{\ell^*}$ (reached in the paper by Clifford theory, \Cref{ch:reps}; its trace form
\Cref{prop:ant:twistdetect} is not used by the paper), \Cref{warn:ant:chil}.
\item \textbf{Lemma lem:h0 and Proposition prop:ns}: $\bar\chi_3^{\mathrm{cyc}}=\chi_{-3}$ and $\chi_{\ell^*}\ne\bar\chi_\ell$ for $\ell\ge5$, \Cref{ex:ant:characters},
\Cref{warn:ant:chil}; the isomorphism $\bar\rho\simeq\bar\rho\otimes\chi_{-3}$ comes from induction from $\Q(\sqrt{-3})$, $\bar\rho\otimes\Fbar_3\simeq\Ind\psi$ with
$\Ind\psi\otimes\chi_{-3}\simeq\Ind\psi$ (\Cref{ch:reps}), not from Frobenius traces; \Cref{prop:ant:twistdetect} gives the equivalent trace criterion
($3\mid a_p$ for almost all $p\equiv2\pmod3$), which the book uses in an example of \Cref{ch:galois} and the paper does not use. Only for the original version of
Lemma lem:h0: complex multiplication and $a_p=0$ at inert primes, \Cref{thm:ant:cm}, \Cref{cor:ant:chebotarev} (\Cref{pl:ant:h0}).
\item \textbf{Proposition prop:Sf}: Clifford theory for $\bar\rho|_{G_K}$, $K=\Q(\sqrt{\ell^*})$ (\Cref{ch:reps}), and the restriction to $I_\ell$ of
$\chi_{\ell^*}=\bar\chi_\ell^{(\ell-1)/2}$, a power of $\bar\chi_\ell|_{I_\ell}$ (the paper's fundamental character $\omega$ of level one), \Cref{ex:ant:characters},
\Cref{warn:ant:chil}.
\item \textbf{Proposition prop:ns-structure}(ii) and (iii): the trace-zero commutant of an element $\tau$ of order $4$ of $\rho(I_p)$ is the line $F\tau$, on which
Frobenius acts by $\pm1$, \Cref{prop:ant:commutant}(iv); in (iii) the tame relation $\tau\mapsto\tau^p$ (\Cref{cor:ant:frobsigma}) makes the sign $\chi_{-4}(p)$, and
$\chi_{-4}(p)=\chi_{-3}(p)$ iff $p\equiv\pm1\pmod{12}$, \Cref{prop:ant:quadchar}(iii).
\item \textbf{Lemma lem:tamadd} (potentially good primes, $\ell=3$): $T$ free of rank one over $\Z_3[X]/(X^2+X+1)$, with $X$ acting as an element $\sigma$ of order
$3$ of $\rho(I_q)$, at every $q\ne3$, $q=2$ included, \Cref{cor:ant:z3zeta3modules}(i), and the basis in which $\sigma$ is the companion matrix, \Cref{pl:ant:IV}.
\item \textbf{Proposition prop:forms}: $q\equiv\chi_{-3}(q)\pmod3$ for $q\ne3$, \Cref{prop:ant:quadchar}(iii).
\item \textbf{Lemma lem:crit} (congruence criterion): coefficients in the ring of integers $\mathcal O$ of a number field, \Cref{prop:ant:coefffield}; $\lambda\mid\ell$ and the
maximal ideal of residue characteristic $\ell$, $\ell\mid r$, \Cref{prop:ant:primesabove}, \Cref{prop:ant:congruencemaximal}; over $\Z_\ell$, \Cref{prop:ant:otimesZl}.
\item \textbf{Lemma lem:old} (oldforms): Frobenius at unramified primes and Chebotarev, \Cref{prop:ant:frobenius}, \Cref{prop:ant:frobGQ}, \Cref{thm:ant:chebotarev}; ``after
enlarging $\mathcal O$ and $\lambda$'' and ``a prime above $\lambda$ in $\mathcal O[\alpha_p]$'', \Cref{prop:ant:primesabove}(iii) to (v).
\item \textbf{Section 3 and Example ex:qual}: the twisting fields $\Q(\sqrt{-7})$, $\Q(\sqrt{-11})$, \Cref{ex:ant:quadexamples}; the congruences
$a_p(g)\equiv a_p(f)\pmod\lambda$ at all $p\nmid3N$ are traces of Frobenius of $\bar\rho_g\simeq\bar\rho$, \Cref{prop:ant:frobGQ}(iv) (Chebotarev is needed only
for the converse, as in Lemma lem:old).
\end{itemize}

\section{Number fields and their rings of integers}\label{sec:ant:integers}

This section sets up the objects behind every ``$\mathcal O$'' and every ``prime $\lambda\mid\ell$'' of the paper: number fields, their rings of
integers, traces and norms, integral bases and the discriminant, and the units. The paper's Lemmas lem:crit and lem:old take the
coefficients of a newform in the ring of integers $\mathcal O$ of a number field; Lemma lem:IV and Lemma lem:euler work inside
$\Z[\zeta_3]$ and $\Z[i]$, the rings of integers of $\Q(\sqrt{-3})$ and $\Q(i)$, whose discriminants $-3$ and $-4$ name the characters $\chi_{-3}$ and
$\chi_{-4}$ of the table of Section 1. Nothing here is specific to elliptic curves.

\begin{convention}\label{conv:ant:conventions}
(a) We fix an algebraic closure $\Qbar$ of $\Q$ inside $\C$ (as in \Cref{ch:galois}); every number field is a subfield of $\Qbar$, and $\bar\Z$
denotes the ring of all algebraic integers of $\Qbar$. We write $\zeta_n=e^{2\pi i/n}$, and $\zeta_3=(-1+\sqrt{-3})/2$.
(b) The letters $p$, $q$, $\ell$ denote rational primes, as in the paper ($\ell$ the coefficient prime, $q$ a bad prime of $E$); the size of a
residue field is written $N\mathfrak p$, never $q$. Uniformisers of discrete valuation rings are written $\varpi$; the paper's element
$\pi=1-\zeta_3$ keeps its name, and in \Cref{sec:ant:quadratic} the letter $\pi$ means nothing else.
(c) $D$ denotes the discriminant of a quadratic field, and $D_{\mathfrak P}$, $D_p$ decomposition groups; no other meaning of $D$ occurs. The paper's $D$
of Theorem thm:twist, the conductor of $\chi_S$, is written out here as $\prod_{q\in S}q$.
(d) The quadratic character of a quadratic field $K$ of discriminant $D$ is written $\chi_D$ (convention B-C6 of
Appendix B): $\chi_{-4}$ for $\Q(i)$, $\chi_{-3}$ for $\Q(\sqrt{-3})$, $\chi_8$ for $\Q(\sqrt2)$, $\chi_{q^*}$ for $\Q(\sqrt{q^*})$; the index is
always a fundamental discriminant, and $\chi_\ell$ with a letter index is the $\ell$-adic cyclotomic character.
(e) $\Frob$ is the arithmetic Frobenius and $\Frob^{-1}$ the geometric one (\Cref{conv:ant:frobenius}). Since the cyclotomic polynomials are
written $\Phi_n$ in this chapter, the geometric Frobenius is always written $\Frob_p^{-1}$ here (Appendix B also writes it $\Phi_p$).
\end{convention}

\begin{definition}\label{def:ant:numberfield}
A \emph{number field} is a field $K\supseteq\Q$ with $n=[K:\Q]<\infty$. By the primitive element theorem (\Cref{ch:galois}) $K=\Q(\theta)$
for some $\theta$, and there are exactly $n$ field embeddings $\sigma_1,\dots,\sigma_n\colon K\to\C$, sending $\theta$ to the $n$ roots of its
minimal polynomial. If $r_1$ of them have image in $\R$, the other $2r_2$ come in complex conjugate pairs and $n=r_1+2r_2$. The field is
\emph{totally real} if $r_2=0$ and \emph{imaginary quadratic} if $n=2$ and $r_1=0$.
\end{definition}

\begin{proposition}[algebraic integers]\label{prop:ant:integral}
Let $K$ be a number field and $\alpha\in K$. The following are equivalent, and $\alpha$ is then called an \emph{algebraic integer}:
(o) $\alpha$ is a root of a monic polynomial in $\Z[X]$; (i) the minimal polynomial of $\alpha$ over $\Q$ lies in $\Z[X]$; (ii) $\Z[\alpha]$ is a
finitely generated $\Z$-module; (iii) $\alpha$ lies in a subring of $K$ that is finitely generated as a $\Z$-module. The algebraic integers of
$K$ form a subring $\mathcal O_K$, the \emph{ring of integers}; $\mathcal O_K\cap\Q=\Z$; for every $\alpha\in K$ there is an integer $m\ge1$ with
$m\alpha\in\mathcal O_K$, so $\Frac(\mathcal O_K)=K$; and $\mathcal O_K$ is \emph{integrally closed} in $K$: an element of $K$ that is a root of a
monic polynomial with coefficients in $\mathcal O_K$ lies in $\mathcal O_K$.
\end{proposition}

\begin{proof}
(o)$\Rightarrow$(i): if $\alpha$ is a root of a monic $g\in\Z[X]$, its minimal polynomial $m_\alpha$ divides $g$ in $\Q[X]$, and a monic factor
in $\Q[X]$ of a monic polynomial of $\Z[X]$ lies in $\Z[X]$ by Gauss's lemma (\Cref{ch:rings}). (i)$\Rightarrow$(ii): if $m_\alpha\in\Z[X]$ is
monic of degree $d$, division by $m_\alpha$ shows $\Z[\alpha]=\Z+\Z\alpha+\dots+\Z\alpha^{d-1}$. (ii)$\Rightarrow$(iii) is trivial.
(iii)$\Rightarrow$(o): let $M\ni\alpha$ be a subring generated as a $\Z$-module by $m_1,\dots,m_r$. Write $\alpha m_i=\sum_ja_{ij}m_j$ with
$a_{ij}\in\Z$, so $(\alpha I-A)\underline m=0$ for $A=(a_{ij})$ and the column $\underline m$. Multiplying by the adjugate of $\alpha I-A$ gives
$\det(\alpha I-A)m_i=0$ for all $i$; as $1\in M$ is a $\Z$-combination of the $m_i$, $\det(\alpha I-A)=0$, and $\det(XI-A)$ is monic in $\Z[X]$.

If $\alpha,\beta$ are algebraic integers, roots of monic integer polynomials of degrees $a$ and $b$, the ring $\Z[\alpha,\beta]$ is generated
as a $\Z$-module by the $\alpha^i\beta^j$, $i<a$, $j<b$ (every element is a $\Z$-combination of monomials $\alpha^r\beta^s$, and the monic equation of degree $a$
rewrites $\alpha^a$, hence by induction every $\alpha^r$ with $r\ge a$, as a $\Z$-combination of $1,\alpha,\dots,\alpha^{a-1}$; likewise for $\beta$); it contains $\alpha\pm\beta$ and $\alpha\beta$, which are therefore integral by (iii).
A rational root $x=u/v$ (in lowest terms) of a monic $X^d+c_{d-1}X^{d-1}+\dots+c_0\in\Z[X]$ satisfies $u^d=-v(c_{d-1}u^{d-1}+\dots+c_0v^{d-1})$, so
$v\mid u^d$ and $v=\pm1$: hence $\mathcal O_K\cap\Q=\Z$. If $\alpha$ is a root of $a_dX^d+\dots+a_0\in\Z[X]$ with $a_d\ne0$, then $a_d\alpha$ is a root
of the monic polynomial $X^d+a_{d-1}X^{d-1}+a_{d-2}a_dX^{d-2}+\dots+a_0a_d^{d-1}$. Finally, if $x\in K$ satisfies $x^d+c_{d-1}x^{d-1}+\dots+c_0=0$
with $c_i\in\mathcal O_K$, then $\Z[c_0,\dots,c_{d-1},x]$ is a finitely generated $\Z$-module (each $c_i$ is integral, and powers of $x$ beyond
$x^{d-1}$ are reduced by the equation), so $x\in\mathcal O_K$ by (iii).
\end{proof}

\begin{definition}\label{def:ant:tracenorm}
For $\alpha\in K$ let $m_\alpha\colon K\to K$, $x\mapsto\alpha x$, a $\Q$-linear map. The \emph{trace} and \emph{norm} of $\alpha$ are
$\Tr_{K/\Q}(\alpha)=\tr(m_\alpha)$ and $N_{K/\Q}(\alpha)=\det(m_\alpha)$. For a tower $L\supseteq K$ of number fields, $\Tr_{L/K}$ and $N_{L/K}$ are
defined in the same way with $L$ viewed as a $K$-vector space.
\end{definition}

The trace and the norm turn an element of $K$ into rational numbers by linear algebra, without choosing an embedding into $\C$; this is what makes them
useful for integrality questions. For example, in $K=\Q(i)$ with basis $1,i$, multiplication by $\alpha=a+bi$ sends $1\mapsto a+bi$ and $i\mapsto-b+ai$, so its
matrix is $\sm a{-b}ba$, with trace $2a$ and determinant $a^2+b^2$: $\Tr(a+bi)=(a+bi)+(a-bi)$ and $N(a+bi)=(a+bi)(a-bi)$, the sum and the product of the two
conjugates, which is the general pattern of \Cref{prop:ant:tracenorm}(i).

\begin{proposition}[trace and norm]\label{prop:ant:tracenorm}
Let $K$ be a number field of degree $n$ with embeddings $\sigma_1,\dots,\sigma_n$, and $\alpha\in K$.
\begin{enumerate}[label=(\roman*)]
\item The characteristic polynomial of $m_\alpha$ is $\prod_i(X-\sigma_i(\alpha))=m_\alpha(X)^{n/d}$, $d=[\Q(\alpha):\Q]$, $m_\alpha$ the minimal
polynomial. Hence $\Tr_{K/\Q}(\alpha)=\sum_i\sigma_i(\alpha)$ and $N_{K/\Q}(\alpha)=\prod_i\sigma_i(\alpha)$; $\Tr$ is $\Q$-linear, $N$ is
multiplicative, and $N(\alpha)=0$ iff $\alpha=0$. If $K=\Q(\alpha)$ with minimal polynomial $X^n+c_{n-1}X^{n-1}+\dots+c_0$, then
$\Tr(\alpha)=-c_{n-1}$ and $N(\alpha)=(-1)^nc_0$.
\item If $\alpha\in\mathcal O_K$ then $\Tr(\alpha),N(\alpha)\in\Z$, and $\alpha\in\mathcal O_K^\times$ if and only if $N(\alpha)=\pm1$.
\item For number fields $L\supseteq K$: $\Tr_{L/\Q}=\Tr_{K/\Q}\circ\Tr_{L/K}$ and $N_{L/\Q}=N_{K/\Q}\circ N_{L/K}$.
\end{enumerate}
\end{proposition}

\begin{proof}
(i) If $\beta_1,\dots,\beta_{n/d}$ is a basis of $K$ over $\Q(\alpha)$, the elements $\beta_j\alpha^i$ ($0\le i<d$) form a basis of $K$ over $\Q$ in which
$m_\alpha$ is block diagonal with $n/d$ copies of the companion matrix of $m_\alpha(X)$; so its characteristic polynomial is $m_\alpha(X)^{n/d}$.
Each of the $d$ embeddings of $\Q(\alpha)$, which send $\alpha$ to the $d$ distinct roots of $m_\alpha$, extends to exactly $n/d$ embeddings of
$K$ (\Cref{ch:galois}), so $\prod_i(X-\sigma_i(\alpha))=m_\alpha(X)^{n/d}$ as well. The trace of a matrix is the sum of the roots of its
characteristic polynomial and the determinant their product, which gives the two formulas. The remaining statements are properties of the trace and
the determinant: $m_{\alpha+\beta}=m_\alpha+m_\beta$ and $m_{c\alpha}=cm_\alpha$ for $c\in\Q$ give linearity of $\Tr$; $m_{\alpha\beta}=m_\alpha m_\beta$ gives
multiplicativity of $N$; $m_\alpha$ is invertible for $\alpha\ne0$ (with inverse $m_{\alpha^{-1}}$) and $m_0=0$. If $K=\Q(\alpha)$, the characteristic polynomial is the
minimal polynomial $X^n+c_{n-1}X^{n-1}+\dots+c_0$, so the sum of its roots is $-c_{n-1}$ and their product is $(-1)^nc_0$.
(ii) By \Cref{prop:ant:integral}(i), $m_\alpha(X)\in\Z[X]$, so the characteristic polynomial has integer coefficients. If $\alpha\beta=1$ with
$\beta\in\mathcal O_K$, then $N(\alpha)N(\beta)=1$ in $\Z$. Conversely, let $\chi(X)=X^n+\dots+c_1X+c_0$ be the characteristic polynomial of
$m_\alpha$, with $c_0=(-1)^nN(\alpha)=\pm1$; by Cayley-Hamilton $\chi(\alpha)=0$, so $\alpha^{-1}=-c_0^{-1}(\alpha^{n-1}+\dots+c_1)\in\Z[\alpha]\subseteq\mathcal O_K$.
(iii) By (i) applied over $K$ (with the same proof), $\Tr_{L/K}(\alpha)=\sum_\sigma\sigma(\alpha)$ over the $[L:K]$ embeddings $\sigma\colon L\to\Qbar$
that are the identity on $K$. Fix an embedding $\tau\colon K\to\Qbar$ and extend it to an automorphism $\tilde\tau$ of $\Qbar$ (\Cref{ch:galois}); the
embeddings of $L$ extending $\tau$ are the $\tilde\tau\circ\sigma$, so their sum at $\alpha$ is $\tau(\Tr_{L/K}\alpha)$. Summing over the $[K:\Q]$
choices of $\tau$ gives $\Tr_{L/\Q}(\alpha)=\Tr_{K/\Q}(\Tr_{L/K}\alpha)$; the norm is the same with products.
\end{proof}

\begin{theorem}[integral bases and the discriminant]\label{thm:ant:integralbasis}
Let $K$ be a number field of degree $n$.
\begin{enumerate}[label=(\roman*)]
\item The trace form $(x,y)\mapsto\Tr_{K/\Q}(xy)$ is a nondegenerate symmetric bilinear form on $K$.
\item $\mathcal O_K$ is a free $\Z$-module of rank $n$. For a $\Z$-basis $\beta_1,\dots,\beta_n$ (an \emph{integral basis}) the integer
$\disc(\beta_1,\dots,\beta_n)=\det\bigl(\Tr(\beta_i\beta_j)\bigr)=\det\bigl(\sigma_k(\beta_j)\bigr)^2$ is nonzero and independent of the basis; it is
the \emph{discriminant} $d_K$ of $K$.
\item If $\alpha_1,\dots,\alpha_n\in\mathcal O_K$ are linearly independent over $\Q$ and $M=\Z\alpha_1+\dots+\Z\alpha_n$, then
$\disc(\alpha_1,\dots,\alpha_n)=[\mathcal O_K:M]^2\,d_K$. In particular, if $\disc(\alpha_1,\dots,\alpha_n)$ is squarefree then $M=\mathcal O_K$.
\item If $K=\Q(\theta)$ with $\theta\in\mathcal O_K$ of minimal polynomial $g$, with roots $\theta_1,\dots,\theta_n$, then
\[
\disc(1,\theta,\dots,\theta^{n-1})=\prod_{i<j}(\theta_i-\theta_j)^2=\disc(g)=(-1)^{n(n-1)/2}N_{K/\Q}(g'(\theta)).
\]
\end{enumerate}
\end{theorem}

\begin{proof}
(i) Symmetry and bilinearity come from the linearity of $\Tr$ (\Cref{prop:ant:tracenorm}(i)). Nondegeneracy means that no $x\ne0$ is orthogonal to
all of $K$; and if $x\ne0$, the element $y=x^{-1}$ gives $\Tr(xy)=\Tr(1)=n\ne0$ ($m_1$ is the identity of an $n$-dimensional space). (ii) By \Cref{prop:ant:integral} a $\Q$-basis of $K$ can be multiplied by integers to lie in
$\mathcal O_K$; let $\alpha_1,\dots,\alpha_n$ be such a basis and $\alpha_1^\vee,\dots,\alpha_n^\vee$ the dual basis for the trace form, which exists
by (i). For $x\in\mathcal O_K$, $x=\sum_i\Tr(x\alpha_i)\alpha_i^\vee$ with $\Tr(x\alpha_i)\in\Z$ (\Cref{prop:ant:tracenorm}(ii)). So
$\bigoplus_i\Z\alpha_i\subseteq\mathcal O_K\subseteq\bigoplus_i\Z\alpha_i^\vee$, and $\mathcal O_K$, a subgroup of a free abelian group of rank $n$
containing one of rank $n$, is free of rank $n$ (\Cref{ch:modules}). Since $\Tr(\beta_i\beta_j)=\sum_k\sigma_k(\beta_i)\sigma_k(\beta_j)$, the Gram matrix is
$S^{\mathsf T}S$ with $S=(\sigma_k(\beta_j))_{k,j}$, which gives the second expression; it is nonzero by (i). If $\beta'=C\beta$ is another basis,
$C\in\GL_n(\Z)$, the Gram matrix becomes $C(\cdot)C^{\mathsf T}$ and the determinant is multiplied by $\det(C)^2=1$.
(iii) Write $\alpha=C\beta$ with $C\in\Mat_n(\Z)$; then $\disc(\alpha)=\det(C)^2d_K$, and $\abs{\det C}=[\mathcal O_K:M]$ by the elementary
divisor theorem (\Cref{ch:modules}). If $\disc(\alpha)$ is squarefree, the square $[\mathcal O_K:M]^2$ dividing it is $1$.
(iv) With $\sigma_k(\theta)=\theta_k$, the matrix $(\sigma_k(\theta^{j-1}))$ is the Vandermonde matrix of the $\theta_k$, of determinant
$\prod_{i<j}(\theta_j-\theta_i)$; its square is $\prod_{i<j}(\theta_i-\theta_j)^2$, which is $\disc(g)$ by definition. Writing $g(X)=\prod_k(X-\theta_k)$, the
product rule gives $g'(\theta_i)=\prod_{k\ne i}(\theta_i-\theta_k)$. In the product $\prod_i\prod_{k\ne i}(\theta_i-\theta_k)$ each unordered pair $\{i,j\}$, $i<j$,
contributes the two factors $(\theta_i-\theta_j)(\theta_j-\theta_i)=-(\theta_i-\theta_j)^2$; there are $n(n-1)/2$ pairs, so the product is
$(-1)^{n(n-1)/2}\prod_{i<j}(\theta_i-\theta_j)^2$. Finally $\prod_ig'(\theta_i)=\prod_i\sigma_i(g'(\theta))=N(g'(\theta))$ by \Cref{prop:ant:tracenorm}(i). For example, for
$g=X^2-d$ the formula gives $(-1)^1N(2\sqrt d)=-(2\sqrt d)(-2\sqrt d)=4d$, the discriminant of $1,\sqrt d$ found again in \Cref{prop:ant:quadraticintegers}.
\end{proof}

\begin{proposition}[quadratic fields]\label{prop:ant:quadraticintegers}
Let $d\in\Z$ be squarefree, $d\ne0,1$, and $K=\Q(\sqrt d)$. If $d\equiv2,3\pmod4$ then $\mathcal O_K=\Z[\sqrt d]$ and $d_K=4d$; if $d\equiv1\pmod4$
then $\mathcal O_K=\Z[\frac{1+\sqrt d}2]$ and $d_K=d$. A \emph{fundamental discriminant} is an integer $D\ne1$ that is either squarefree with
$D\equiv1\pmod4$, or $D=4m$ with $m$ squarefree and $m\equiv2,3\pmod4$; the map $K\mapsto d_K$ is a bijection from quadratic fields to fundamental
discriminants, with inverse $D\mapsto\Q(\sqrt D)$. The norm forms are $N(a+b\sqrt d)=a^2-db^2$ and, for $d\equiv1\pmod4$ and $\theta_d=\frac{1+\sqrt d}2$,
$N(a+b\theta_d)=a^2+ab+\frac{1-d}4b^2$.
\end{proposition}

\begin{proof}
Let $\alpha=a+b\sqrt d$ with $a,b\in\Q$, $b\ne0$. Its minimal polynomial is $X^2-2aX+(a^2-db^2)$, so $\alpha\in\mathcal O_K$ iff $2a\in\Z$ and
$a^2-db^2\in\Z$. If $a\in\Z$ then $db^2\in\Z$, and $b\in\Z$ because $d$ is squarefree (if $b=u/v$ in lowest terms, $v^2\mid du^2$, so $v^2\mid d$). If
$a=A/2$ with $A$ odd, then $4(a^2-db^2)=A^2-d(2b)^2\in4\Z$, so $d(2b)^2\in\Z$, $B:=2b\in\Z$ as before, and $A^2\equiv dB^2\pmod4$ with $A$ odd forces
$B$ odd and $d\equiv1\pmod4$. This gives $\mathcal O_K$. For the basis $1,\sqrt d$ the Gram matrix of the trace form is $\mathrm{diag}(2,2d)$, of
determinant $4d$; for $1,\theta_d$ it is $\sm21 1{(1+d)/2}$, of determinant $d$. The fundamental discriminants are exactly these values: $d$ for
$d\equiv1\pmod4$ and $4d$ for $d\equiv2,3\pmod4$, and $\Q(\sqrt{d_K})=K$. The norm forms follow from $N(x)=x\bar x$.
\end{proof}

\begin{example}\label{ex:ant:quadraticrings}
$\Q(i)$: $\mathcal O=\Z[i]$, $D=-4$, $N(a+bi)=a^2+b^2$. $\Q(\sqrt{-3})$: $\mathcal O=\Z[\zeta_3]$ since $\zeta_3=\frac{-1+\sqrt{-3}}2=\theta_{-3}-1$, $D=-3$,
and $N(a+b\zeta_3)=a^2-ab+b^2$; the subring $\Z[\sqrt{-3}]$ has index $2$ (its discriminant is $-12=2^2\cdot(-3)$, \Cref{thm:ant:integralbasis}(iii)) and is
not integrally closed, which is why unique factorisation fails in it (\Cref{ch:rings}, where $4=2\cdot2=(1+\sqrt{-3})(1-\sqrt{-3})$). $\Q(\sqrt5)$:
$\mathcal O=\Z[\frac{1+\sqrt5}2]$, $D=5$. $\Q(\sqrt{-7})$, $\Q(\sqrt{-11})$: $D=-7$, $-11$. $\Q(\sqrt{-5})$, $\Q(\sqrt2)$: $D=-20$, $8$.
\end{example}

\begin{theorem}[units]\label{thm:ant:units}
(i) (Dirichlet's unit theorem.) $\mathcal O_K^\times\simeq\mu(K)\times\Z^{r_1+r_2-1}$, where $\mu(K)$, the group of roots of unity in $K$, is finite
and cyclic. (ii) If $K$ is imaginary quadratic, $\mathcal O_K^\times=\mu(K)$, which is $\{\pm1\}$ except for $\mu(\Q(i))=\{\pm1,\pm i\}=\mu_4$ and
$\mu(\Q(\sqrt{-3}))=\{\pm1,\pm\zeta_3,\pm\zeta_3^2\}=\mu_6$. (iii) $\Z[\sqrt2]^\times=\{\pm(1+\sqrt2)^k:k\in\Z\}$.
\end{theorem}

\begin{proof}
(i) We do not prove this; see \cite[Chapter I, Section 7]{Neukirch} or \cite[Chapter 5]{Marcus77}. (ii) A unit has norm $\pm1$
(\Cref{prop:ant:tracenorm}(ii)), and norms are $\ge0$ in an imaginary quadratic field, so units are the elements of norm $1$. For $d\equiv2,3\pmod4$,
$d<0$: $a^2+\abs db^2=1$ has the solutions $(\pm1,0)$, and also $(0,\pm1)$ when $d=-1$. For $d\equiv1\pmod4$, $d<0$: with $\alpha=(A+B\sqrt d)/2$, $A\equiv B
\pmod2$, the equation is $A^2+\abs dB^2=4$; for $\abs d\ge7$ only $B=0$, $A=\pm2$; for $d=-3$ also $A=\pm1$, $B=\pm1$, which are the six elements
$\pm1,\pm\zeta_3,\pm\zeta_3^2$. (iii) Let $u_0=1+\sqrt2$, a unit of norm $1-2=-1$. If $u=a+b\sqrt2>1$ is a unit ($a,b\in\Z$), its conjugate
$\bar u=a-b\sqrt2$ satisfies $u\bar u=N(u)=\pm1$, so $\bar u=\pm u^{-1}$ has $\abs{\bar u}<1$; hence $a=(u+\bar u)/2>(1-1)/2=0$ and $b\sqrt2=(u-\bar u)/2>0$.
So $a$ and $b$ are positive integers, $a,b\ge1$, and $u\ge1+\sqrt2=u_0$: $u_0$ is the smallest unit greater than $1$. Given any unit $u>0$, choose $k$ with
$u_0^k\le u<u_0^{k+1}$; then $u\,u_0^{-k}$ is a unit in $[1,u_0)$, so it is $1$. Units $u<0$ are $-1$ times positive ones.
\end{proof}

In the paper the units of $\Z[\zeta_3]$ appear in ``$\pi^2=-3$ up to a unit'' (Lemma lem:IV):
$3=-\zeta_3^2(1-\zeta_3)^2$ and $\zeta_3^2\in\mu_6$ (\Cref{prop:ant:above3}).

So far $\mathcal O_K$ is known as an abelian group ($\Z^n$) and through its units. What it lacks in general is unique factorisation of elements, and
the next section replaces it by unique factorisation of ideals; that is the structure behind the primes $\lambda$, the valuations and the residue fields used
everywhere later.

\section{Dedekind domains, ideals, norms and class groups}\label{sec:ant:dedekind}

Unique factorisation of elements fails in rings of integers, but unique factorisation of ideals holds, and it is what gives meaning to
``a prime $\lambda\mid\ell$'', to the valuation $v_\lambda$ and to the residue field $\mathcal O/\lambda$ in Lemmas lem:crit and lem:old of the paper.
This section proves it (\Cref{thm:ant:idealfactorisation}), with the norm of an ideal, the Chinese remainder theorem and the localisations,
which are discrete valuation rings in the sense of \Cref{ch:modules}. The class group is described but not used by the paper: the rings
$\Z[i]$ and $\Z[\zeta_3]$ that the paper needs are principal ideal domains.

\begin{definition}\label{def:ant:dedekind}
A \emph{Dedekind domain} is a Noetherian integral domain $R$, not a field, which is integrally closed in its fraction field and in which every
nonzero prime ideal is maximal.
\end{definition}

The three conditions are exactly what the proof of unique factorisation of ideals uses, and each can fail on its own. Noetherian guarantees that the
inductions over ideals terminate. Integral closedness fails for $\Z[\sqrt{-3}]$, which misses the integral element $\zeta_3$, and
\Cref{ex:ant:ex-sqrtminus3} shows that ideals cannot be cancelled there. ``Nonzero primes are maximal'' (dimension one) fails for $\Z[X]$, which is
Noetherian and integrally closed (a UFD) but has the chain of primes $0\subsetneq(X)\subsetneq(2,X)$; and in $\Z[X]$ the ideal $(4,X)$ is not a product of
primes at all: a prime factor would contain $(4,X)$, hence $2$ and $X$, hence be the maximal ideal $(2,X)$, but $(2,X)\ne(4,X)$ (as $2\notin(4,X)$) and
$(2,X)^k\subseteq(4,2X,X^2)\not\ni X$ for $k\ge2$. The rings of the paper, $\mathcal O_K$, $\Z[i]$, $\Z[\zeta_3]$, $\Z_3[\zeta_3]$,
satisfy all three.

\begin{theorem}\label{thm:ant:dedekind}
(i) For every number field $K$, $\mathcal O_K$ is a Dedekind domain; for a nonzero prime $\mathfrak p$ of $\mathcal O_K$, $\mathfrak p\cap\Z=p\Z$ for a
prime $p$, and $\mathcal O_K/\mathfrak p$ is a finite field of characteristic $p$. (ii) A principal ideal domain that is not a field is a Dedekind
domain. (iii) A Noetherian domain $R$ that is not a field is a Dedekind domain if and only if the localisation $R_{\mathfrak p}$ is a discrete
valuation ring for every nonzero prime $\mathfrak p$.
\end{theorem}

\begin{proof}
(i) $\mathcal O_K$ is a free $\Z$-module of finite rank (\Cref{thm:ant:integralbasis}), so every ideal, a subgroup, is finitely generated:
$\mathcal O_K$ is Noetherian. It is integrally closed (\Cref{prop:ant:integral}) and not a field ($\frac12\notin\mathcal O_K$). Let $\mathfrak p\ne0$
be prime and $0\ne\alpha\in\mathfrak p$, with minimal polynomial $X^d+\dots+c_1X+c_0$; then $c_0\ne0$ by irreducibility, and
$c_0=-\alpha(\alpha^{d-1}+\dots+c_1)\in\mathfrak p\cap\Z$. So $\mathfrak p\cap\Z$ is a nonzero prime ideal of $\Z$, equal to $p\Z$. Then
$\mathcal O_K/\mathfrak p$ is a quotient of $\mathcal O_K/p\mathcal O_K\simeq(\Z/p)^n$, a finite integral domain, hence a field: $\mathfrak p$ is
maximal. (ii) A PID is Noetherian, is a UFD and hence integrally closed (\Cref{ch:rings}), and a nonzero prime $(\varpi)$ is maximal: if
$(\varpi)\subseteq(a)$ then $\varpi=ab$, and $\varpi\mid a$ or $\varpi\mid b$, so $(a)=(\varpi)$ or $(a)=R$.
(iii) ``Only if'' is \Cref{prop:ant:crt}(iv) below, whose proof uses only \Cref{thm:ant:idealfactorisation}. ``If'': a discrete valuation ring is
integrally closed (\Cref{ch:modules}). If $x\in K$ lies in $R_{\mathfrak m}$ for every maximal ideal $\mathfrak m$, the ideal
$\{r\in R:rx\in R\}$ is contained in no maximal ideal (for each $\mathfrak m$ write $x=a/s$ with $a\in R$, $s\notin\mathfrak m$; then $s$ lies in the ideal but not in
$\mathfrak m$), so it is $R$, it contains $1$, and $x\in R$. Hence $R=\bigcap_{\mathfrak m}R_{\mathfrak m}$. Now if $x\in K$ is a root of a monic polynomial with
coefficients in $R$, it is integral over each $R_{\mathfrak m}\supseteq R$, so it lies in each $R_{\mathfrak m}$, hence in $R$: $R$ is integrally
closed. If $0\ne\mathfrak q\subsetneq\mathfrak m$ were primes, $\mathfrak qR_{\mathfrak m}\subsetneq\mathfrak mR_{\mathfrak m}$ would be two distinct
nonzero primes of the discrete valuation ring $R_{\mathfrak m}$, which has only one.
\end{proof}

\begin{lemma}\label{lem:ant:primeproducts}
Let $R$ be a Dedekind domain with fraction field $K$, and for a nonzero prime $\mathfrak p$ put $\mathfrak p^{-1}=\{x\in K:x\mathfrak p\subseteq R\}$.
(i) Every nonzero ideal of $R$ contains a product of nonzero prime ideals. (ii) For every nonzero ideal $\mathfrak a$, $\mathfrak a\mathfrak p^{-1}\ne\mathfrak a$;
in particular $\mathfrak p^{-1}\ne R$. (iii) $\mathfrak p\mathfrak p^{-1}=R$.
\end{lemma}

\begin{proof}
(i) If not, the set of nonzero ideals not containing such a product is nonempty and has a maximal element $\mathfrak a$ ($R$ is Noetherian). It is not
prime (a nonzero prime contains the product of one prime, itself), and not $R$ (which contains every product), so there are $b,c\notin\mathfrak a$ with
$bc\in\mathfrak a$; the ideals $\mathfrak a+(b)$ and $\mathfrak a+(c)$ are strictly larger, so by maximality they contain products of primes, and their product
$(\mathfrak a+(b))(\mathfrak a+(c))\subseteq\mathfrak a+(bc)=\mathfrak a$ contains the product of these two products: a contradiction.
(ii) First $\mathfrak p^{-1}\ne R$: take $0\ne a\in\mathfrak p$ and, by (i), primes with $\mathfrak p_1\cdots\mathfrak p_r\subseteq(a)\subseteq\mathfrak p$, $r$
minimal. Some $\mathfrak p_i\subseteq\mathfrak p$: otherwise choose $x_i\in\mathfrak p_i\setminus\mathfrak p$, and $x_1\cdots x_r\in\mathfrak p$ contradicts the
primality of $\mathfrak p$. Say $\mathfrak p_1\subseteq\mathfrak p$; as $\mathfrak p_1$ is a nonzero prime, it is maximal, so $\mathfrak p_1=\mathfrak p$. By
minimality of $r$, $\mathfrak p_2\cdots\mathfrak p_r\not\subseteq(a)$ (for $r=1$ read $\mathfrak p_2\cdots\mathfrak p_r=R$), so there is $b\in\mathfrak p_2\cdots\mathfrak p_r$
with $b\notin(a)$. Then $b\mathfrak p\subseteq\mathfrak p_1\mathfrak p_2\cdots\mathfrak p_r\subseteq(a)$, so $(b/a)\mathfrak p\subseteq R$, that is
$b/a\in\mathfrak p^{-1}$, while $b/a\notin R$ because $b\notin aR$.
Now let $\mathfrak a=(\alpha_1,\dots,\alpha_m)\ne0$ with $\mathfrak a\mathfrak p^{-1}=\mathfrak a$. For $x\in\mathfrak p^{-1}$, $x\alpha_i\in\mathfrak a$, so
$x\alpha_i=\sum_ja_{ij}\alpha_j$ with $a_{ij}\in R$, and the adjugate argument of \Cref{prop:ant:integral} gives $\det(xI-(a_{ij}))\alpha_j=0$ for all $j$, hence
$\det(xI-(a_{ij}))=0$ (as $R$ is a domain and some $\alpha_j\ne0$): $x$ is a root of a monic polynomial over $R$, and $x\in R$ because $R$ is integrally closed. So
$\mathfrak p^{-1}\subseteq R$, contradicting the first step. (iii) $\mathfrak p\mathfrak p^{-1}$ is an ideal of $R$ (by the definition of $\mathfrak p^{-1}$), it
contains $\mathfrak p$ (as $1\in\mathfrak p^{-1}$), and $\mathfrak p\mathfrak p^{-1}\ne\mathfrak p$ by (ii) with $\mathfrak a=\mathfrak p$. Since $\mathfrak p$ is
maximal, $\mathfrak p\mathfrak p^{-1}=R$.
\end{proof}

\begin{theorem}[unique factorisation of ideals]\label{thm:ant:idealfactorisation}
Let $R$ be a Dedekind domain with fraction field $K$.
\begin{enumerate}[label=(\roman*)]
\item Every nonzero ideal $\mathfrak a\ne R$ is a product $\mathfrak a=\mathfrak p_1^{a_1}\cdots\mathfrak p_r^{a_r}$ of distinct nonzero primes, $a_i\ge1$,
uniquely up to order. We write $v_{\mathfrak p}(\mathfrak a)$ for the exponent of $\mathfrak p$ ($0$ if $\mathfrak p$ does not occur).
\item A \emph{fractional ideal} is a nonzero finitely generated $R$-submodule of $K$, equivalently a set $x\mathfrak b$ with $x\in K^\times$ and
$\mathfrak b$ a nonzero ideal. The fractional ideals form an abelian group $J_R$ under multiplication, free on the nonzero primes; the inverse of
$\mathfrak a$ is $\mathfrak a^{-1}=\{x\in K:x\mathfrak a\subseteq R\}$.
\item For nonzero ideals, $\mathfrak a\mid\mathfrak b$ (that is, $\mathfrak b=\mathfrak a\mathfrak c$ for an ideal $\mathfrak c$) if and only if
$\mathfrak a\supseteq\mathfrak b$; $\mathfrak a+\mathfrak b$ and $\mathfrak a\cap\mathfrak b$ have exponents $\min$ and $\max$ of those of $\mathfrak a$, $\mathfrak b$.
\item For $\alpha\in K^\times$ put $v_{\mathfrak p}(\alpha)=v_{\mathfrak p}(\alpha R)$. Then $v_{\mathfrak p}(\alpha\beta)=v_{\mathfrak p}(\alpha)+v_{\mathfrak p}(\beta)$
and $v_{\mathfrak p}(\alpha+\beta)\ge\min(v_{\mathfrak p}(\alpha),v_{\mathfrak p}(\beta))$: $v_{\mathfrak p}$ is a discrete valuation of $K$.
\end{enumerate}
\end{theorem}

\begin{proof}
(i) Existence: otherwise let $\mathfrak a$ be maximal among the nonzero proper ideals that are not products of primes, and $\mathfrak p\supseteq\mathfrak a$ a
maximal ideal. Then $\mathfrak a\subseteq\mathfrak a\mathfrak p^{-1}\subseteq\mathfrak p\mathfrak p^{-1}=R$, with $\mathfrak a\ne\mathfrak a\mathfrak p^{-1}$
(\Cref{lem:ant:primeproducts}). If $\mathfrak a\mathfrak p^{-1}=R$ then $\mathfrak a=\mathfrak a\mathfrak p^{-1}\mathfrak p=\mathfrak p$, a prime; otherwise
$\mathfrak a\mathfrak p^{-1}=\mathfrak p_1\cdots\mathfrak p_s$ by maximality, and $\mathfrak a=\mathfrak p_1\cdots\mathfrak p_s\mathfrak p$. Either way a contradiction.
Uniqueness: if $\mathfrak p_1\cdots\mathfrak p_r=\mathfrak q_1\cdots\mathfrak q_s$, then $\mathfrak p_1$ contains the right side, hence some $\mathfrak q_j$, and
$\mathfrak p_1=\mathfrak q_j$ by maximality; multiply by $\mathfrak p_1^{-1}$ (\Cref{lem:ant:primeproducts}(iii)) and induct.
(ii) A fractional ideal $\mathfrak a$ with generators $x_i=a_i/s$ satisfies $s\mathfrak a\subseteq R$, so $\mathfrak a=s^{-1}\mathfrak b$; writing $sR$ and
$\mathfrak b$ as products of primes, $\mathfrak a=\prod\mathfrak p^{a_{\mathfrak p}}$ with $a_{\mathfrak p}\in\Z$, almost all zero, and the exponents are unique
by (i) after multiplying by a common denominator. Put $\mathfrak a'=\prod\mathfrak p^{-a_{\mathfrak p}}$; then $\mathfrak a\mathfrak a'=R$, so $\mathfrak a'\subseteq
\mathfrak c:=\{x:x\mathfrak a\subseteq R\}$, and $\mathfrak c=\mathfrak c\mathfrak a\mathfrak a'\subseteq\mathfrak a'$. (iii) If $\mathfrak a\supseteq\mathfrak b$,
then $\mathfrak c=\mathfrak a^{-1}\mathfrak b\subseteq\mathfrak a^{-1}\mathfrak a=R$ is an ideal with $\mathfrak b=\mathfrak a\mathfrak c$; the converse is clear. So
$\mathfrak a\supseteq\mathfrak b$ iff $v_{\mathfrak p}(\mathfrak a)\le v_{\mathfrak p}(\mathfrak b)$ for all $\mathfrak p$; $\mathfrak a+\mathfrak b$ is the smallest
ideal containing both, $\mathfrak a\cap\mathfrak b$ the largest contained in both, whence the exponents. (iv) The first identity is (i); and
$(\alpha+\beta)\subseteq(\alpha)+(\beta)$, whose exponent at $\mathfrak p$ is the minimum, by (iii) extended to fractional ideals by scaling.
\end{proof}

Fractional ideals are introduced so that ideals can be inverted, as rational numbers invert integers. For $R=\Z$ the fractional ideals are the groups
$r\Z$, $r\in\Q_{>0}$, and $J_\Z\simeq\Q_{>0}^\times$, free on the primes: $\frac{12}5\Z=(2)^2(3)(5)^{-1}$. In $\Z[\sqrt{-5}]$, with the primes of
\Cref{ex:ant:sqrtminus5}, $\mathfrak p_2^{-1}=\frac12\mathfrak p_2$ (because $\mathfrak p_2\cdot\frac12\mathfrak p_2=\frac12(2)=R$), a fractional ideal that is not
of the form $xR$. The valuation $v_{\mathfrak p}$ of (iv) is the exponent of $\mathfrak p$ in $\alpha R$; for $\mathcal O_K=\Z[i]$ and $\mathfrak p=(1+i)$, for
instance, $v_{\mathfrak p}(2)=2$ since $2=-i(1+i)^2$, and $v_{\mathfrak p}(3)=0$.

\begin{proposition}[the norm of an ideal]\label{prop:ant:idealnorm}
Let $K$ be a number field. For a nonzero ideal $\mathfrak a\subseteq\mathcal O_K$ the \emph{norm} $N\mathfrak a=\#(\mathcal O_K/\mathfrak a)$ is finite, and:
\begin{enumerate}[label=(\roman*)]
\item $N(\mathfrak a\mathfrak b)=N\mathfrak a\,N\mathfrak b$; for a prime $\mathfrak p$ and $k\ge0$, $\mathfrak p^k/\mathfrak p^{k+1}\simeq\mathcal O_K/\mathfrak p$ as $\mathcal O_K$-modules;
\item $N(\alpha\mathcal O_K)=\abs{N_{K/\Q}(\alpha)}$ for $\alpha\ne0$;
\item if $\mathfrak p\cap\Z=p\Z$ then $\mathcal O_K/\mathfrak p\simeq\F_{p^f}$, where $f=f(\mathfrak p\mid p)=[\mathcal O_K/\mathfrak p:\F_p]$, and $N\mathfrak p=p^f$;
\item an ideal of prime norm is prime; $N\mathfrak a\in\mathfrak a$; and if $\gcd(N\mathfrak a,m)=1$ then $\mathfrak a+m\mathcal O_K=\mathcal O_K$;
\item the discriminant of a $\Z$-basis of $\mathfrak a$ is $(N\mathfrak a)^2d_K$.
\end{enumerate}
\end{proposition}

\begin{proof}
$\mathfrak a$ contains a nonzero integer $m$ (the constant term $c_0$ of the proof of \Cref{thm:ant:dedekind}), and $\mathcal O_K/m\mathcal O_K\simeq(\Z/m)^n$
is finite. (i) The module $M=\mathfrak a/\mathfrak a\mathfrak p$ is a vector space over $\mathcal O_K/\mathfrak p$; it is nonzero
(\Cref{thm:ant:idealfactorisation}) and has no subspace other than $0$ and $M$, because an ideal $\mathfrak b$ with $\mathfrak a\mathfrak p\subseteq
\mathfrak b\subseteq\mathfrak a$ satisfies $\mathfrak a\mid\mathfrak b\mid\mathfrak a\mathfrak p$, so $\mathfrak b\in\{\mathfrak a,\mathfrak a\mathfrak p\}$. Hence $M$ is
one-dimensional, $M\simeq\mathcal O_K/\mathfrak p$, and $N(\mathfrak a\mathfrak p)=N\mathfrak a\cdot\#(\mathfrak a/\mathfrak a\mathfrak p)=N\mathfrak a\,N\mathfrak p$;
induction on the number of prime factors of $\mathfrak b$ gives the claim. (ii) $[\mathcal O_K:\alpha\mathcal O_K]$ is $\abs{\det}$ of the matrix of
$m_\alpha$ on a $\Z$-basis (\Cref{ch:modules}), which is $\abs{N(\alpha)}$. (iii) $\mathcal O_K/\mathfrak p$ is a finite field
(\Cref{thm:ant:dedekind}) containing $\Z/p$. (iv) If $N\mathfrak a=p$ is prime, $\mathfrak a\ne\mathcal O_K$, and a factorisation
$\mathfrak a=\mathfrak b\mathfrak c$ into proper ideals would give $N\mathfrak a=N\mathfrak b\,N\mathfrak c$ composite. The additive group
$\mathcal O_K/\mathfrak a$ has order $N\mathfrak a$, so $N\mathfrak a\cdot1\in\mathfrak a$; if $uN\mathfrak a+vm=1$, then $1\in\mathfrak a+m\mathcal O_K$.
(v) is \Cref{thm:ant:integralbasis}(iii), with $[\mathcal O_K:\mathfrak a]=N\mathfrak a$.
\end{proof}

\begin{proposition}[Chinese remainder theorem, generators, localisation]\label{prop:ant:crt}
In (ii) to (iv), $R$ is a Dedekind domain.
\begin{enumerate}[label=(\roman*)]
\item In a commutative ring $R$, let $\mathfrak a_1$, \dots, $\mathfrak a_r$ be pairwise coprime ideals, that is,
$\mathfrak a_i+\mathfrak a_j=R$ for $i\ne j$. Then
\[
\mathfrak a_1\cdots\mathfrak a_r=\bigcap_i\mathfrak a_i\quad\text{and}\quad R/\mathfrak a_1\cdots\mathfrak a_r\simeq\prod_iR/\mathfrak a_i.
\] In a Dedekind domain, powers of
distinct primes are coprime, so $R/\mathfrak a\simeq\prod_iR/\mathfrak p_i^{a_i}$ for $\mathfrak a=\prod\mathfrak p_i^{a_i}$, and $R/\mathfrak p^a$ has length
$a$, with composition factors $\mathfrak p^k/\mathfrak p^{k+1}\simeq R/\mathfrak p$.
\item In a Dedekind domain every nonzero ideal $\mathfrak a$ is generated by two elements, the first an arbitrary nonzero $\alpha\in\mathfrak a$.
\item A Dedekind domain with finitely many primes is a principal ideal domain; a Dedekind domain is a UFD if and only if it is a PID.
\item For a nonzero prime $\mathfrak p$ of a Dedekind domain $R$, $R_{\mathfrak p}$ is a discrete valuation ring with maximal ideal
$\mathfrak pR_{\mathfrak p}$, residue field $R/\mathfrak p$, and valuation $v_{\mathfrak p}$.
\end{enumerate}
\end{proposition}

\begin{proof}
(i) For two ideals, $1=a+b$ with $a\in\mathfrak a_1$, $b\in\mathfrak a_2$; then $\mathfrak a_1\cap\mathfrak a_2=(\mathfrak a_1\cap\mathfrak a_2)(a+b)\subseteq\mathfrak a_1
\mathfrak a_2$, and $x\mapsto(x\bmod\mathfrak a_1,x\bmod\mathfrak a_2)$ is surjective since $(x,y)$ is the image of $xb+ya$. Induct, using that
$\mathfrak a_1$ is coprime to $\mathfrak a_2\cdots\mathfrak a_r$ (expand $1=\prod_{j\ge2}(a_j+b_j)$). In a Dedekind domain, a maximal ideal containing
$\mathfrak p^a+\mathfrak q^b$ contains $\mathfrak p$ and $\mathfrak q$, so equals both; hence distinct primes give coprime powers. The filtration
$R\supseteq\mathfrak p\supseteq\dots\supseteq\mathfrak p^a$ has the quotients of \Cref{prop:ant:idealnorm}(i), whose proof holds in any Dedekind domain.
(ii) Let $\mathfrak p_1,\dots,\mathfrak p_r$ be the primes dividing $\alpha R$, $e_i=v_{\mathfrak p_i}(\mathfrak a)$, and choose $x_i\in\mathfrak p_i^{e_i}\setminus
\mathfrak p_i^{e_i+1}$. By (i) there is $\beta\in R$ with $\beta\equiv x_i\pmod{\mathfrak p_i^{e_i+1}}$ for all $i$; then $v_{\mathfrak p_i}(\beta)=e_i$. Every
prime dividing $\mathfrak a$ divides $\alpha R$, so $\beta\in\mathfrak a$ (\Cref{thm:ant:idealfactorisation}(iii)), and $(\alpha,\beta)=\alpha R+\beta R$
has exponent $\min(v_{\mathfrak p_i}(\alpha),e_i)=e_i$ at $\mathfrak p_i$ and $0$ elsewhere: it is $\mathfrak a$.
(iii) If the primes are $\mathfrak p_1,\dots,\mathfrak p_r$, choose by (i) $\pi_i\equiv x_i\pmod{\mathfrak p_i^2}$ with $x_i\in\mathfrak p_i\setminus\mathfrak p_i^2$
and $\pi_i\equiv1\pmod{\mathfrak p_j}$ for $j\ne i$ (the ideals $\mathfrak p_i^2$ and $\mathfrak p_j$, $j\ne i$, are pairwise coprime). Then $\pi_i\in\mathfrak p_i$,
$\pi_i\notin\mathfrak p_i^2$ (since $\pi_i-x_i\in\mathfrak p_i^2$ and $x_i\notin\mathfrak p_i^2$) and $\pi_i\notin\mathfrak p_j$; so $v_{\mathfrak p_i}(\pi_i)=1$ and
$v_{\mathfrak p_j}(\pi_i)=0$, and as there are no other primes, $\pi_iR=\mathfrak p_i$ by \Cref{thm:ant:idealfactorisation}(i). Every ideal, a product of primes, is
then principal. A PID is a UFD
(\Cref{ch:rings}). Conversely, in a UFD a nonzero prime $\mathfrak p$ contains a nonzero element, hence one of its irreducible factors $\varpi$, which
is a prime element; $(\varpi)\subseteq\mathfrak p$ is a nonzero prime ideal, hence maximal, so $\mathfrak p=(\varpi)$, and all ideals are principal.
(iv) A nonzero ideal $A$ of $R_{\mathfrak p}$ is $(A\cap R)R_{\mathfrak p}$, and for a prime $\mathfrak q\ne\mathfrak p$, $\mathfrak qR_{\mathfrak p}=R_{\mathfrak p}$
($\mathfrak q\not\subseteq\mathfrak p$). So $A=\mathfrak p^aR_{\mathfrak p}$ with $a=v_{\mathfrak p}(A\cap R)$. If $\varpi\in\mathfrak p\setminus\mathfrak p^2$ then
$\varpi R=\mathfrak p\mathfrak b$ with $\mathfrak b$ prime to $\mathfrak p$, so $\varpi R_{\mathfrak p}=\mathfrak pR_{\mathfrak p}$, and every nonzero ideal is
$\varpi^aR_{\mathfrak p}$: $R_{\mathfrak p}$ is a local principal ideal domain, not a field, that is, a discrete valuation ring (\Cref{ch:modules}). Its
valuation of $\alpha$ is the exponent $a$ with $\alpha R_{\mathfrak p}=\varpi^aR_{\mathfrak p}$, which is $v_{\mathfrak p}(\alpha)$. The map $R/\mathfrak p\to
R_{\mathfrak p}/\mathfrak pR_{\mathfrak p}$ is injective ($\mathfrak pR_{\mathfrak p}\cap R=\mathfrak p$) and surjective (an $s\notin\mathfrak p$ is invertible
modulo the maximal ideal $\mathfrak p$).
\end{proof}

\begin{theorem}[the class group]\label{thm:ant:classgroup}
Let $K$ be a number field, $J_K$ its group of fractional ideals and $P_K$ the subgroup of principal fractional ideals $\alpha\mathcal O_K$. The
\emph{class group} is $\Cl(K)=J_K/P_K$ and the \emph{class number} is $h_K=\#\Cl(K)$; $h_K=1$ iff $\mathcal O_K$ is a PID iff it is a UFD
(\Cref{prop:ant:crt}(iii)). (Minkowski.) $\Cl(K)$ is finite: every class contains an integral ideal of norm at most
$M_K=(4/\pi)^{r_2}\,(n!/n^n)\,\abs{d_K}^{1/2}$. Moreover $\abs{d_K}>1$ for $K\ne\Q$.
\end{theorem}

\begin{proof}[Reference]
The definitions and the equivalences are proved above. We do not prove Minkowski's theorem; see \cite[Chapter I, Sections 5 and 6]{Neukirch} or
\cite[Chapter 5]{Marcus77}. The idea: the embeddings $(\sigma_1,\dots,\sigma_{r_1+r_2})$ map an ideal $\mathfrak a$ to a lattice in
$\R^{r_1}\times\C^{r_2}\simeq\R^n$ of covolume $2^{-r_2}\abs{d_K}^{1/2}N\mathfrak a$ (the Gram determinant of \Cref{thm:ant:integralbasis}, up to the
factor $2^{-r_2}$ from the coordinates of $\C$), and Minkowski's convex body theorem gives $\alpha\in\mathfrak a$ with $\abs{N(\alpha)}\le
M_KN\mathfrak a$. The inequality $\abs{d_K}>1$ follows: an integral ideal has norm at least $1$, so $M_K\ge1$, that is
$\abs{d_K}^{1/2}\ge(\pi/4)^{r_2}n^n/n!\ge(\pi/4)^{n/2}n^n/n!$, and the last expression is $\pi/2>1$ for $n=2$ and increases with $n$ (the ratio of consecutive
terms is $(\pi/4)^{1/2}(1+\frac1n)^n\ge2(\pi/4)^{1/2}>1$).
\end{proof}

\begin{example}[class numbers of small quadratic fields]\label{ex:ant:classnumbers}
For a quadratic field $n!/n^n=\frac12$, so $M_K=\frac2\pi\abs D^{1/2}$ if $D<0$ and $M_K=\frac12D^{1/2}$ if $D>0$. Hence (all values computed with PARI's
\texttt{quadclassunit}, which also confirms the class numbers):
\begin{center}\small
\begin{tabular}{lccccccc}
\toprule
$K$ & $\Q(\sqrt{-3})$ & $\Q(i)$ & $\Q(\sqrt{-7})$ & $\Q(\sqrt5)$ & $\Q(\sqrt2)$ & $\Q(\sqrt{-11})$ & $\Q(\sqrt{-5})$\\
$D$ & $-3$ & $-4$ & $-7$ & $5$ & $8$ & $-11$ & $-20$\\
$M_K$ & $1.10$ & $1.27$ & $1.68$ & $1.12$ & $1.41$ & $2.11$ & $2.85$\\
$h_K$ & $1$ & $1$ & $1$ & $1$ & $1$ & $1$ & $2$\\
\bottomrule
\end{tabular}
\end{center}
In the first five cases $M_K<2$, so every class contains the ideal of norm $1$, which is $\mathcal O_K$. For $\Q(\sqrt{-11})$, $M_K>2$, but there is
no ideal of norm $2$, because $2$ is inert (\Cref{thm:ant:quadsplit}: $-11\equiv5\pmod8$); so again $h=1$. For $\Q(\sqrt{-5})$ the ideals of norm $\le2$ are
$\mathcal O_K$ and $\mathfrak p_2$ of \Cref{ex:ant:sqrtminus5} (the norm $2$ ideals are prime, \Cref{prop:ant:idealnorm}(iv), and $\mathfrak p_2$ is the only
prime above the ramified prime $2$), which is not principal, and $\mathfrak p_2^2=(2)$. Every class contains one of these two ideals, and they lie in different
classes, so there are exactly two classes, and $h=2$. Also $h(\Q(\sqrt{-23}))=3$
(stated; PARI). For $\Z[i]$ and $\Z[\zeta_3]$ the class number one is also a consequence of the Euclidean algorithm (\Cref{ch:rings}).
\end{example}

\begin{example}[$\Z\lbrack\sqrt{-5}\rbrack$]\label{ex:ant:sqrtminus5}
$R=\Z[\sqrt{-5}]$ is the ring of integers of $\Q(\sqrt{-5})$ ($-5\equiv3\pmod4$), with $D=-20$. In $R$, $6=2\cdot3=(1+\sqrt{-5})(1-\sqrt{-5})$
are two factorisations into irreducibles, no factor of one being associate to a factor of the other (\Cref{ch:rings}). Put
$\mathfrak p_2=(2,1+\sqrt{-5})$, $\mathfrak p_3=(3,1+\sqrt{-5})$, $\bar{\mathfrak p}_3=(3,1-\sqrt{-5})$. The ring homomorphisms $R\to\F_2$,
$a+b\sqrt{-5}\mapsto a+b$, and $R\to\F_3$, $a+b\sqrt{-5}\mapsto a\mp b$ (well defined because $(\pm1)^2+5\equiv0$ modulo $2$, resp. $3$), are surjective
with kernels $\mathfrak p_2$, $\mathfrak p_3$, $\bar{\mathfrak p}_3$ (an element $a+b\sqrt{-5}$ with $a\equiv b\pmod 3$ is $b(1+\sqrt{-5})+(a-b)\in\mathfrak p_3$,
and similarly for the others). So these are primes of norms $2,3,3$. Next, $\mathfrak p_2^2=(4,2+2\sqrt{-5},-4+2\sqrt{-5})$ contains
$6=(2+2\sqrt{-5})-(-4+2\sqrt{-5})$, hence $2$, and is contained in $(2)$: $\mathfrak p_2^2=(2)$. Similarly $\mathfrak p_3\bar{\mathfrak p}_3=(9,3+3\sqrt{-5},
3-3\sqrt{-5},6)=(3)$. Finally $1+\sqrt{-5}\in\mathfrak p_2\cap\mathfrak p_3=\mathfrak p_2\mathfrak p_3$ and both ideals have norm $6$, so
$(1+\sqrt{-5})=\mathfrak p_2\mathfrak p_3$, and likewise $(1-\sqrt{-5})=\mathfrak p_2\bar{\mathfrak p}_3$. Thus
\[
(6)=\mathfrak p_2^2\,\mathfrak p_3\,\bar{\mathfrak p}_3,
\]
and the two factorisations of $6$ are two groupings of the same four prime ideals. None of $\mathfrak p_2,\mathfrak p_3,\bar{\mathfrak p}_3$ is principal: a
generator would have norm $a^2+5b^2\in\{2,3\}$, which is impossible. This is the ideal-theoretic repair of \Cref{ch:rings}, and the pattern is
the one predicted by \Cref{thm:ant:dedekindkummer}: $X^2+5\equiv(X+1)^2\pmod2$ and $\equiv(X-1)(X+1)\pmod3$.
\end{example}

\section{Splitting of primes, ramification and Frobenius}\label{sec:ant:splitting}

This section explains how a prime of $K$ factors in an extension $L$, defines the decomposition and inertia groups and the Frobenius element,
and proves, given the tame relation (sketched here, proved in \Cref{ch:local}), the form of it that the paper uses. The paper's Section 1.1 is written in terms of $I_p$, $\Frob_p$ and
$V^{I_p}$ (\Cref{prop:ant:frobGQ}); its local $L$-factors use the geometric Frobenius (\Cref{conv:ant:frobenius}); Lemma lem:old uses the
Frobenius at unramified primes and Chebotarev (\Cref{prop:ant:frobenius}, \Cref{sec:ant:chebotarev}); Lemmas lem:IV and lem:euler use that
a Frobenius at a tamely ramified prime $q$ conjugates inertia to its $q$-th power (\Cref{cor:ant:frobsigma}). The results on Frobenius elements
are the black boxes that \Cref{ch:galois} states.

\begin{theorem}[$\sum e_if_i=n$]\label{thm:ant:efg}
Let $L/K$ be an extension of number fields of degree $n$, $\mathfrak p$ a nonzero prime of $\mathcal O_K$, and $\mathfrak p\mathcal O_L=
\mathfrak P_1^{e_1}\cdots\mathfrak P_g^{e_g}$ with distinct primes $\mathfrak P_i$ (\Cref{thm:ant:idealfactorisation}). The $\mathfrak P_i$ are exactly the
primes $\mathfrak P$ of $\mathcal O_L$ with $\mathfrak P\cap\mathcal O_K=\mathfrak p$ (``$\mathfrak P$ lies above $\mathfrak p$'', $\mathfrak P\mid\mathfrak p$). With the
\emph{ramification index} $e_i=e(\mathfrak P_i\mid\mathfrak p)$ and the \emph{residue degree} $f_i=f(\mathfrak P_i\mid\mathfrak p)=[\mathcal O_L/\mathfrak P_i:
\mathcal O_K/\mathfrak p]$,
\[
\sum_{i=1}^ge_if_i=n.
\]
$\mathfrak p$ is \emph{ramified} in $L$ if some $e_i>1$ and \emph{unramified} otherwise; \emph{totally ramified} if $g=1$, $e_1=n$; \emph{inert} if
$g=1$, $f_1=n$; \emph{split completely} if $g=n$. For a quadratic $L/K$ each $\mathfrak p$ is exactly one of ramified, inert, split. In a tower
$K\subseteq M\subseteq L$ with $\mathfrak P\mid\mathfrak q\mid\mathfrak p$: $e(\mathfrak P\mid\mathfrak p)=e(\mathfrak P\mid\mathfrak q)e(\mathfrak q\mid\mathfrak p)$ and
$f(\mathfrak P\mid\mathfrak p)=f(\mathfrak P\mid\mathfrak q)f(\mathfrak q\mid\mathfrak p)$.
\end{theorem}

\begin{proof}
A prime $\mathfrak P$ of $\mathcal O_L$ contains $\mathfrak p\mathcal O_L$ iff $\mathfrak P\cap\mathcal O_K\supseteq\mathfrak p$ iff $\mathfrak P\cap\mathcal O_K=\mathfrak p$
(the left side is a prime of $\mathcal O_K$, nonzero as it contains $\mathfrak P\cap\Z$, hence maximal), and $\mathfrak P\supseteq\mathfrak p\mathcal O_L$
iff $\mathfrak P$ divides $\mathfrak p\mathcal O_L$ (\Cref{thm:ant:idealfactorisation}(iii)). Let $\kappa=\mathcal O_K/\mathfrak p$. By
\Cref{prop:ant:crt}(i), $\mathcal O_L/\mathfrak p\mathcal O_L\simeq\prod_i\mathcal O_L/\mathfrak P_i^{e_i}$, and $\mathcal O_L/\mathfrak P_i^{e_i}$ has a filtration
with $e_i$ quotients $\mathfrak P_i^k/\mathfrak P_i^{k+1}\simeq\mathcal O_L/\mathfrak P_i$, each of dimension $f_i$ over $\kappa$; so $\dim_\kappa\mathcal O_L/
\mathfrak p\mathcal O_L=\sum e_if_i$. To see that this dimension is $n$, let $S=\mathcal O_K\setminus\mathfrak p$ and $A=S^{-1}\mathcal O_K$, a discrete
valuation ring (\Cref{prop:ant:crt}(iv)), hence a PID. $B=S^{-1}\mathcal O_L$ is a finitely generated torsion-free $A$-module ($\mathcal O_L$ is finitely
generated over $\Z$), hence free (\Cref{ch:modules}), of rank $n$ because $B\otimes_AK=L$ (\Cref{prop:ant:integral}). So $B/\mathfrak pB\simeq\kappa^n$.
Finally $\mathcal O_L/\mathfrak p\mathcal O_L\to B/\mathfrak pB$ is an isomorphism: every $s\in S$ has an inverse $s'$ modulo $\mathfrak p$, so $x/s\equiv s'x$
(surjectivity), and if $x\in\mathcal O_L$ and $sx\in\mathfrak p\mathcal O_L$ with $s\in S$, then $x=s'sx-(ss'-1)x\in\mathfrak p\mathcal O_L$ (injectivity). For
$K=\Q$ this is simply $\mathcal O_L\simeq\Z^n$. A quadratic $L/K$ has $\sum e_if_i=2$, so $(e,f,g)$ is $(2,1,1)$, $(1,2,1)$ or $(1,1,2)$. In a tower,
substitute $\mathfrak q\mathcal O_L=\prod\mathfrak P^{e(\mathfrak P\mid\mathfrak q)}$ into $\mathfrak p\mathcal O_M=\prod\mathfrak q^{e(\mathfrak q\mid\mathfrak p)}$ and use
uniqueness; the residue degrees multiply as degrees of the tower of residue fields.
\end{proof}

\begin{theorem}[Dedekind-Kummer]\label{thm:ant:dedekindkummer}
Let $K=\Q(\theta)$ with $\theta\in\mathcal O_K$ of minimal polynomial $g\in\Z[X]$, and let $p$ be a prime with $p\nmid[\mathcal O_K:\Z[\theta]]$ (for
instance $p^2\nmid\disc(g)$, by \Cref{thm:ant:integralbasis}(iii)). Let $\bar g=\bar g_1^{e_1}\cdots\bar g_r^{e_r}$ be the factorisation of $g\bmod p$
into distinct monic irreducibles of $\F_p[X]$, and $g_i\in\Z[X]$ monic lifts of $\bar g_i$. Then $\mathfrak P_i=(p,g_i(\theta))$ are distinct primes,
$f(\mathfrak P_i\mid p)=\deg\bar g_i$, and
\[
p\mathcal O_K=\mathfrak P_1^{e_1}\cdots\mathfrak P_r^{e_r}.
\]
The same holds for a relative extension $L=K(\theta)$, $\theta\in\mathcal O_L$, and a prime $\mathfrak p$ of $\mathcal O_K$ prime to the conductor
$\{x\in\mathcal O_L:x\mathcal O_L\subseteq\mathcal O_K[\theta]\}$ (stated; \cite[Chapter I, Proposition 8.3]{Neukirch}).
\end{theorem}

\begin{proof}
Let $m=[\mathcal O_K:\Z[\theta]]$, prime to $p$. The inclusion induces $\Z[\theta]/p\Z[\theta]\to\mathcal O_K/p\mathcal O_K$. It is injective: if
$x\in\Z[\theta]$ and $x=py$ with $y\in\mathcal O_K$, then $my\in\Z[\theta]$, and with $um+vp=1$, $y=u(my)+vx\in\Z[\theta]$. Both sides have $p^n$
elements, so it is an isomorphism, and $\Z[\theta]/p\simeq\Z[X]/(p,g)\simeq\F_p[X]/(\bar g)$. By the Chinese remainder theorem in $\F_p[X]$,
\[
A:=\mathcal O_K/p\mathcal O_K\simeq\F_p[X]/(\bar g)\simeq\prod_i\F_p[X]/(\bar g_i^{e_i}),
\]
whose maximal ideals are the $(\bar g_i)$, with residue fields $\F_p[X]/(\bar g_i)$ of degree $\deg\bar g_i$. The primes of $\mathcal O_K$ containing $p$
correspond to the maximal ideals of $A$; the one corresponding to $(\bar g_i)$ is the preimage of $(\bar g_i(X))$, which is $\mathfrak P_i=(p,g_i(\theta))$.
So these are the primes above $p$, distinct, with $f(\mathfrak P_i\mid p)=\deg\bar g_i$. In $A$, $\prod_i(\bar g_i)^{e_i}=(\bar g)=0$, so
$\prod_i\mathfrak P_i^{e_i}\subseteq p\mathcal O_K$; writing $p\mathcal O_K=\prod_i\mathfrak P_i^{a_i}$, this says $a_i\le e_i$ (\Cref{thm:ant:idealfactorisation}(iii)),
while $\sum a_if_i=n=\deg g=\sum e_if_i$ (\Cref{thm:ant:efg}). Hence $a_i=e_i$.
\end{proof}

\begin{example}[worked factorisations]\label{ex:ant:dkexamples}
All of the following factorisations were checked with PARI.
(a) $\Q(i)$, $g=X^2+1$, index $1$. Modulo $2$, $g\equiv(X+1)^2$: $2\Z[i]=(2,1+i)^2=(1+i)^2$, and indeed $2=-i(1+i)^2$. Modulo $5$, $g\equiv(X-2)(X+2)$:
$5\Z[i]=(5,i-2)(5,i+2)=(2-i)(2+i)$. Modulo $3$, $g$ is irreducible: $3$ is inert, with residue field $\F_9$.
(b) $\Q(\zeta_3)$, $g=X^2+X+1$, index $1$. Modulo $3$, $g\equiv(X-1)^2$: $3\Z[\zeta_3]=(3,\zeta_3-1)^2=(1-\zeta_3)^2$. Modulo $7$, $g\equiv(X-2)(X-4)$:
$7=(7,\zeta_3-2)(7,\zeta_3-4)$, and these ideals are $(2-\zeta_3)$ and $(2-\zeta_3^2)$, of norm $4+2+1=7$, with $(2-\zeta_3)(2-\zeta_3^2)=7$ (note
$\zeta_3^2=-1-\zeta_3\equiv4$ when $\zeta_3\equiv2$). Modulo $2$ and $5$, $g$ is irreducible: $2$ and $5$ are inert, with residue fields $\F_4$ and $\F_{25}$.
(c) $\Q(\sqrt{-5})$: \Cref{ex:ant:sqrtminus5}.
(d) The index matters. In $K=\Q(\sqrt{-7})$, $\mathcal O_K=\Z[\theta]$ with $\theta=\frac{1+\sqrt{-7}}2$, $g=X^2-X+2$, and $[\mathcal O_K:\Z[\sqrt{-7}]]=2$. Modulo $2$,
$X^2+7\equiv(X+1)^2$ would wrongly suggest that $2$ ramifies; but $X^2-X+2\equiv X(X-1)\pmod2$, so $2$ splits:
$2=\frac{1+\sqrt{-7}}2\cdot\frac{1-\sqrt{-7}}2$, and indeed $2\nmid d_K=-7$ (\Cref{thm:ant:discram}).
(e) A non-Galois cubic. $K=\Q(2^{1/3})$ has $\mathcal O_K=\Z[2^{1/3}]$ and $d_K=\disc(X^3-2)=-108=-2^2\cdot3^3$ (stated; PARI's \texttt{nfbasis}; the
index is prime to $5$, $7$ and $31$ since $\disc$ is divisible only by $2$ and $3$). Modulo $5$, $X^3-2\equiv(X+2)(X^2+3X+4)$ with the quadratic
irreducible, so $5\mathcal O_K=\mathfrak P_1\mathfrak P_2$ with $f=1$ and $f=2$; modulo $7$, $X^3-2$ is irreducible ($2$ is not a cube modulo $7$), so $7$ is
inert; modulo $31$ it has three roots ($4,7,20$), so $31$ splits completely. At $2$ and $3$ PARI finds $2\mathcal O_K=\mathfrak P^3$ and $3\mathcal O_K=\mathfrak P^3$,
totally ramified ($X^3-2\equiv X^3$, resp. $(X+1)^3$, consistently with this, although the theorem applies at $2$ and $3$ only once the index is
known to be $1$). Unequal residue degrees above $5$ cannot occur in a Galois extension (\Cref{prop:ant:galoissplitting}).
\end{example}

\begin{theorem}[Dedekind's discriminant theorem]\label{thm:ant:discram}
A prime $p$ ramifies in the number field $K$ if and only if $p\mid d_K$. Hence only finitely many primes ramify in $K$, and at least one does if
$K\ne\Q$ (\Cref{thm:ant:classgroup}). For an extension $L/K$, a prime $\mathfrak p$ of $K$ ramifies in $L$ iff $\mathfrak p$ divides the relative
discriminant ideal $\mathfrak d_{L/K}$ (stated; \cite[Chapter III, Section 2]{Neukirch}).
\end{theorem}

\begin{proof}
Let $\beta_1,\dots,\beta_n$ be an integral basis and $A=\mathcal O_K/p\mathcal O_K$, with $\F_p$-basis $\bar\beta_i$. The matrix of multiplication by
$\bar x$ on $A$ is the reduction of the matrix of $m_x$, so $\Tr_{A/\F_p}(\bar x)=\Tr_{K/\Q}(x)\bmod p$ and $d_K\bmod p=\det(\Tr_{A/\F_p}(\bar\beta_i
\bar\beta_j))$. Thus $p\mid d_K$ iff the trace form of the $\F_p$-algebra $A$ is degenerate. By \Cref{prop:ant:crt}, $A\simeq\prod_iA_i$ with
$A_i=\mathcal O_K/\mathfrak P_i^{e_i}$, and the trace form of $A$ is the orthogonal sum of those of the $A_i$. If some $e_i>1$, $A_i$ contains a
nonzero nilpotent $x$ (the image of an element $\xi\in\mathfrak P_i\setminus\mathfrak P_i^{e_i}$, nonzero in $A_i$, with $\xi^{e_i}\in\mathfrak P_i^{e_i}$); then $xy$
is nilpotent for all $y$ ($A$ is commutative), multiplication by a nilpotent element is a nilpotent endomorphism, a nilpotent endomorphism has trace $0$
(its only eigenvalue is $0$), and $x$ lies in the kernel of the trace form. If all $e_i=1$, each $A_i=\F_{p^{f_i}}$ is a field, whose trace form over
$\F_p$ is nondegenerate: $\Tr(x)=x+x^p+\dots+x^{p^{f-1}}$ (the sum of the conjugates, as in \Cref{prop:ant:tracenorm}(i), the Galois group being
generated by $x\mapsto x^p$, \Cref{ch:rings}) is a polynomial function of degree $p^{f-1}<p^f$, hence not identically zero (a nonzero polynomial of degree
$<p^f$ cannot vanish at all $p^f$ elements of the field), and if $\Tr(xy)=0$ for all $y$ with $x\ne0$ then, as $xy$ runs over the whole field when $y$ does,
$\Tr$ would vanish identically.
\end{proof}

For example, for $K=\Q(i)$ and $p=2$ the algebra is $A=\F_2[X]/(X+1)^2$, the nilpotent is $x=1+i$, and the Gram matrix of the trace form in the basis $1,i$ is
$\sm200{-2}$, zero modulo $2$; for $p=3$, $A=\F_9$ and the same matrix is invertible modulo $3$. The theorem says that this mechanism, nilpotents killing the
trace form, is the whole story.

From now on in this section $L/K$ is a \emph{Galois} extension of number fields with group $G$, $\mathfrak p$ is a nonzero prime of $K$ and
$\mathfrak P\mid\mathfrak p$. Write $\kappa(\mathfrak p)=\mathcal O_K/\mathfrak p$ and $\kappa(\mathfrak P)=\mathcal O_L/\mathfrak P$, finite fields with
$\#\kappa(\mathfrak p)=N\mathfrak p$. Every $\sigma\in G$ maps $\mathcal O_L$ to itself and primes above $\mathfrak p$ to primes above $\mathfrak p$.

\begin{proposition}[decomposition and inertia groups]\label{prop:ant:galoissplitting}
(i) $G$ acts transitively on the primes of $L$ above $\mathfrak p$. Hence all $e(\mathfrak P\mid\mathfrak p)=e$ are equal, all
$f(\mathfrak P\mid\mathfrak p)=f$ are equal, and $efg=n=\#G$.
(ii) The \emph{decomposition group} $D_{\mathfrak P}=\{\sigma\in G:\sigma\mathfrak P=\mathfrak P\}$ is the stabiliser of $\mathfrak P$ (\Cref{ch:groups}); each
$\sigma\in D_{\mathfrak P}$ induces an automorphism $\bar\sigma$ of $\kappa(\mathfrak P)$ over $\kappa(\mathfrak p)$, and the \emph{inertia group}
$I_{\mathfrak P}=\{\sigma\in D_{\mathfrak P}:\sigma(x)\equiv x\pmod{\mathfrak P}\text{ for all }x\in\mathcal O_L\}$ is the kernel of $\sigma\mapsto\bar\sigma$, a
normal subgroup of $D_{\mathfrak P}$. Moreover $[G:D_{\mathfrak P}]=g$, $\#D_{\mathfrak P}=ef$, $D_{\sigma\mathfrak P}=\sigma D_{\mathfrak P}\sigma^{-1}$ and
$I_{\sigma\mathfrak P}=\sigma I_{\mathfrak P}\sigma^{-1}$; if $G$ is abelian, $D_{\mathfrak P}$ and $I_{\mathfrak P}$ depend only on $\mathfrak p$.
\end{proposition}

\begin{proof}
(i) Suppose $\mathfrak P'\mid\mathfrak p$ and $\mathfrak P'\ne\sigma\mathfrak P$ for all $\sigma\in G$. By the Chinese remainder theorem there is $x\in\mathcal O_L$
with $x\in\mathfrak P'$ and $x\equiv1\pmod{\sigma\mathfrak P}$ for all $\sigma$. The norm $N_{L/K}(x)=\prod_{\sigma\in G}\sigma(x)$ is fixed by $G$ and integral,
so lies in $\mathcal O_K$; it lies in $\mathfrak P'$ (the factor $x$ does), so in $\mathfrak P'\cap\mathcal O_K=\mathfrak p\subseteq\mathfrak P$. As $\mathfrak P$ is
prime, some $\sigma(x)\in\mathfrak P$, that is $x\in\sigma^{-1}\mathfrak P$, contradicting $x\equiv1\pmod{\sigma^{-1}\mathfrak P}$. Each $\sigma$ induces
an isomorphism $\mathcal O_L/\mathfrak P\simeq\mathcal O_L/\sigma\mathfrak P$ over $\kappa(\mathfrak p)$, and applying $\sigma$ to
$\mathfrak p\mathcal O_L=\prod\mathfrak P_i^{e_i}$ permutes the factors keeping exponents; so $e$ and $f$ are constant on the orbit, which is everything, and
\Cref{thm:ant:efg} gives $efg=n$. (ii) Orbit-stabiliser (\Cref{ch:groups}) gives $[G:D_{\mathfrak P}]=g$, so $\#D_{\mathfrak P}=n/g=ef$. A kernel is normal.
$\tau\in I_{\mathfrak P}$ iff $\tau x-x\in\mathfrak P$ for all $x$ iff $(\sigma\tau\sigma^{-1})(\sigma x)-\sigma x\in\sigma\mathfrak P$ for all $x$ iff
$\sigma\tau\sigma^{-1}\in I_{\sigma\mathfrak P}$; the same for $D$. If $G$ is abelian, conjugation is trivial, and transitivity gives the last claim.
\end{proof}

The smallest example shows all three behaviours. Let $L=\Q(i)$, $K=\Q$, $G=\{1,c\}$ with $c$ complex conjugation (\Cref{ex:ant:dkexamples}(a)). At $p=5$,
$5=(2+i)(2-i)$ and $c$ exchanges the two primes, so $g=2$ and $D_{\mathfrak P}=I_{\mathfrak P}=1$. At $p=3$, $(3)$ is prime with residue field $\F_9$, so
$g=1$, $f=2$ and $D_{\mathfrak P}=G$; and $c$ induces the nontrivial automorphism of $\F_9=\F_3(i)$, as $c(i)=-i\ne i$ modulo $3$, so $I_{\mathfrak P}=1$. At
$p=2$, $(2)=(1+i)^2$, so $e=2$, $D_{\mathfrak P}=G$, and $c$ acts trivially on $\Z[i]/(1+i)=\F_2$, so $I_{\mathfrak P}=G$. In each case $\#D_{\mathfrak P}=ef$.

\begin{theorem}[the residue field extension]\label{thm:ant:residuegalois}
In this situation:
\begin{enumerate}[label=(\roman*)]
\item The homomorphism $D_{\mathfrak P}\to\Gal(\kappa(\mathfrak P)/\kappa(\mathfrak p))$, $\sigma\mapsto\bar\sigma$, is surjective with kernel $I_{\mathfrak P}$.
Hence $\#I_{\mathfrak P}=e$, and $D_{\mathfrak P}/I_{\mathfrak P}$ is cyclic of order $f$, generated by the class of any $\sigma$ with
$\bar\sigma=(x\mapsto x^{N\mathfrak p})$.
\item For $K\subseteq M\subseteq L$ with $H=\Gal(L/M)$ and $\mathfrak q=\mathfrak P\cap M$: $D_{\mathfrak P}(L/M)=D_{\mathfrak P}\cap H$ and
$I_{\mathfrak P}(L/M)=I_{\mathfrak P}\cap H$; if $M/K$ is Galois, the images of $D_{\mathfrak P}$ and $I_{\mathfrak P}$ in $G/H=\Gal(M/K)$ are $D_{\mathfrak q}$ and
$I_{\mathfrak q}$.
\item $\mathfrak p$ splits completely in $L$ iff $D_{\mathfrak P}=1$, and $\mathfrak p$ is unramified in $L$ iff $I_{\mathfrak P}=1$.
\end{enumerate}
\end{theorem}

\begin{proof}
(i) $\kappa(\mathfrak P)/\kappa(\mathfrak p)$ is an extension of finite fields of degree $f$, Galois with cyclic group generated by $x\mapsto x^{N\mathfrak p}$
(\Cref{ch:rings}). Choose $\bar\theta$ generating the cyclic group $\kappa(\mathfrak P)^\times$, so $\kappa(\mathfrak P)=\kappa(\mathfrak p)(\bar\theta)$, and by
the Chinese remainder theorem a lift $\theta\in\mathcal O_L$ of $\bar\theta$ lying in every $\tau\mathfrak P\ne\mathfrak P$. The polynomial
$h(X)=\prod_{\sigma\in G}(X-\sigma\theta)$ has coefficients in $\mathcal O_K$, and its reduction $\bar h=\prod_\sigma(X-\overline{\sigma\theta})$ modulo $\mathfrak P$ has
coefficients in $\kappa(\mathfrak p)$ and the root $\bar\theta$; so $\bar\theta^{N\mathfrak p}$ is also a root: $\bar\theta^{N\mathfrak p}=\overline{\sigma\theta}$
for some $\sigma$. Since $\bar\theta^{N\mathfrak p}\ne0$, $\sigma\theta\notin\mathfrak P$, so $\theta\notin\sigma^{-1}\mathfrak P$, which forces
$\sigma^{-1}\mathfrak P=\mathfrak P$: $\sigma\in D_{\mathfrak P}$, and $\bar\sigma(\bar\theta)=\bar\theta^{N\mathfrak p}$, so $\bar\sigma$ is the generator
$x\mapsto x^{N\mathfrak p}$. The image is therefore everything; the kernel is $I_{\mathfrak P}$ by definition, and $\#I_{\mathfrak P}=ef/f=e$.
(ii) The first two equalities hold by definition. Restriction maps $D_{\mathfrak P}$ into $D_{\mathfrak q}$ and $I_{\mathfrak P}$ into $I_{\mathfrak q}$, with
kernels $D_{\mathfrak P}\cap H$ and $I_{\mathfrak P}\cap H$, of orders $e(\mathfrak P\mid\mathfrak q)f(\mathfrak P\mid\mathfrak q)$ and $e(\mathfrak P\mid\mathfrak q)$ by (i) for
$L/M$. So the images have orders $e(\mathfrak q\mid\mathfrak p)f(\mathfrak q\mid\mathfrak p)=\#D_{\mathfrak q}$ and $e(\mathfrak q\mid\mathfrak p)=\#I_{\mathfrak q}$
(\Cref{thm:ant:efg}, and (i) for $M/K$). (iii) $D_{\mathfrak P}=1$ iff $ef=1$ iff $g=n$; $I_{\mathfrak P}=1$ iff $e=1$.
\end{proof}

\begin{convention}[arithmetic and geometric Frobenius]\label{conv:ant:frobenius}
For $\mathfrak P$ unramified over $\mathfrak p$, the \emph{arithmetic Frobenius} $\Frob_{\mathfrak P}$ is the element of \Cref{prop:ant:frobenius}, characterised
by $\Frob_{\mathfrak P}(x)\equiv x^{N\mathfrak p}\pmod{\mathfrak P}$; its inverse is the \emph{geometric Frobenius}. The book writes $\Frob_p$ for an
arithmetic Frobenius at $p$ (\Cref{prop:ant:frobGQ}), as \Cref{ch:galois} does and as the paper does (Section 1.1: ``$\Frob_p$ is an arithmetic Frobenius
at $p$ and $\Frob_p^{-1}$ a geometric one''). So $\Frob_p$ acts on $\mu_n$ ($p\nmid n$) by $\zeta\mapsto\zeta^p$, $\chi_\ell(\Frob_p)=p$,
$\bar\chi_\ell(\Frob_p)=p\bmod\ell$, and for $E/\Q$ with good reduction at $p\ne\ell$ the characteristic polynomial of $\rho_{E,\ell}(\Frob_p)$ on $V_\ell E$ is
$X^2-a_pX+p$ (\Cref{ch:ellcurves,ch:galrep}). Local $L$-factors of a representation $V$ of $G_\Q$ on a finite-dimensional $\Q_\ell$-vector space are
written with the geometric Frobenius:
\[
L_p(V,s)=\det\bigl(1-\Frob_p^{-1}p^{-s}\bigm|V^{I_p}\bigr)^{-1}\qquad(p\ne\ell),
\]
so that $L_p(\Q_\ell(1),s)=(1-p^{-1-s})^{-1}$ and $L_p(V(1),s)=L_p(V,s+1)$, as $\chi_\ell(\Frob_p^{-1})=p^{-1}$. In detail: on $V(1)=V\otimes\chi_\ell$ the
geometric Frobenius acts as $p^{-1}\Frob_p^{-1}$, so $\det(1-p^{-1}\Frob_p^{-1}p^{-s}\mid V^{I_p})=\det(1-\Frob_p^{-1}p^{-(s+1)}\mid V^{I_p})$ (the twist by the
unramified $\chi_\ell$ does not change the inertia invariants). On $V_\ell E$ the geometric Frobenius has the
eigenvalues $\alpha^{-1},\beta^{-1}$ (with $\alpha+\beta=a_p$, $\alpha\beta=p$, the eigenvalues of $\rho_{E,\ell}(\Frob_p)$), and on the dual $V_\ell E^\vee$, where
$g$ acts by the inverse transpose of $\rho(g)$, the geometric Frobenius acts by the transpose of $\rho(\Frob_p)$ and has the eigenvalues $\alpha,\beta$; so
$\det(1-\Frob_p^{-1}p^{-s}\mid V_\ell E^\vee)=(1-\alpha p^{-s})(1-\beta p^{-s})=1-a_pp^{-s}+p^{1-2s}$, and $L_p(V_\ell E^\vee,s)$ is the Euler factor of $L(E,s)$. \emph{Dictionary.} When the multiset of Frobenius eigenvalues on
$V^{I_p}$ is stable under inversion, the two conventions give the same factor: this is the case for $\ad V$ at a good prime (eigenvalues
$\alpha/\beta,1,\beta/\alpha$) and on the lines on which Frobenius acts by $\pm1$ (types $\kod{II}$, $\kod{III}$, $\kod{IV}$ and their duals). When it is not
stable they differ: if $E$ has multiplicative reduction at $q$, $(\ad V)^{I_q}$ is a line on which the arithmetic Frobenius acts by $q$ and the geometric one by
$q^{-1}$ (\Cref{ex:ant:multiplicative}), and the factor $1-q^{-1-s}$ of the paper's table requires the geometric one.
\end{convention}

\begin{example}[the multiplicative case]\label{ex:ant:multiplicative}
Let $E$ have split multiplicative reduction at $q\ne\ell$. In a suitable basis $e_1,e_2$ of $T_\ell E$ (\Cref{ch:reduction}), inertia acts through
$\sm1{*}01$ with a nonzero corner, and an arithmetic Frobenius through $\rho(\Frob)=\sm q{c'}01$. On $\ad V$, with basis $e=e_1\otimes e_2^\vee$,
$h$, $f$ (\Cref{ch:reps}), the invariants of a nontrivial unipotent $u=\sm1t01$, $t\ne0$, form the line $\Q_\ell e$. Indeed, with $e=\sm0100$,
$h=\sm100{-1}$, $f=\sm0010$, a direct multiplication gives
\[
ueu^{-1}=e,\qquad uhu^{-1}=h-2te,\qquad ufu^{-1}=f+th-t^2e,
\]
so $xe+yh+zf$ is sent to $(x-2ty-t^2z)e+(y+tz)h+zf$, and it is fixed iff $tz=0$ and $2ty+t^2z=0$, that is $z=y=0$ (as $t\ne0$ and $2$ is invertible). The
image of inertia consists of such $u$, and contains one with $t\ne0$, so $(\ad V)^{I_q}=\Q_\ell e$. Next $\rho(\Frob)\,e\,\rho(\Frob)^{-1}=qe$: with
$F=\sm q{c'}01$ and $F^{-1}=\sm{q^{-1}}{-c'q^{-1}}01$, $Fe=\sm0q00$ and $FeF^{-1}=\sm0q00=qe$ (in general conjugation by $\sm a*0d$ scales $e$ by $a/d$). So $\Frob$ acts on
$(\ad V)^{I_q}$ by $q$, and $\Frob^{-1}$ by $q^{-1}$: $\det(1-\Frob_q^{-1}q^{-s}\mid(\ad V)^{I_q})=1-q^{-1-s}$, the entry of the paper's table, whereas the
arithmetic Frobenius would give $1-q^{1-s}$. The nonsplit case is the same, the unramified quadratic character $\pm1$ cancelling in the adjoint.
\end{example}

\begin{proposition}[the Frobenius element]\label{prop:ant:frobenius}
Let $L/K$ be Galois with group $G$, $\mathfrak p$ unramified in $L$ and $\mathfrak P\mid\mathfrak p$.
\begin{enumerate}[label=(\roman*)]
\item There is a unique $\Frob_{\mathfrak P}\in G$ with $\Frob_{\mathfrak P}(x)\equiv x^{N\mathfrak p}\pmod{\mathfrak P}$ for all $x\in\mathcal O_L$; it lies in
$D_{\mathfrak P}$, generates it, and has order $f$.
\item $\Frob_{\sigma\mathfrak P}=\sigma\Frob_{\mathfrak P}\sigma^{-1}$; so $\{\Frob_{\mathfrak P}:\mathfrak P\mid\mathfrak p\}$ is a conjugacy class of $G$, the
\emph{Frobenius class} $\Frob_{\mathfrak p}$.
\item $\mathfrak p$ splits completely in $L$ iff $\Frob_{\mathfrak p}=\{1\}$.
\item If $G$ is abelian, $\Frob_{\mathfrak P}$ depends only on $\mathfrak p$ and is written $\Frob_{\mathfrak p}\in G$.
\item If $K\subseteq M\subseteq L$ with $M/K$ Galois, then $\Frob_{\mathfrak P}|_M=\Frob_{\mathfrak P\cap M}$. If $K\subseteq K'\subseteq L$ and
$\mathfrak p'=\mathfrak P\cap K'$ has residue degree $f'$ over $\mathfrak p$, then $\Frob_{\mathfrak P}(L/K')=\Frob_{\mathfrak P}(L/K)^{f'}$.
\item If $\mathfrak p$ is ramified, the elements of $D_{\mathfrak P}$ inducing $x\mapsto x^{N\mathfrak p}$ on $\kappa(\mathfrak P)$ form a coset of
$I_{\mathfrak P}$; each of them is called \emph{a Frobenius at $\mathfrak P$}, and there is no distinguished one.
\end{enumerate}
\end{proposition}

\begin{proof}
(i) By \Cref{thm:ant:residuegalois}, $I_{\mathfrak P}=1$ and $D_{\mathfrak P}\simeq\Gal(\kappa(\mathfrak P)/\kappa(\mathfrak p))$, cyclic of order $f$; let
$\Frob_{\mathfrak P}$ be the preimage of $x\mapsto x^{N\mathfrak p}$. If $\sigma\in G$ satisfies the congruence, then $\sigma(\mathfrak P)\subseteq\mathfrak P$
(take $x\in\mathfrak P$), so $\sigma\mathfrak P=\mathfrak P$ and $\sigma\in D_{\mathfrak P}$ induces $x\mapsto x^{N\mathfrak p}$: $\sigma=\Frob_{\mathfrak P}$.
(ii) $\sigma\Frob_{\mathfrak P}\sigma^{-1}(x)=\sigma(\Frob_{\mathfrak P}(\sigma^{-1}x))\equiv\sigma((\sigma^{-1}x)^{N\mathfrak p})=x^{N\mathfrak p}\pmod{\sigma\mathfrak P}$,
and use uniqueness; all primes above $\mathfrak p$ are conjugate. (iii) $\Frob_{\mathfrak P}$ generates $D_{\mathfrak P}$; use
\Cref{thm:ant:residuegalois}(iii). (iv) follows from (ii). (v) $\Frob_{\mathfrak P}|_M$ satisfies the defining congruence modulo $\mathfrak P\cap\mathcal O_M$ on
$\mathcal O_M$, and $\mathfrak P\cap M$ is unramified over $\mathfrak p$ (\Cref{thm:ant:efg}); use uniqueness in $\Gal(M/K)$. For $K'$: both elements lie in
$D_{\mathfrak P}(L/K)$ (as $D_{\mathfrak P}(L/K')\subseteq D_{\mathfrak P}(L/K)$) and induce $x\mapsto x^{(N\mathfrak p)^{f'}}=x^{N\mathfrak p'}$ on $\kappa(\mathfrak P)$;
$D_{\mathfrak P}(L/K)\to\Gal(\kappa(\mathfrak P)/\kappa(\mathfrak p))$ is injective. (vi) is \Cref{thm:ant:residuegalois}(i).
\end{proof}

Examples, proved in \Cref{thm:ant:cyclotomicsplitting} and \Cref{prop:ant:quadchar}: in $\Q(\zeta_n)$, for $p\nmid n$, $\Frob_p=\sigma_p\colon\zeta_n\mapsto\zeta_n^p$;
in $\Q(\sqrt d)$ of discriminant $D$, for $p\nmid D$, $\Frob_p(\sqrt d)=\bigl(\frac Dp\bigr)\sqrt d$. The defining congruence can be checked by hand in
$\Q(i)$: for an odd prime $p$ and $x=a+bi\in\Z[i]$, $x^p\equiv a^p+b^pi^p\equiv a+b\,i^p\pmod p$ (binomial theorem modulo $p$ and Fermat), and $i^p=(-1)^{(p-1)/2}i$.
So $\Frob_p$ is the identity for $p\equiv1\pmod4$ and complex conjugation for $p\equiv3\pmod4$, matching the decomposition groups found after
\Cref{prop:ant:galoissplitting}.

\begin{proposition}[Frobenius and subfields]\label{prop:ant:frobsubfield}
Let $L/\Q$ be finite Galois with group $G$, $H\le G$, $M=L^H$, $p$ a prime unramified in $L$, $\mathfrak P\mid p$ and $F=\Frob_{\mathfrak P}$. The map
$\sigma H\mapsto\sigma^{-1}\mathfrak P\cap M$ induces a bijection from the orbits of $\langle F\rangle$ acting on $G/H$ by left multiplication to the primes
of $M$ above $p$; the residue degree of the prime is the length of the orbit, and all ramification indices are $1$.
\end{proposition}

\begin{proof}
Every prime of $M$ above $p$ lies under a prime of $L$ above $p$ (\Cref{thm:ant:efg} for $L/M$), and these are the $\tau\mathfrak P$, $\tau\in G$. Two of
them, $\tau\mathfrak P$ and $\tau'\mathfrak P$, lie over the same prime of $M$ iff they are conjugate under $H=\Gal(L/M)$
(\Cref{prop:ant:galoissplitting}(i) for $L/M$), iff $\tau'\in H\tau D$ with $D=D_{\mathfrak P}=\langle F\rangle$. So the primes of $M$ above $p$ correspond to the
double cosets $H\backslash G/D$, and inversion $H\tau D\mapsto D\tau^{-1}H$ identifies these with the $D$-orbits on $G/H$; with $\sigma=\tau^{-1}$ this is the
stated map. For $\mathfrak q=\tau\mathfrak P\cap M$: $D_{\tau\mathfrak P}(L/M)=\tau D\tau^{-1}\cap H$ has order $f(\tau\mathfrak P\mid\mathfrak q)$ (all ramification is
trivial), so $f(\mathfrak q\mid p)=f/\#(\tau D\tau^{-1}\cap H)=\#D/\#(D\cap\tau^{-1}H\tau)$, which is the length of the $D$-orbit of $\tau^{-1}H$ (its stabiliser in
$D$ is $D\cap\tau^{-1}H\tau$).
\end{proof}

\begin{corollary}[Dedekind's theorem on cycle types]\label{cor:ant:cycletype}
Let $f\in\Z[X]$ be separable of degree $n$ with leading coefficient $a$, $p$ a prime with $p\nmid a\disc(f)$, $L$ the splitting field of $f$, and
$d_1,\dots,d_r$ the degrees of the irreducible factors of $f\bmod p$ in $\F_p[X]$. Then $p$ is unramified in $L$, and for $\mathfrak P\mid p$ the
permutation of the roots of $f$ induced by $\Frob_{\mathfrak P}$ has cycle type $(d_1,\dots,d_r)$.
\end{corollary}

\begin{proof}
Replacing $f$ by $a^{n-1}f(X/a)$, whose roots are $a$ times those of $f$, whose discriminant is a power of $a$ times $\disc(f)$, and whose
factorisation modulo $p$ has the same degrees ($a$ is invertible modulo $p$), we may assume $f$ monic. Its roots $\beta_1,\dots,\beta_n$ lie in $\mathcal O_L$,
and $\disc(f)=\prod_{i<j}(\beta_i-\beta_j)^2\notin\mathfrak P$, so their reductions are distinct roots of $\bar f$ in $\kappa(\mathfrak P)$: reduction is a
bijection from the roots of $f$ to the roots of $\bar f$, compatible with the action of $D_{\mathfrak P}$. An element of $I_{\mathfrak P}$ acts trivially on
$\kappa(\mathfrak P)$, so fixes each reduction, hence each $\beta_i$, hence $L$: $I_{\mathfrak P}=1$ and $p$ is unramified. $\Frob_{\mathfrak P}$ acts on the
reductions as $x\mapsto x^p$, and the orbit of a root of $\bar f$ under $x\mapsto x^p$ consists of the roots of its irreducible factor over $\F_p$
(\Cref{ch:rings}); so the cycles have lengths $d_1,\dots,d_r$.
\end{proof}

For example, in $\Q(2^{1/3})$ (\Cref{ex:ant:dkexamples}(e)), $X^3-2\equiv(X+2)(X^2+3X+4)\pmod5$, so $\Frob_5$ in $\Gal(\Q(2^{1/3},\zeta_3)/\Q)\simeq S_3$ is a
transposition: it fixes one root and swaps the other two, and the primes above $5$ in $\Q(2^{1/3})$ have residue degrees $1$ and $2$, the orbit
lengths of \Cref{prop:ant:frobsubfield}. To see the dictionary concretely, take $H=\Gal(L/\Q(2^{1/3}))$, the stabiliser of the real root; then $G/H$ is
identified with the set of the three cube roots of $2$, and $\langle\Frob_{\mathfrak P}\rangle$-orbits on $G/H$ are orbits on the roots. At $p=31$, $X^3-2$ has
three roots modulo $31$, $\Frob$ is trivial, there are three orbits of length $1$, and $31$ splits completely in $\Q(2^{1/3})$; at $p=7$ it is irreducible,
$\Frob$ is a $3$-cycle, there is one orbit of length $3$, and $7$ is inert, in agreement with \Cref{ex:ant:dkexamples}(e).

\begin{proposition}[Frobenius elements in $G_\Q$]\label{prop:ant:frobGQ}
Let $p$ be a prime. A \emph{prime of $\bar\Z$ above $p$} is a maximal ideal $\bar{\mathfrak P}$ of $\bar\Z$ containing $p$.
\begin{enumerate}[label=(\roman*)]
\item $\bar{\mathfrak P}\mapsto(\bar{\mathfrak P}\cap\mathcal O_L)_L$ is a bijection from the primes of $\bar\Z$ above $p$ to the compatible families of primes
above $p$ in the finite Galois extensions $L/\Q$; $G_\Q$ acts transitively on them, and $\bar\Z/\bar{\mathfrak P}\simeq\Fbar_p$.
\item The \emph{decomposition group} $D_{\bar{\mathfrak P}}=\{\sigma\in G_\Q:\sigma\bar{\mathfrak P}=\bar{\mathfrak P}\}$ and the \emph{inertia group}
$I_{\bar{\mathfrak P}}$ (the kernel of the action of $D_{\bar{\mathfrak P}}$ on $\bar\Z/\bar{\mathfrak P}$) are closed subgroups, equal to
$\varprojlim D_{\mathfrak P_L}$ and $\varprojlim I_{\mathfrak P_L}$, with surjective projections to $D_{\mathfrak P_L}$ and $I_{\mathfrak P_L}$ for every finite
Galois $L$; for another choice $\sigma\bar{\mathfrak P}$ they are conjugated by $\sigma$. Once $\bar{\mathfrak P}$ is fixed we write $D_p$, $I_p$. The
sequence $1\to I_p\to D_p\to\Gal(\Fbar_p/\F_p)\to1$ is exact, and $\Gal(\Fbar_p/\F_p)\simeq\hat\Z$ is topologically generated by $x\mapsto x^p$.
\item \emph{A Frobenius at $p$} is any $\Frob_p\in D_p$ mapping to $x\mapsto x^p$; it is well defined up to $I_p$ and, when $\bar{\mathfrak P}$ varies, up
to conjugation in $G_\Q$. Its image in $\Gal(L/\Q)$, for $L$ finite Galois and unramified at $p$, lies in the class $\Frob_p(L/\Q)$.
\item A continuous representation $\rho$ of $G_\Q$ is \emph{unramified at $p$} if $\rho(I_p)=1$ for one, equivalently every, $\bar{\mathfrak P}\mid p$; then
the conjugacy class of $\rho(\Frob_p)$, and in particular $\det(X-\rho(\Frob_p))$, is well defined. If $\rho$ has finite image, with $L$ the field cut
out by $\rho$ (the fixed field of $\ker\rho$), then $\rho$ is unramified at $p$ iff $p$ is unramified in $L$, and $\rho(\Frob_p)$ is the image of the
class $\Frob_p(L/\Q)$.
\item (Stated; \Cref{ch:local}.) An embedding $\Qbar\to\Qbar_p$ identifies $D_{\bar{\mathfrak P}}$, for the prime it defines, with $G_{\Q_p}$ and
$I_{\bar{\mathfrak P}}$ with the inertia subgroup of $G_{\Q_p}$.
\end{enumerate}
\end{proposition}

\begin{proof}
(i) $\bar{\mathfrak P}\cap\mathcal O_L$ is a prime of $\mathcal O_L$ containing $p$, and the family is compatible. Conversely the union of a compatible family is
an ideal of $\bar\Z=\bigcup_L\mathcal O_L$ (a directed union) which is prime and contains $p$, and every nonzero prime of $\bar\Z$ is maximal (its trace on
each $\mathcal O_L$ is). Given $\bar{\mathfrak P},\bar{\mathfrak P}'$, the sets $S_L=\{\sigma\in\Gal(L/\Q):\sigma\mathfrak P_L=\mathfrak P'_L\}$ are nonempty
(\Cref{prop:ant:galoissplitting}) and finite, and restriction maps $S_{L'}$ into $S_L$ for $L\subseteq L'$; an inverse limit of nonempty finite sets is
nonempty (\Cref{ch:galois}), and an element of it is a $\sigma\in G_\Q$ with $\sigma\bar{\mathfrak P}=\bar{\mathfrak P}'$. The residue ring $\bigcup_L\kappa(\mathfrak P_L)$
is an algebraic extension of $\F_p$; it contains every $\F_{p^f}$, since a root in $\bar\Z$ of a monic lift of an irreducible polynomial of degree $f$
over $\F_p$ reduces to a root of that polynomial. (ii) By definition $D_{\bar{\mathfrak P}}$ is the set of $\sigma$ whose restrictions lie in all
$D_{\mathfrak P_L}$, which is the inverse limit, a closed subgroup; likewise for $I$. Surjectivity of the projections: for $L\subseteq L'$ the restriction
maps $D_{\mathfrak P_{L'}}\to D_{\mathfrak P_L}$ and $I_{\mathfrak P_{L'}}\to I_{\mathfrak P_L}$ are surjective (\Cref{thm:ant:residuegalois}(ii)), and the
fibres over a given element form an inverse system of nonempty finite sets. The exact sequence is the inverse limit of those of
\Cref{thm:ant:residuegalois}(i), surjectivity again by compactness; and $\Gal(\Fbar_p/\F_p)=\varprojlim\Z/f\simeq\hat\Z$ (\Cref{ch:galois}).
(iii) follows from (ii) and \Cref{prop:ant:frobenius}. (iv) Inertia groups for different $\bar{\mathfrak P}$ are conjugate. If $\rho(I_p)=1$, two Frobenius
elements differ by $I_p$ and conjugation, so $\rho(\Frob_p)$ is defined up to conjugacy. If $\rho$ factors through $\Gal(L/\Q)$ injectively, then
$\rho(I_p)$ is the image of $I_{\mathfrak P_L}$ by (ii), trivial iff $p$ is unramified in $L$, and the image of $\Frob_p$ is $\Frob_{\mathfrak P_L}$.
(v) We do not prove this here; see \Cref{ch:local}.
\end{proof}

\begin{proposition}[tame inertia]\label{prop:ant:tameconj}
Let $L/K$ be finite Galois, $\mathfrak P\mid\mathfrak p$, $p$ the residue characteristic, and suppose that $\#I_{\mathfrak P}$ is prime to $p$ ($\mathfrak P$ is
\emph{tamely ramified}). Then $I_{\mathfrak P}$ is cyclic, and for every $F\in D_{\mathfrak P}$ inducing $x\mapsto x^{N\mathfrak p}$ on $\kappa(\mathfrak P)$ and every
$\tau\in I_{\mathfrak P}$,
\[
F\tau F^{-1}=\tau^{N\mathfrak p}.
\]
\end{proposition}

\begin{proof}[Sketch of proof]
Let $\varpi$ be a uniformiser of the discrete valuation ring $\mathcal O_{L,\mathfrak P}$ (\Cref{prop:ant:crt}(iv)), on which $D_{\mathfrak P}$ acts. For
$\tau\in I_{\mathfrak P}$, $\tau(\varpi)$ is again a uniformiser, so $\tau(\varpi)/\varpi$ is a unit; put $\theta(\tau)=\tau(\varpi)/\varpi\bmod\mathfrak P
\in\kappa(\mathfrak P)^\times$. Since $\tau$ acts trivially modulo $\mathfrak P$ on $\mathcal O_{L,\mathfrak P}$ (the congruence $\tau x\equiv x$ extends from
$\mathcal O_L$ to fractions with denominators prime to $\mathfrak P$), $\theta$ does not depend on $\varpi$ (replacing $\varpi$ by $u\varpi$ multiplies by
$\tau(u)/u\equiv1$) and is a homomorphism: for $\sigma,\tau\in I_{\mathfrak P}$,
$\sigma\tau(\varpi)/\varpi=\sigma\bigl(\tau(\varpi)/\varpi\bigr)\cdot\sigma(\varpi)/\varpi$, and $\sigma$ fixes the unit $\tau(\varpi)/\varpi$ modulo $\mathfrak P$, so
$\theta(\sigma\tau)\equiv\theta(\tau)\theta(\sigma)$. The kernel
of $\theta$ is the first ramification group, which is a $p$-group; this is the step we do not prove (it uses the filtration of $I_{\mathfrak P}$ by the
higher ramification groups, whose successive quotients embed in the additive groups $\mathfrak P^i/\mathfrak P^{i+1}$; \cite[Chapter IV, Section 2]{SerreLocalFields}).
So when $\#I_{\mathfrak P}$ is prime to $p$, $\theta$ is injective and $I_{\mathfrak P}$, a subgroup of the cyclic group $\kappa(\mathfrak P)^\times$, is cyclic.
For $F$ as stated, $\varpi'=F^{-1}(\varpi)$ is a uniformiser ($F\in D_{\mathfrak P}$ preserves $\mathfrak P$ and its powers), and
$F\tau F^{-1}(\varpi)/\varpi=F(\tau(\varpi'))/F(\varpi')=F\bigl(\tau(\varpi')/\varpi'\bigr)$. Reducing modulo $\mathfrak P$: the unit $\tau(\varpi')/\varpi'$ reduces to
$\theta(\tau)$ (computed with the uniformiser $\varpi'$, which is allowed), and $F$ acts on $\kappa(\mathfrak P)$ as $x\mapsto x^{N\mathfrak p}$, so
$\theta(F\tau F^{-1})=\theta(\tau)^{N\mathfrak p}=\theta(\tau^{N\mathfrak p})$. Injectivity of $\theta$ gives $F\tau F^{-1}=\tau^{N\mathfrak p}$. The profinite form, $\Frob\,t\,\Frob^{-1}=t^q$ on the tame inertia
quotient of $G_{\Q_q}$, is proved in \Cref{ch:local}.
\end{proof}

\begin{corollary}[Frobenius on a tame image of inertia]\label{cor:ant:frobsigma}
Let $q$ be a prime, $\ell\ne q$ an odd prime, and $\rho\colon G_\Q\to\GL_2(\Z_\ell)$ continuous with $\rho(I_q)$ finite of order prime to $q$ (for
$\rho=\rho_{E,\ell}$ at a prime $q\ge5$ of additive potentially good reduction the order is $2$, $3$, $4$ or $6$, by Kraus's theorem, \Cref{ch:reduction,ch:galrep}).
Then $\rho(I_q)$ is cyclic, and for every Frobenius $F\in D_q$ (\Cref{prop:ant:frobGQ}) and every $\tau\in I_q$,
\[
\rho(F)\rho(\tau)\rho(F)^{-1}=\rho(\tau)^q.
\]
In particular, if $\rho(I_q)$ is generated by $\sigma$ of order $3$ (or by $-\sigma$, of order $6$), then $\rho(F)$ commutes with $\sigma$ if $q\equiv1\pmod3$
and conjugates it to $\sigma^{-1}$ if $q\equiv2\pmod3$; if $\rho(I_q)$ is generated by $\tau$ of order $4$, $\rho(F)$ commutes with $\tau$ if $q\equiv1\pmod4$ and
conjugates it to $\tau^{-1}$ if $q\equiv3\pmod4$. None of this depends on the choice of $F$.
\end{corollary}

\begin{proof}
The kernel $1+\ell\Mat_2(\Z_\ell)$ of reduction modulo $\ell$ contains no nontrivial element of finite order when $\ell$ is odd (\Cref{ch:groups,ch:reps}).
For the reader's convenience, the argument: write $g\ne1$ as $g=1+\ell^kx$ with $k\ge1$ and $x\not\equiv0\pmod\ell$. For $m$ prime to $\ell$,
$g^m\equiv1+m\ell^kx\not\equiv1\pmod{\ell^{k+1}}$. For $m=\ell$, $g^\ell=1+\ell^{k+1}x+\binom\ell2\ell^{2k}x^2+\sum_{j\ge3}\binom\ell j\ell^{jk}x^j$, where
$\ell\mid\binom\ell2$ ($\ell$ odd) and $jk\ge k+2$ for $j\ge3$; so $g^\ell=1+\ell^{k+1}x'$ with $x'\equiv x\not\equiv0\pmod\ell$, again of the same shape. By induction
$g^{\ell^a}=1+\ell^{k+a}x_a$ with $x_a\not\equiv0$, and then $g^{\ell^am'}\ne1$ for $m'$ prime to $\ell$; so $g$ has infinite order. (For $\ell=2$ the step
$m=\ell$ fails at $k=1$: $-1\in1+2\Mat_2(\Z_2)$.) Hence
reduction $\rho\mapsto\bar\rho$ is injective on the finite group $\rho(I_q)$. Let $L$ be the field cut out by $\bar\rho$, a finite Galois extension of $\Q$
with $\bar\rho\colon\Gal(L/\Q)\simeq\bar\rho(G_\Q)$, and $\mathfrak P=\bar{\mathfrak P}\cap\mathcal O_L$. By \Cref{prop:ant:frobGQ}(ii), $I_{\mathfrak P}$ is the
image of $I_q$, so $\bar\rho(I_{\mathfrak P})=\bar\rho(I_q)\simeq\rho(I_q)$ has order prime to $q$, and $F|_L\in D_{\mathfrak P}$ induces $x\mapsto x^q$ on
$\kappa(\mathfrak P)$. \Cref{prop:ant:tameconj} gives that $I_{\mathfrak P}$ is cyclic and $F|_L\,\tau|_L\,F|_L^{-1}=(\tau|_L)^q$, that is,
$\bar\rho(F\tau F^{-1})=\bar\rho(\tau^q)$. Both $F\tau F^{-1}$ and $\tau^q$ lie in $I_q$, which is normal in $D_q$; so $\rho(F\tau F^{-1})$ and
$\rho(\tau)^q$ lie in $\rho(I_q)$ and have the same reduction, hence are equal. The special cases: $\sigma^q=\sigma$ or $\sigma^2=\sigma^{-1}$ according to
$q\bmod3$ ($q\ne3$ since $3\mid\#\rho(I_q)$ is prime to $q$); $(-\sigma)^q=-\sigma^q$ for $q$ odd, and for $q=2$ the order $6$ is not prime to $q$; $\tau^q=\tau$
or $\tau^3=\tau^{-1}$ according to $q\bmod4$. Two Frobenius elements differ by an element of $I_q$, whose image commutes with the cyclic group $\rho(I_q)$.
\end{proof}

\begin{paperlink}\label{pl:ant:frobsigma}
Lemma lem:IV: ``The arithmetic Frobenius acts on the tame quotient of $I_q$ by $x\mapsto x^q$, so $\rho(\Frob_q)\sigma\rho(\Frob_q)^{-1}=\sigma^q=
\sigma^{\chi_{-3}(q)}$'' is \Cref{cor:ant:frobsigma}, at every prime $q\ne3$ at which the image of inertia has order $3$, including $q=2$ (the order $3$ is
prime to $2$); $\sigma^{\chi_{-3}(q)}=\sigma^q$ because $\chi_{-3}(q)=1$ iff $q\equiv1\pmod3$ (\Cref{prop:ant:quadchar}). Lemma lem:euler: ``Frobenius
centralises $\sigma$ if $q\equiv1\pmod3$ and inverts it if $q\equiv2$, for either choice of Frobenius'' is the same corollary, and the types
$\kod{III}$, $\kod{III}^*$ use its order-$4$ case. The action on the invariant line of the adjoint is \Cref{prop:ant:commutant}.
\end{paperlink}

\section{Quadratic fields, \texorpdfstring{$\Z[i]$}{Z[i]}, \texorpdfstring{$\Z[\zeta_3]$}{Z[zeta3]} and the rings \texorpdfstring{$\Z_\ell[\zeta_3]$}{Zl[zeta3]}}\label{sec:ant:quadratic}

Quadratic fields carry all the characters of the paper: $\chi_{-3}$ and $\chi_{-4}$ in the table of Section 1 and in Lemma lem:euler,
$\chi_{q^*}$ in the twists of Section 1.4 and Lemma lem:twist, $\chi_S$ in Theorems thm:twist and thm:quant, and $\Q(\sqrt{\ell^*})$ in the hypothesis
of Theorem thm:main. This section proves the splitting law (\Cref{thm:ant:quadsplit}), identifies the character of a quadratic field with a
Kronecker symbol (\Cref{prop:ant:quadchar}), and then studies the two rings $\Z[i]$, $\Z[\zeta_3]$ and their $\ell$-adic versions, which carry the
algebra of Lemma lem:IV ($\Z_3[\zeta_3]$ a discrete valuation ring with uniformiser $\pi=1-\zeta_3$, $T_3E$ free of rank one over it,
$E[3]\simeq\Z_3[\zeta_3]/(\pi^2)$) and of Lemma lem:euler (the line spanned by $\sqrt{-3}=2\sigma+1$ and the sign $\chi_{-3}(q)$). The Kronecker symbol
$\bigl(\frac Dp\bigr)$ is that of \Cref{ch:rings}: the Legendre symbol for odd $p$, and for $p=2$ it is $0$ if $2\mid D$, $1$ if $D\equiv\pm1\pmod8$ and $-1$ if
$D\equiv\pm3\pmod8$.

\begin{theorem}[splitting in quadratic fields]\label{thm:ant:quadsplit}
Let $K=\Q(\sqrt d)$, $d\ne0,1$ squarefree, of discriminant $D$, and $p$ a prime. Then $p$ ramifies in $K$ iff $\bigl(\frac Dp\bigr)=0$ (iff $p\mid D$), splits
iff $\bigl(\frac Dp\bigr)=1$, and is inert iff $\bigl(\frac Dp\bigr)=-1$. Explicitly: for odd $p\mid D$, $p\mathcal O_K=(p,\sqrt d)^2$; for odd $p\nmid D$ with
$a^2\equiv d\pmod p$, $p\mathcal O_K=(p,\sqrt d-a)(p,\sqrt d+a)$ with distinct factors; for $p=2$: if $d\equiv3\pmod4$, $(2)=(2,1+\sqrt d)^2$; if
$d\equiv2\pmod4$, $(2)=(2,\sqrt d)^2$; if $d\equiv1\pmod8$, $(2)=(2,\frac{1+\sqrt d}2)(2,\frac{1-\sqrt d}2)$; if $d\equiv5\pmod8$, $2$ is inert.
\end{theorem}

\begin{proof}
Apply \Cref{thm:ant:dedekindkummer}. If $p$ is odd, take $\theta=\sqrt d$, $g=X^2-d$: the index $[\mathcal O_K:\Z[\sqrt d]]$ is $1$ or $2$, prime to $p$. If
$p\mid d$, $\bar g=X^2$ and $p=(p,\sqrt d)^2$; if $d$ is a nonzero square $a^2$ modulo $p$, $\bar g=(X-a)(X+a)$ with $a\not\equiv-a$; otherwise $\bar g$ is
irreducible. For odd $p$, $p\mid D$ iff $p\mid d$, and $\bigl(\frac Dp\bigr)=\bigl(\frac dp\bigr)$ since $D/d\in\{1,4\}$. If $p=2$ and $d\equiv2,3\pmod4$, take
$\theta=\sqrt d$ (index $1$): $X^2-d\equiv X^2$ or $(X+1)^2\pmod 2$, and $2\mid D=4d$. If $d\equiv1\pmod4$ take $\theta_d=\frac{1+\sqrt d}2$, $g=X^2-X+\frac{1-d}4$
(index $1$): modulo $2$, $g\equiv X(X+1)$ if $\frac{1-d}4$ is even, that is $d\equiv1\pmod8$, and $g\equiv X^2+X+1$, irreducible, if $d\equiv5\pmod8$; here
$D=d$ is odd and $\bigl(\frac D2\bigr)=1$, resp. $-1$. The primes above $2$ for $d\equiv1\pmod8$ are $(2,\theta_d)$ and $(2,\theta_d-1)=(2,\frac{1-\sqrt d}2)$, as $\theta_d-1=-\frac{1-\sqrt d}2$.
\end{proof}

\begin{proposition}[the character of a quadratic field]\label{prop:ant:quadchar}
Let $K=\Q(\sqrt d)$ have discriminant $D$ and let $\chi_D\colon G_\Q\to\{\pm1\}$ be its character, $\sigma(\sqrt d)=\chi_D(\sigma)\sqrt d$ (\Cref{ch:galois}).
\begin{enumerate}[label=(\roman*)]
\item For every prime $p\nmid D$, $\chi_D(\Frob_p)=\bigl(\frac Dp\bigr)$, the value at $p$ of the primitive real Dirichlet character $n\mapsto\bigl(\frac Dn\bigr)$ of
conductor $\abs D$ (\Cref{ch:rings}).
\item $\chi_D$ is unramified exactly at the primes $p\nmid D$; for $p\mid D$, $\chi_D(I_p)=\{\pm1\}$.
\item In particular: $\chi_{-4}(\Frob_p)=(-1)^{(p-1)/2}$ for odd $p$ (conductor $4$); $\chi_{-3}(\Frob_p)=\bigl(\frac p3\bigr)$, which is $1$ iff $p\equiv1\pmod3$, for
$p\ne3$, including $\chi_{-3}(\Frob_2)=-1$ (conductor $3$); for an odd prime $q$, $\chi_{q^*}(\Frob_p)=\bigl(\frac{q^*}p\bigr)=\bigl(\frac pq\bigr)$ for every prime
$p\ne q$, including $p=2$ (conductor $q$); $\chi_8$ and $\chi_{-8}$, the characters of $\Q(\sqrt2)$ and $\Q(\sqrt{-2})$, have conductor $8$.
\end{enumerate}
\end{proposition}

\begin{proof}
(i) $p$ is unramified (\Cref{thm:ant:discram}, $d_K=D$), and $\Frob_p\in\Gal(K/\Q)$ is trivial iff $D_{\mathfrak P}=1$ iff $p$ splits
(\Cref{prop:ant:frobenius}(iii)) iff $\bigl(\frac Dp\bigr)=1$ (\Cref{thm:ant:quadsplit}); since both sides are $\pm1$ they agree. That $n\mapsto\bigl(\frac Dn\bigr)$ is a
primitive Dirichlet character of conductor $\abs D$ for a fundamental discriminant $D$ is proved in \Cref{ch:rings}. (ii) For $p\mid D$, $p$ ramifies, the inertia
group in $\Gal(K/\Q)$ has order $e=2$, and $\chi_D(I_p)$ is its image (\Cref{prop:ant:frobGQ}(ii)). (iii) is (i) for $D=-4,-3,q^*,\pm8$, together with
$\bigl(\frac{-4}p\bigr)=\bigl(\frac{-1}p\bigr)=(-1)^{(p-1)/2}$, $\bigl(\frac{-3}p\bigr)=\bigl(\frac p3\bigr)$ and $\bigl(\frac{q^*}p\bigr)=\bigl(\frac pq\bigr)$ (quadratic reciprocity with its
supplements, \Cref{ch:rings}, or \Cref{cor:ant:qr}); at $p=2$, $\bigl(\frac{q^*}2\bigr)=1$ iff $q^*\equiv1\pmod8$ iff $q\equiv\pm1\pmod8$ iff $\bigl(\frac2q\bigr)=1$. All
these congruence conditions were also checked numerically for $p<3000$.
\end{proof}

\begin{proposition}[quadratic fields unramified outside a set of primes]\label{prop:ant:ramonlyat}
(i) For an odd prime $q$, $\Q(\sqrt{q^*})$, $q^*=(-1)^{(q-1)/2}q$, is the unique quadratic field ramified only at $q$; its discriminant is $q^*$.
(ii) Every fundamental discriminant is uniquely a product $D=D_0\prod_{q\in T}q^*$ with $T$ a finite set of odd primes and $D_0\in\{1,-4,8,-8\}$ (the
\emph{prime discriminants}). A quadratic field is unramified outside a finite set $S$ of primes iff its discriminant has this form with $T\subseteq S$ and
$D_0=1$ unless $2\in S$.
(iii) For a finite set $S$ of odd primes, $\chi_S:=\prod_{q\in S}\chi_{q^*}$ is the character of $\Q(\sqrt{\prod_{q\in S}q^*})$, of discriminant
$\prod_{q\in S}q^*$ and conductor $\prod_{q\in S}q$; it is ramified exactly at the primes of $S$.
(iv) Hence twisting a representation of $G_\Q$ by $\chi_{q^*}$ or by $\chi_S$ does not change it on $I_p$ for $p\notin S$.
\end{proposition}

\begin{proof}
(i) $q^*\equiv1\pmod4$ and is squarefree, so it is a fundamental discriminant, and $\Q(\sqrt{q^*})$ is ramified exactly at $q$ (\Cref{thm:ant:discram}).
Conversely, a quadratic field ramified only at $q$ has odd discriminant $D$ divisible only by $q$, so $D=\pm q$; an odd fundamental discriminant is
$\equiv1\pmod4$, which fixes the sign: $D=q^*$. (ii) If $D$ is odd, it is squarefree and $\equiv1\pmod 4$; writing $D=\pm\prod_{q\mid D}q$, the product
$\prod_{q\mid D}q^*$ has the same absolute value and is $\equiv1\pmod4$, so equals $D$. If $D=4m$ with $m\equiv3\pmod4$, then $-m\equiv1\pmod 4$ is an odd
fundamental discriminant (or $1$) and $D=-4\cdot(-m)$; if $m\equiv2\pmod4$, $m=2m'$ with $m'$ odd, and $D=8m'=\pm8\cdot(\pm m')$ with the sign chosen so that
$\pm m'\equiv1\pmod 4$. Uniqueness: the odd part determines $T$, the sign and the $2$-part then determine $D_0$. The second statement follows from
\Cref{thm:ant:discram}. (iii) $\prod_{q\in S}\chi_{q^*}(\sigma)$ is the sign by which $\sigma$ multiplies $\prod_{q\in S}\sqrt{q^*}$, a square root of $\prod q^*$;
the product is a fundamental discriminant by (ii), whose conductor is its absolute value. (iv) $\chi_S$ is trivial on $I_p$ for $p\notin S$ by (iii) and
\Cref{prop:ant:quadchar}(ii).
\end{proof}

\begin{example}[splitting tables]\label{ex:ant:quadexamples}
By \Cref{thm:ant:quadsplit} and quadratic reciprocity:
\begin{center}\small
\begin{tabular}{llll}
\toprule
field & $D$ & ramified & $p$ splits iff\\
\midrule
$\Q(i)$ & $-4$ & $2$ & $p\equiv1\pmod4$\\
$\Q(\sqrt{-3})=\Q(\zeta_3)$ & $-3$ & $3$ & $p\equiv1\pmod3$ ($2$ is inert)\\
$\Q(\sqrt{-7})$ & $-7$ & $7$ & $p\equiv1,2,4\pmod7$ ($2$ splits, $-7\equiv1\pmod8$)\\
$\Q(\sqrt{-11})$ & $-11$ & $11$ & $p$ a square modulo $11$ ($2$, $7$ inert; $3$, $5$ split)\\
$\Q(\sqrt5)$ & $5$ & $5$ & $p\equiv\pm1\pmod5$\\
$\Q(\sqrt{-5})$ & $-20$ & $2,5$ & $p\equiv1,3,7,9\pmod{20}$ ($3$ and $7$ split)\\
\bottomrule
\end{tabular}
\end{center}
The paper's twists by these fields: 121a1 and its twist by $\chi_{-11}$ ($\Q(\sqrt{-11})=\Q(\sqrt{11^*})$), 539c1 and 980i1 twisted by $\chi_{-7}$
($\Q(\sqrt{-7})=\Q(\sqrt{7^*})$), in the Example after the proof of Theorem thm:qual (Example ex:qual).
\end{example}

\begin{proposition}[Gaussian and Eisenstein primes]\label{prop:ant:gausseisenstein}
(i) $\Z[i]$ and $\Z[\zeta_3]$ are Euclidean for the norms $a^2+b^2$ and $a^2-ab+b^2$ (\Cref{ch:rings}), hence principal ideal domains; their unit groups are
$\mu_4$ and $\mu_6$ (\Cref{thm:ant:units}).
(ii) Up to units, the primes of $\Z[i]$ are: $1+i$, with $2=-i(1+i)^2$; for each $p\equiv1\pmod4$, the two non-associate factors $a\pm bi$ of $p=a^2+b^2$;
for each $p\equiv3\pmod4$, $p$ itself. For example $5=(2+i)(2-i)$, $13=(3+2i)(3-2i)$.
(iii) Up to units, the primes of $\Z[\zeta_3]$ are: $\pi=1-\zeta_3$ (\Cref{prop:ant:above3}); for each $p\equiv1\pmod3$, the two non-associate factors
$a+b\zeta_3$, $a+b\zeta_3^2$ of $p=a^2-ab+b^2$; for each $p\equiv2\pmod3$, $p$ itself. For example $7=(3+\zeta_3)(3+\zeta_3^2)$, $13=(4+\zeta_3)(4+\zeta_3^2)$.
(iv) A prime $p$ is a sum of two squares iff $p=2$ or $p\equiv1\pmod4$ (Fermat); $p=a^2-ab+b^2$ for integers $a,b$ iff $p=3$ or $p\equiv1\pmod3$, and then
$p=x^2+3y^2$ for some integers $x,y$.
\end{proposition}

\begin{proof}
(ii) Every prime $\varpi$ of $\Z[i]$ lies above a unique rational $p$ (\Cref{thm:ant:dedekind}), and the primes above $p$ are given by
\Cref{ex:ant:quadexamples}: $p=2$ is ramified, $p\equiv1\pmod4$ split, $p\equiv3\pmod4$ inert. In the split case $p\Z[i]=\mathfrak P\bar{\mathfrak P}$ with
$\mathfrak P=(a+bi)$ principal of norm $p$, so $a^2+b^2=p$, and $\bar{\mathfrak P}=(a-bi)\ne\mathfrak P$. In the inert case $(p)$ is prime. (iii) The same with
\Cref{ex:ant:dkexamples}(b) and $N(a+b\zeta_3)=a^2-ab+b^2$, $\overline{a+b\zeta_3}=a+b\zeta_3^2$. (iv) $p$ is a norm from $\Z[i]$ (resp. $\Z[\zeta_3]$) iff some
prime above $p$ has norm $p$, iff $p$ is ramified or split. If $p=a^2-ab+b^2$, then $p$ is odd (it is $3$ or $\equiv1\pmod3$), so not both $a,b$ are even, and one of $a$, $b$, $a-b$ is
even; by the symmetries $(a,b)\mapsto(b,a)$ and $(a,b)\mapsto(a-b,-b)$, which preserve the form, we may assume $b=2y$ even, and then
$p=(a-y)^2+3y^2$. The small cases were also checked by computer for $p<300$.
\end{proof}

\begin{proposition}[the prime above $3$]\label{prop:ant:above3}
In $\Z[\zeta_3]$ put $\pi=1-\zeta_3$ and $\sqrt{-3}:=1+2\zeta_3$.
\begin{enumerate}[label=(\roman*)]
\item $(\sqrt{-3})^2=-3$, $\sqrt{-3}=\zeta_3-\zeta_3^2=\zeta_3\pi$, $\pi^2=-3\zeta_3$ and $3=-\zeta_3^2\pi^2$. So $\pi$ and $\sqrt{-3}$ are associates, $N(\pi)=3$,
$3\Z[\zeta_3]=(\pi)^2$ ($3$ is totally ramified), and $\pi^2=-3$ up to the unit $\zeta_3$.
\item $\Z[\zeta_3]/(\pi)=\F_3$, and $\zeta_3\equiv1\pmod\pi$.
\item $\Z[\zeta_3]/(3)=\Z[\zeta_3]/(\pi^2)$ is a local ring with nine elements, with $\F_3$-basis $1,\pi$; multiplication by $\zeta_3=1-\pi$ has the matrix
$\sm10{-1}1$ in this basis. It is a nontrivial unipotent map, and its fixed points form the line $\pi\Z[\zeta_3]/(\pi^2)$.
\item The analogue at $2$ in $\Z[i]$: $(2)=(1+i)^2$, $\Z[i]/(2)\simeq\F_2[x]/(x^2)$ with $x=1+i$, and multiplication by $i=1+x$ is a nontrivial unipotent map.
\end{enumerate}
\end{proposition}

\begin{proof}
All identities follow from $\zeta_3^2=-1-\zeta_3$: $(1+2\zeta_3)^2=1+4\zeta_3+4\zeta_3^2=-3$; $\zeta_3-\zeta_3^2=1+2\zeta_3$; $\zeta_3\pi=\zeta_3-\zeta_3^2$;
$\pi^2=1-2\zeta_3+\zeta_3^2=-3\zeta_3$; $-\zeta_3^2\pi^2=3\zeta_3^3=3$. $N(\pi)=1+1+1=3$. (ii) $\Z[\zeta_3]/(\pi)$ has $N(\pi)=3$ elements, and
$\zeta_3=1-\pi$. (iii) $(3)=(\pi^2)$ by (i); the quotient has $9$ elements, its ideals are $0$, $(\pi)/(\pi^2)$ and the whole ring
(\Cref{prop:ant:crt}(i)), and $1,\pi$ is a basis since $\pi\notin(\pi^2)$. $\zeta_3\cdot1=1-\pi$ and $\zeta_3\pi=\pi-\pi^2\equiv\pi$, giving the matrix (columns are
images). The fixed points of $x\mapsto\zeta_3x$ are the $x$ with $\pi x\equiv0$, that is $x\in(\pi)$. (iv) $(1+i)^2=2i$, so $(2)=(1+i)^2$; $\Z[i]/(2)$ has basis
$1,x$ with $x^2=2i\equiv0$, and $i\cdot1=1+x$, $i\cdot x=x+x^2=x$.
\end{proof}

\begin{lemma}[quadratic algebras over $\Z_\ell$]\label{lem:ant:quadraticZl}
Let $\ell$ be a prime, $g\in\Z_\ell[X]$ monic of degree $2$, $R=\Z_\ell[X]/(g)$ (free over $\Z_\ell$ with basis $1,x$, $x$ the class of $X$) and $\bar g=g\bmod\ell$.
\begin{enumerate}[label=(\roman*)]
\item If $\bar g$ is irreducible over $\F_\ell$, $R$ is a discrete valuation ring with uniformiser $\ell$ and residue field $\F_{\ell^2}$.
\item If $g(X+c)$ is an Eisenstein polynomial at $\ell$ for some $c\in\Z_\ell$, $R$ is a discrete valuation ring with uniformiser $x-c$ and residue field
$\F_\ell$, and $\ell R=(x-c)^2R$.
\item If $\bar g=(X-\bar a)(X-\bar b)$ with $\bar a\ne\bar b$, then $g$ has roots $a,b\in\Z_\ell$ lifting $\bar a,\bar b$, and $R\simeq\Z_\ell\times\Z_\ell$, $x\mapsto(a,b)$.
\end{enumerate}
In each case $R$ is a principal ideal ring. If $T$ is an $R$-module which is free of rank $2$ over $\Z_\ell$ and satisfies $T\otimes\Q_\ell\simeq R\otimes\Q_\ell$ as
$R\otimes\Q_\ell$-modules (automatic in cases (i) and (ii)), then $T$ is free of rank one over $R$.
\end{lemma}

\begin{proof}
$R$ is Noetherian (finite over $\Z_\ell$) and $\ell$-adically complete, so an element $1-\ell z$ is a unit ($\sum_k(\ell z)^k$ converges).
(i) $R/\ell R=\F_\ell[X]/(\bar g)=\F_{\ell^2}$ is a field, so $\ell R$ is maximal. If $y\notin\ell R$, $\bar y$ is invertible in $R/\ell R$: $yy'=1-\ell z$, and
$y$ is a unit. So $R$ is local with maximal ideal $\ell R$. Every nonzero $y$ is $\ell^ky_0$ with $y_0\notin\ell R$ (as $\bigcap_k\ell^kR=0$, $R$ being free over
$\Z_\ell$), so a product of nonzero elements is $\ell^{k+k'}$ times a unit, nonzero: $R$ is a domain. A Noetherian local domain, not a field, with
principal maximal ideal is a discrete valuation ring (\Cref{ch:modules}).
(ii) Write $y=x-c$, so $R=\Z_\ell[Y]/(Y^2+uY+v)$ with $u=\ell u'$, $u'\in\Z_\ell$, and $v=\ell w_0$, $w_0\in\Z_\ell^\times$. Then $R/(y)=\Z_\ell/(v)=\F_\ell$,
so $(y)$ is maximal. If $z=s+ty$ ($s,t\in\Z_\ell$) is not in $(y)$, then $\ell\nmid s$, and the determinant of multiplication by $z$ in the basis $1,y$ is
$s^2-stu+t^2v\equiv s^2\not\equiv0\pmod\ell$; so $z$ is a unit (Cayley-Hamilton, as in \Cref{prop:ant:tracenorm}(ii)), and $R$ is local with maximal ideal $(y)$.
Next $y^2=-\ell(u'y+w_0)$ with $u'y+w_0\notin(y)$ a unit, so $\ell R=y^2R$. Hence $\bigcap_ky^kR=\bigcap_k\ell^kR=0$, every nonzero $z$ is $y^kw$ with $w$ a unit
($k$ the largest exponent with $z\in y^kR$), and $y^k\ne0$ for all $k$; so a product of nonzero elements is nonzero, $R$ is a domain, and it is a discrete
valuation ring with uniformiser $y$ (\Cref{ch:modules}).
(iii) $R/\ell R\simeq\F_\ell[X]/(X-\bar a)\times\F_\ell[X]/(X-\bar b)=\F_\ell\times\F_\ell$ by the Chinese remainder theorem. Its idempotent $(1,0)$ lifts to an
idempotent $e\in R$ (\Cref{ch:modules}, idempotent lifting over the complete ring $\Z_\ell$). Then $R=Re\times R(1-e)$, each factor is a direct summand of
$R\simeq\Z_\ell^2$, hence free over $\Z_\ell$, and nonzero modulo $\ell$, hence of rank $1$; as $Re/\ell Re\simeq\F_\ell$ is spanned by $\bar e$, Nakayama gives
$Re=\Z_\ell e$, and similarly $R(1-e)=\Z_\ell(1-e)$. So $R\simeq\Z_\ell\times\Z_\ell$ as rings, and the image of $x$ is a pair $(a,b)$ of roots of $g$ in $\Z_\ell$
reducing to $\bar a,\bar b$.
Modules: a discrete valuation ring and $\Z_\ell\times\Z_\ell$ are principal ideal rings. Let $T$ be as stated. If $R$ is a discrete valuation ring, $T$ is
finitely generated and torsion-free over $R$ (if $zt=0$ with $z\ne0$, put $z'=\Tr(z)-z$, with $\Tr$ and $N$ the trace and determinant of multiplication by $z$ on $R$; by Cayley-Hamilton
$zz'=N(z)$, which is nonzero as $R$ is a domain, so $N(z)t=z'zt=0$ and $t=0$), hence free (\Cref{ch:modules}), of rank $\rank_{\Z_\ell}T/\rank_{\Z_\ell}R=1$; in this case
$R\otimes\Q_\ell$ is a field and the hypothesis on $T\otimes\Q_\ell$ holds automatically. If $R=\Z_\ell e\times\Z_\ell(1-e)$, then $T=eT\oplus(1-e)T$, each summand free over $\Z_\ell$, of
ranks $\dim(e(T\otimes\Q_\ell))=1$ and $1$ by the hypothesis; a generator $t_1$ of $eT$ and $t_2$ of $(1-e)T$ give $T=R(t_1+t_2)$.
\end{proof}

\begin{theorem}[$\Z_3\lbrack\zeta_3\rbrack$]\label{thm:ant:z3zeta3dvr}
Let $\Z_3[\zeta_3]:=\Z_3[X]/(X^2+X+1)$, $\zeta_3$ the class of $X$, $\pi=1-\zeta_3$. Then:
\begin{enumerate}[label=(\roman*)]
\item $\Z_3[\zeta_3]$ is an integral domain, free over $\Z_3$ with basis $1,\zeta_3$ (or $1,\pi$), and $\Z_3[\zeta_3]\simeq\Z_3\otimes_\Z\Z[\zeta_3]$;
\item it is a discrete valuation ring with maximal ideal $(\pi)$, uniformiser $\pi$ and residue field $\F_3$; every nonzero element is $u\pi^k$ with $u$ a
unit and $k\ge0$;
\item $\pi^2=-3\zeta_3$, so $3\Z_3[\zeta_3]=\pi^2\Z_3[\zeta_3]$ and $v_\pi(3)=2$;
\item (stated; \Cref{ch:local}) it is the $\pi$-adic completion of $\Z[\zeta_3]$ and the ring of integers of $\Q_3(\zeta_3)=\Q_3(\sqrt{-3})$, a totally
ramified quadratic extension of $\Q_3$.
\end{enumerate}
\end{theorem}

\begin{proof}
With $g=X^2+X+1$ and $c=1$, $g(X+1)=X^2+3X+3$ is Eisenstein at $3$; \Cref{lem:ant:quadraticZl}(ii) gives (i) except the last isomorphism, (ii) with uniformiser
$\zeta_3-1=-\pi$, and $3\Z_3[\zeta_3]=\pi^2\Z_3[\zeta_3]$; the identity $\pi^2=-3\zeta_3$ is \Cref{prop:ant:above3}(i). The isomorphism
$\Z_3\otimes_\Z\Z[\zeta_3]=\Z_3\otimes_\Z\Z[X]/(X^2+X+1)\simeq\Z_3[X]/(X^2+X+1)$ is right exactness of the tensor product. A direct check of locality: if
$a+b\zeta_3\notin(\pi)$, i.e. $3\nmid a+b$, then $N(a+b\zeta_3)=a^2-ab+b^2\equiv(a+b)^2\not\equiv0\pmod3$ and $(a+b\zeta_3)^{-1}=(a+b\zeta_3^2)/(a^2-ab+b^2)$.
(iv) We do not prove this here; see \Cref{ch:local} and \cite[Chapter I, Section 6]{SerreLocalFields}.
\end{proof}

\begin{corollary}[the lattice of Lemma lem:IV]\label{cor:ant:z3zeta3modules}
Let $T$ be a free $\Z_3$-module of rank $2$ and $\sigma\in\Aut_{\Z_3}(T)$ with $\sigma^2+\sigma+1=0$ (for $\sigma\in\SL_2(\Z_3)$ of order $3$ this holds, its
eigenvalues being the primitive cube roots of unity, \Cref{ch:reps}). Make $T$ a $\Z_3[\zeta_3]$-module with $\zeta_3$ acting as $\sigma$. Then:
\begin{enumerate}[label=(\roman*)]
\item $T$ is free of rank one over $\Z_3[\zeta_3]$, and $\Z_3[\sigma]\simeq\Z_3[X]/(X^2+X+1)$;
\item $T/3T\simeq\Z_3[\zeta_3]/(\pi^2)$, and $\sigma$ acts on $T/3T$ as the nontrivial unipotent map of \Cref{prop:ant:above3}(iii), with fixed line $\pi T/3T$;
\item for every $k\ge1$ the $\sigma$-fixed vectors of $T/\pi^kT$ form the line $\pi^{k-1}T/\pi^kT$ over $\F_3$, and $T^\sigma=0$.
\end{enumerate}
\end{corollary}

\begin{proof}
(i) $\sigma$ is not a scalar: a scalar $s\in\Z_3$ with $s^2+s+1=0$ does not exist (\Cref{thm:ant:z3zeta3dvr}(i); directly, $s^2+s+1\equiv(s-1)^2\pmod3$ forces
$s=1+3t$, and then $s^2+s+1=3+9t+9t^2\equiv3\pmod9$ is not $0$). So the polynomials over $\Z_3$ killing $\sigma$ are exactly the multiples of $X^2+X+1$ (division
by the monic $X^2+X+1$ leaves a remainder $a+bX$ with $a+b\sigma=0$, and $b\ne0$ would make $\sigma=-a/b$ a scalar in $\Q_3$ satisfying $s^2+s+1=0$, hence in $\Z_3$ since $\Z_3$ is integrally closed, which was just excluded), and
$\Z_3[\sigma]\simeq\Z_3[X]/(X^2+X+1)$. $T$ is torsion-free over the domain $\Z_3[\zeta_3]$ by the norm argument of \Cref{lem:ant:quadraticZl}, hence
free, of rank $2/2=1$. (ii) $T\simeq\Z_3[\zeta_3]$ gives $T/3T\simeq\Z_3[\zeta_3]/(3)=\Z_3[\zeta_3]/(\pi^2)\simeq\Z[\zeta_3]/(\pi^2)$, and \Cref{prop:ant:above3}(iii)
applies. (iii) In $\Z_3[\zeta_3]/(\pi^k)$, $\zeta_3x=x$ iff $\pi x\equiv0\pmod{\pi^k}$ iff $x\in(\pi^{k-1})$ (write $x=u\pi^j$ with $u$ a unit,
\Cref{thm:ant:z3zeta3dvr}(ii)); and in $T$ itself $\pi x=0$ forces $x=0$, as $T$ is torsion-free over the domain $\Z_3[\zeta_3]$.
\end{proof}

A worked case: let $T=\Z_3^2$ and let $\sigma=\sm0{-1}1{-1}$ be the companion matrix of $X^2+X+1$ ($\sigma e_1=e_2$, $\sigma e_2=-e_1-e_2$). Then $e_1$ generates $T$
over $\Z_3[\sigma]$, as $e_1,\sigma e_1=e_2$ is a basis, and the isomorphism $\Z_3[\zeta_3]\to T$, $z\mapsto z\cdot e_1$, sends $1,\zeta_3$ to $e_1,e_2$. Modulo $3$,
$\sigma-1\equiv\sm2211$ has rank one and $(\sigma-1)^2\equiv\sm6633\equiv0$, so $\sigma$ is a nontrivial unipotent on $T/3T$; its fixed line is the kernel of
$\sm2211$, spanned by $e_1-e_2=(1-\sigma)e_1=\pi e_1$, which is the line $\pi T/3T$ of (ii).

\begin{proposition}[$\Z_\ell\lbrack\zeta_3\rbrack$ and $\Z_\ell\lbrack i\rbrack$]\label{prop:ant:zlzeta3}
For a prime $\ell$ put $\Z_\ell[\zeta_3]:=\Z_\ell[X]/(X^2+X+1)\simeq\Z_\ell\otimes_\Z\Z[\zeta_3]$ and $\Z_\ell[i]:=\Z_\ell[X]/(X^2+1)\simeq\Z_\ell\otimes_\Z\Z[i]$. These are
defined as these quotient rings, not as subrings of an algebraic closure of $\Q_\ell$.
\begin{enumerate}[label=(\roman*)]
\item $\ell=3$: $\Z_3[\zeta_3]$ is the ramified discrete valuation ring of \Cref{thm:ant:z3zeta3dvr}.
\item $\ell\equiv1\pmod3$: $X^2+X+1$ has two distinct roots in $\Z_\ell$, and $\Z_\ell[\zeta_3]\simeq\Z_\ell\times\Z_\ell$, which is not a domain.
\item $\ell\equiv2\pmod3$: $\Z_\ell[\zeta_3]$ is an unramified discrete valuation ring with uniformiser $\ell$ and residue field $\F_{\ell^2}$.
\item Likewise $\Z_\ell[i]\simeq\Z_\ell\times\Z_\ell$ for $\ell\equiv1\pmod4$, an unramified discrete valuation ring with residue field $\F_{\ell^2}$ for
$\ell\equiv3\pmod4$, and a ramified discrete valuation ring with uniformiser $1+i$ for $\ell=2$.
\item In all cases $\Q_\ell[\zeta_3]:=\Z_\ell[\zeta_3][1/\ell]=\Q_\ell\otimes_\Q\Q(\zeta_3)$ carries the $\Q_\ell$-linear involution $\zeta_3\mapsto\zeta_3^2$ (the
nontrivial element of $\Gal(\Q(\zeta_3)/\Q)$, extended by scalars), which negates $\sqrt{-3}=1+2\zeta_3$; and $\sqrt{-3}$ is a unit of $\Z_\ell[\zeta_3]$ iff
$\ell\ne3$. The same holds for $\Q_\ell[i]$ with $i\mapsto-i$, and $2i$ is a unit of $\Z_\ell[i]$ iff $\ell\ne2$.
\end{enumerate}
The trichotomy is the splitting of $\ell$ in $\Q(\zeta_3)$ and in $\Q(i)$ (\Cref{ex:ant:quadexamples}), and is the case $K=\Q(\zeta_3)$, $\Q(i)$ of
\Cref{prop:ant:otimesZl}.
\end{proposition}

\begin{proof}
(i) is \Cref{thm:ant:z3zeta3dvr}. For $\ell\ne3$, $X^2+X+1$ is separable modulo $\ell$ (its discriminant is $-3$), and it has a root in $\F_\ell$ iff $\F_\ell^\times$
contains an element of order $3$ iff $3\mid\ell-1$ (\Cref{ch:rings}). So (ii) and (iii) are \Cref{lem:ant:quadraticZl}(iii) and (i) (for $\ell=2$,
$X^2+X+1$ is irreducible over $\F_2$, and $2\equiv2\pmod 3$). (iv) The same with $X^2+1$, of discriminant $-4$: for odd $\ell$ it has a root modulo $\ell$
iff $4\mid\ell-1$; for $\ell=2$, $(X+1)^2+1=X^2+2X+2$ is Eisenstein. (v) The tensor product with $\Q_\ell$ of the automorphism $\zeta_3\mapsto\zeta_3^2$ of
$\Z[\zeta_3]$ preserves $\Z_\ell[\zeta_3]$ and sends $1+2\zeta_3$ to $1+2\zeta_3^2=-(1+2\zeta_3)$. $(\sqrt{-3})^2=-3$ is a unit iff $\ell\ne3$, and then
$\sqrt{-3}$ is a unit; for $\ell=3$ it lies in the maximal ideal $(\pi)$. Likewise $(2i)^2=-4$.
\end{proof}

\begin{proposition}[lattices with an automorphism of order $3$, $4$ or $6$]\label{prop:ant:latticefree}
Let $\ell$ be an odd prime, $T$ a free $\Z_\ell$-module of rank $2$, $m\in\{3,4,6\}$, and $g\in\Aut_{\Z_\ell}(T)$ with characteristic polynomial the
cyclotomic polynomial $\Phi_m$ ($X^2+X+1$, $X^2+1$, $X^2-X+1$). This holds for every $g\in\SL_2(\Z_\ell)$ of order $m$ (its eigenvalues are $\zeta$, $\zeta^{-1}$
with $\zeta$ a primitive $m$-th root of unity). Then $\Z_\ell[g]\simeq\Z_\ell[X]/(\Phi_m)$, which is $\Z_\ell[\zeta_3]$ for $m=3,6$ (via $X\mapsto\zeta_3$,
resp. $X\mapsto-\zeta_3$) and $\Z_\ell[i]$ for $m=4$; $T$ is free of rank one over $\Z_\ell[g]$; and the commutant of $g$ in $\End_{\Z_\ell}(T)$ is $\Z_\ell[g]$.
For $\ell=3$ and $m=3$ this is \Cref{cor:ant:z3zeta3modules}.
\end{proposition}

\begin{proof}
$g$ is not a scalar (its characteristic polynomial has distinct roots), so $\Phi_m$ is its minimal polynomial and $\Z_\ell[g]\simeq\Z_\ell[X]/(\Phi_m)=:R$; and
$\Phi_6(X)=\Phi_3(-X)$. By \Cref{prop:ant:zlzeta3}, $R$ is a discrete valuation ring or $\Z_\ell\times\Z_\ell$. Over $\Q_\ell$, $T\otimes\Q_\ell$ is a cyclic
$\Q_\ell[g]$-module (a non-scalar endomorphism of a plane has a cyclic vector: if every nonzero vector were an eigenvector, two independent eigenvectors and their
sum would have to share one eigenvalue and $g$ would be a scalar; so some $v$ is not an eigenvector, and $v,gv$ is a basis), so $T\otimes\Q_\ell\simeq R\otimes\Q_\ell$
via $z\mapsto zv$, and \Cref{lem:ant:quadraticZl} gives $T\simeq R$. Then the commutant of $g$ is $\End_R(T)=\End_R(R)=R$, since an $R$-linear endomorphism of
$R$ is multiplication by the image of $1$. (In the split case the idempotents are explicit: if $\zeta\in\Z_\ell$ is a root of $\Phi_m$, the other root is
$\zeta^{-1}$, and $e=(g-\zeta^{-1})/(\zeta-\zeta^{-1})$ is $1$ on the $\zeta$-eigenline and $0$ on the $\zeta^{-1}$-eigenline. It lies in $\Z_\ell[g]$ because its
denominator is a unit: $(\zeta-\zeta^{-1})^2=(\zeta+\zeta^{-1})^2-4$ is $1-4=-3$ for $m=3,6$ and $0-4=-4$ for $m=4$, units for $\ell$ odd, $\ell\ne3$; and the split
case requires $\ell\ne3$.)
\end{proof}

\begin{proposition}[the commutant and the invariant line of the adjoint]\label{prop:ant:commutant}
Let $F$ be a field of characteristic $0$, $V$ a two-dimensional $F$-vector space, and $\sigma\in\GL(V)$ with characteristic polynomial $X^2+X+1$ (so
$\sigma^3=1$, $\det\sigma=1$).
\begin{enumerate}[label=(\roman*)]
\item The commutant $\{x\in\End_F(V):x\sigma=\sigma x\}$ is $F[\sigma]\simeq F[X]/(X^2+X+1)=F\otimes_\Q\Q(\zeta_3)$, of dimension $2$; it is a field iff $F$
contains no primitive cube root of unity (for $F=\Q_\ell$: iff $\ell\not\equiv1\pmod3$). The same holds for $-\sigma$.
\item The trace-zero part of $F[\sigma]$ is the line $F\cdot(2\sigma+1)$, and $(2\sigma+1)^2=-3$: it is ``the line spanned by $\sqrt{-3}$''.
\item If $A\in\GL(V)$ satisfies $A\sigma A^{-1}=\sigma^q$ for a prime $q\ne3$, then conjugation by $A$ preserves $F[\sigma]$ and acts on it as
$\id_F\otimes\Frob_q$, where $\Frob_q\in\Gal(\Q(\zeta_3)/\Q)$ is $\zeta_3\mapsto\zeta_3^q$; hence it acts on the line $F(2\sigma+1)$ by $\chi_{-3}(\Frob_q)=\chi_{-3}(q)$,
and so does conjugation by $A^{-1}$.
\item If $\tau\in\GL(V)$ has characteristic polynomial $X^2+1$, the commutant is $F[\tau]\simeq F\otimes_\Q\Q(i)$, its trace-zero part is $F\tau$, and
$A\tau A^{-1}=\tau^q$ with $q$ an odd prime gives the action of conjugation by $A^{\pm1}$ on $F\tau$ by $\chi_{-4}(q)$.
\item Integral form: if $F=\Q_\ell$, $\ell$ odd, and $T\subset V$ is a $\sigma$-stable $\Z_\ell$-lattice, the commutant of $\sigma$ in $\End_{\Z_\ell}(T)$ is
$\Z_\ell[\sigma]\simeq\Z_\ell\otimes_\Z\Z[\zeta_3]$, and $T$ is free of rank one over it (\Cref{prop:ant:latticefree}); $\End_{\langle\sigma\rangle}(V)=\Q_\ell[\sigma]\simeq
\Q_\ell\otimes_\Q\Q(\zeta_3)$. For $\ell\equiv1\pmod3$ neither is a domain.
\end{enumerate}
\end{proposition}

\begin{proof}
(i) $\sigma$ is not a scalar, so $V$ has a cyclic vector $v$ ($v,\sigma v$ a basis); an $x$ commuting with $\sigma$ is determined by $x(v)=av+b\sigma v$, so
$x=a+b\sigma$. $F[X]/(X^2+X+1)$ is a field iff $X^2+X+1$ has no root in $F$; for $\Q_\ell$ use \Cref{prop:ant:zlzeta3}. (ii) $\tr(a+b\sigma)=2a-b$, which vanishes
iff $a+b\sigma=\frac b2(2\sigma+1)$; and $(2\sigma+1)^2=4(\sigma^2+\sigma+1)-3=-3$. (iii) $A F[\sigma]A^{-1}=F[\sigma^q]=F[\sigma]$, and conjugation is an
$F$-algebra automorphism sending $\sigma\mapsto\sigma^q$, which under $\sigma\leftrightarrow1\otimes\zeta_3$ is $\id\otimes(\zeta_3\mapsto\zeta_3^q)$. Explicitly
$A(2\sigma+1)A^{-1}=2\sigma^q+1$, which is $2\sigma+1$ if $q\equiv1\pmod3$ and $2\sigma^2+1=-(2\sigma+1)$ if $q\equiv2\pmod 3$; and $\chi_{-3}(q)=\pm1$
accordingly (\Cref{prop:ant:quadchar}(iii)). Since $\pm1$ is its own inverse, conjugation by $A^{-1}$ acts by the same sign. (iv) The same with $\tr(a+b\tau)=2a$ and
$\tau^q=\tau$ or $\tau^3=-\tau$ according to $q\bmod4$. (v) is \Cref{prop:ant:latticefree} with $m=3$, tensored with $\Q_\ell$ for $V$; and
\Cref{prop:ant:zlzeta3}(ii) for $\ell\equiv1\pmod3$.
\end{proof}

\begin{paperlink}\label{pl:ant:IV}
Lemma lem:IV and the paragraph after it: ``$T_3E$ is then free of rank one over $\Z_3[X]/(X^2+X+1)\simeq\Z_3[\zeta_3]$, with $X$ acting as $\sigma$'' is
\Cref{cor:ant:z3zeta3modules}(i), with $\sigma$ the generator of $\rho_{E,3}(I_q)$, of order $3$ and determinant $1$ (\Cref{ch:galrep}); here $\Z_3[\sigma]$, the
subring of $\End(T)$ generated by $\sigma$, is isomorphic to $\Z_3[X]/(X^2+X+1)$ because $\sigma$ is not a scalar. ``Its image in $\Aut(E[3])$ is a nontrivial
unipotent subgroup. Consequently $E[3]^{I_q}$ is a line'' is \Cref{cor:ant:z3zeta3modules}(ii): $E[3]=T/3T\simeq\Z_3[\zeta_3]/(\pi^2)$, $\sigma$ acts as
multiplication by $\zeta_3=1-\pi$, and its fixed space is $\pi T/3T$. That $\Z_3[\zeta_3]$ is a discrete valuation ring with uniformiser $\pi=1-\zeta_3$ and
residue field $\F_3$ is \Cref{thm:ant:z3zeta3dvr}; ``$\pi^2=-3$ up to a unit'' is \Cref{prop:ant:above3}(i). The same freeness is used in Lemma lem:tamadd,
with the basis in which $\sigma=\sm0{-1}1{-1}$, the companion matrix of $X^2+X+1$ acting on the basis $t,\sigma t$ for a generator $t$.
The paper's sentence attaches the freeness to $q\ge5$, but it holds at every $q\ne3$ of type $\kod{IV}$ or $\kod{IV}^*$, $q=2$ included, and wherever an element of
order $3$ generates the $3$-Sylow subgroup of $\rho(I_q)$, as in Lemma lem:tamadd: \Cref{cor:ant:z3zeta3modules} needs only $\sigma^2+\sigma+1=0$. Only the
identification of the order $3$ with Kraus's semistability defect needs $q\ge5$.
\end{paperlink}

\begin{paperlink}\label{pl:ant:euler}
Lemma lem:euler, types $\kod{II}$, $\kod{II}^*$, $\kod{IV}$, $\kod{IV}^*$: ``the invariants form the line spanned by $2\sigma+1$, whose square is $-3$, inside
$\End_{I_q}(T)=\Z_\ell[\sigma]\simeq\Z_\ell\otimes\Z[\zeta_3]$ (not a field when $\ell\equiv1\pmod3$)''. The invariants of $I_q$ on $\ad V$ are the trace-zero
elements of $\End_{I_q}(V)$, the commutant of $\rho(I_q)=\langle\pm\sigma\rangle$, which is the line $\Q_\ell(2\sigma+1)$ of \Cref{prop:ant:commutant}(ii); the
integral statement is \Cref{prop:ant:commutant}(v). ``Frobenius ... acts on the line by $\chi_{-3}(q)$'' is \Cref{prop:ant:commutant}(iii) with $A=\rho(F)$
supplied by \Cref{cor:ant:frobsigma}. Types $\kod{III}$, $\kod{III}^*$: the line spanned by $\tau$ and the sign $\chi_{-4}(q)$, \Cref{prop:ant:commutant}(iv); no
diagonalisation of $\tau$ over $\Q_\ell$ is needed (it is impossible for $\ell\equiv3\pmod4$, $\Q_\ell[i]$ being a field). The table entries
$1-\chi_{-3}(q)q^{-s}$ and $1-\chi_{-4}(q)q^{-s}$ follow with either Frobenius (\Cref{conv:ant:frobenius}). The original version of the note wrote ``$\Z_\ell[\zeta_3]\simeq
\End_{I_q}(V)$''; the correct reading, adopted in the revised version, is that $\End_{I_q}(V)$ is the $\Q_\ell$-algebra $\Q_\ell\otimes\Q(\zeta_3)$ and
$\Z_\ell\otimes\Z[\zeta_3]$ is the endomorphism ring $\End_{I_q}(T)$ of the lattice.
\end{paperlink}

\section{Cyclotomic fields and class field theory over \texorpdfstring{$\Q$}{Q}}\label{sec:ant:cyclotomic}

The cyclotomic fields are where the characters of the paper live: $\chi_\ell$ and $\bar\chi_\ell$ describe the action on $\mu_{\ell^\infty}$, the quadratic
characters $\chi_{-3}$, $\chi_{-4}$, $\chi_{q^*}$, $\chi_S$ factor through $\Gal(\Q(\zeta_n)/\Q)\simeq(\Z/n)^\times$, and the theorem of Kronecker and Weber
says that every character of $G_\Q$ of finite order does. This section computes the splitting of primes in $\Q(\zeta_n)$, shows that
$\Q(\sqrt{\ell^*})$ is the quadratic subfield of $\Q(\zeta_\ell)$ with its splitting law, and sets up the dictionary between Dirichlet characters and
characters of $G_\Q$ that the paper uses in both directions: as Galois characters ($\rho\otimes\chi$ in Proposition prop:ns and Lemma lem:h0) and as
Dirichlet characters ($a_n(g\otimes\chi)=\chi(n)a_n(g)$ in Lemma lem:twist(iii) and Theorems thm:twist and thm:qual). The Galois theory of $\Q(\zeta_n)$,
$\Gal(\Q(\zeta_n)/\Q)\simeq(\Z/n)^\times$, $a\mapsto\sigma_a$ with $\sigma_a(\zeta_n)=\zeta_n^a$, is in \Cref{ch:galois}; the cyclotomic polynomials $\Phi_n$ and
their factorisation modulo $p$ are in \Cref{ch:rings}.

\begin{theorem}[rings of integers of cyclotomic fields]\label{thm:ant:cyclotomicintegers}
Let $n\ge1$ and $K=\Q(\zeta_n)$, of degree $\varphi(n)$.
\begin{enumerate}[label=(\roman*)]
\item $\mathcal O_K=\Z[\zeta_n]$.
\item If $n=\ell^a$ with $\ell$ prime and $\zeta=\zeta_{\ell^a}$, then $\prod_{b\in(\Z/\ell^a)^\times}(1-\zeta^b)=\Phi_{\ell^a}(1)=\ell$, each $1-\zeta^b$ is a unit times
$1-\zeta$, and $\ell\mathcal O_K=(1-\zeta)^{\varphi(\ell^a)}\mathcal O_K$: $\ell$ is totally ramified, and $(1-\zeta)$ is a prime with residue field $\F_\ell$.
\item $d_K=(-1)^{\varphi(n)/2}\,n^{\varphi(n)}\big/\prod_{p\mid n}p^{\varphi(n)/(p-1)}$; in particular $d(\Q(\zeta_\ell))=(-1)^{(\ell-1)/2}\ell^{\ell-2}$ for an odd prime
$\ell$, $d(\Q(\zeta_3))=-3$, $d(\Q(\zeta_4))=-4$, $d(\Q(\zeta_8))=256$.
\item For $n\not\equiv2\pmod4$, a prime $p$ ramifies in $\Q(\zeta_n)$ iff $p\mid n$. (Recall $\Q(\zeta_n)=\Q(\zeta_{2n})$ for $n$ odd.)
\end{enumerate}
\end{theorem}

\begin{proof}
We prove (i) and (ii) for $n=\ell^a$, (iii) for $n$ an odd prime, and (iv) after \Cref{thm:ant:cyclotomicsplitting}; the general case of (i) and (iii)
is stated, see \cite[Chapter 2]{Washington} or \cite[Chapter I, Section 10]{Neukirch} (the formula of (iii) was checked against PARI's \texttt{nfdisc} for
$n\in\{3,4,5,7,8,9,12,15,16,20\}$). (ii) $\Phi_{\ell^a}(X)=\sum_{j=0}^{\ell-1}X^{j\ell^{a-1}}=\prod_b(X-\zeta^b)$, so $\prod_b(1-\zeta^b)=\Phi_{\ell^a}(1)=\ell$. For
$\ell\nmid b$, $(1-\zeta^b)/(1-\zeta)=1+\zeta+\dots+\zeta^{b-1}\in\Z[\zeta]$, and its inverse $(1-\zeta)/(1-\zeta^b)$ is of the same form since $\zeta=(\zeta^b)^{b'}$ with
$bb'\equiv1$. So $\ell=u(1-\zeta)^{\varphi(\ell^a)}$ with $u$ a unit. Writing $(1-\zeta)=\prod\mathfrak P^{a_{\mathfrak P}}$, $\ell\mathcal O_K=\prod\mathfrak P^{\varphi(\ell^a)a_{\mathfrak P}}$
and $\sum_{\mathfrak P}\varphi(\ell^a)a_{\mathfrak P}f_{\mathfrak P}=\varphi(\ell^a)$ (\Cref{thm:ant:efg}) forces a single $\mathfrak P$ with $a_{\mathfrak P}=f_{\mathfrak P}=1$.
(i) for $n=\ell^a$: differentiating $(X^{\ell^{a-1}}-1)\Phi_{\ell^a}(X)=X^{\ell^a}-1$ and evaluating at $\zeta$, where $\Phi_{\ell^a}(\zeta)=0$ kills the term
containing $\Phi_{\ell^a}$, gives $(\xi-1)\Phi'_{\ell^a}(\zeta)=\ell^a\zeta^{\ell^a-1}=\ell^a\zeta^{-1}$ with $\xi=\zeta^{\ell^{a-1}}$ a primitive $\ell$-th
root of unity, whose norm $N_{K/\Q}(\xi-1)$ is $\pm$ a power of $\ell$ by (ii). So $\disc(\Phi_{\ell^a})=\pm N(\Phi'_{\ell^a}(\zeta))$ is $\pm$ a power of $\ell$
(\Cref{thm:ant:integralbasis}(iv)), and the index $[\mathcal O_K:\Z[\zeta]]$ is a power of $\ell$. Suppose $\ell$ divides it. Then some $x\in\mathcal O_K\setminus\Z[\zeta]$
has $\ell x\in\Z[\zeta]$. With $\pi_\ell=1-\zeta$, the elements $1,\pi_\ell,\dots,\pi_\ell^{\varphi-1}$ ($\varphi=\varphi(\ell^a)$) form a $\Z$-basis of $\Z[\zeta]$; write
$\ell x=\sum_ic_i\pi_\ell^i$ with not all $c_i$ divisible by $\ell$, and let $j$ be the least index with $\ell\nmid c_j$. Then $\sum_{i\ge j}c_i\pi_\ell^i=\ell x-\sum_{i<j}c_i\pi_\ell^i$
lies in $\ell\mathcal O_K=\pi_\ell^\varphi\mathcal O_K$ (the terms with $i<j$ have $c_i\in\ell\Z$, and $\ell x\in\ell\mathcal O_K$), so its valuation $v_{\pi_\ell}$,
the valuation at the prime $(\pi_\ell)$, is at least $\varphi$. Here $j<\varphi$: otherwise all $c_i$ ($i<\varphi$) would be divisible by $\ell$. But
$v_{\pi_\ell}(c_j\pi_\ell^j)=j$, because an integer prime to $\ell$ lies in no prime above $\ell$, and $(\pi_\ell)$ is the only such prime, so $v_{\pi_\ell}(c_j)=0$;
while the terms with $i>j$ have valuation $\ge i>j$. If $v(a)<v(b)$ then $v(a+b)=v(a)$, since $v(a+b)\ge v(a)$ and $v(a)=v((a+b)-b)\ge\min(v(a+b),v(b))$ (\Cref{thm:ant:idealfactorisation}(iv)); so the sum
has valuation exactly $j<\varphi$:
a contradiction. Hence $\mathcal O_K=\Z[\zeta]$.
(iii) for an odd prime $\ell$: with $a=1$, $(\zeta-1)\Phi_\ell'(\zeta)=\ell\zeta^{\ell-1}$, $N(\zeta-1)=\prod_b(\zeta^b-1)=\ell$ (as $\ell-1$ is even) and
$N(\ell\zeta^{\ell-1})=\ell^{\ell-1}N(\zeta)^{\ell-1}=\ell^{\ell-1}$ ($\zeta$ is a unit of norm $\prod_b\zeta^b=\zeta^{\ell(\ell-1)/2}=1$), so by multiplicativity
$N(\Phi'_\ell(\zeta))=\ell^{\ell-1}/\ell=\ell^{\ell-2}$, and \Cref{thm:ant:integralbasis}(iv) with $n=\ell-1$ gives $d=(-1)^{(\ell-1)(\ell-2)/2}\ell^{\ell-2}$. Since
$(\ell-1)(\ell-2)/2=\frac{\ell-1}2(\ell-2)$ with $\ell-2$ odd, the sign is $(-1)^{(\ell-1)/2}$. For $\ell=3,5,7$ this gives $-3$, $125$ and $-7^5=-16807$.
\end{proof}

\begin{theorem}[splitting in cyclotomic fields]\label{thm:ant:cyclotomicsplitting}
Let $n=p^am$ with $p\nmid m$, and let $f$ be the order of $p$ in $(\Z/m)^\times$. In $\Q(\zeta_n)$ the prime $p$ has $e=\varphi(p^a)$, residue degree $f$, and
$g=\varphi(m)/f$ primes above it. Under $(\Z/n)^\times=(\Z/p^a)^\times\times(\Z/m)^\times$ the inertia group of $p$ is $(\Z/p^a)^\times\times\{1\}=\Gal(\Q(\zeta_n)/\Q(\zeta_m))$
and the decomposition group is $(\Z/p^a)^\times\times\langle p\rangle$. For $p\nmid n$, $\Frob_p=\sigma_p$; so $p$ splits completely iff $p\equiv1\pmod n$, and $p$ is
inert iff $p$ generates $(\Z/n)^\times$ (possible only if this group is cyclic).
\end{theorem}

\begin{proof}
First let $p\nmid n$. $\disc(\Phi_n)$ divides $\disc(X^n-1)=\pm n^n$ (the discriminant of a product is the product of the discriminants times a square of a
resultant), so $p\nmid\disc(\Phi_n)$; then $p$ is unramified (\Cref{thm:ant:discram}, as $d_K\mid\disc(\Phi_n)$) and $p\nmid[\mathcal O_K:\Z[\zeta_n]]$
(\Cref{thm:ant:integralbasis}(iii)). Modulo $p$, $\Phi_n$ is a product of $\varphi(n)/f$ distinct irreducible factors of degree $f$, the order of $p$ modulo $n$
(\Cref{ch:rings}); \Cref{thm:ant:dedekindkummer} gives the splitting. For $x=\sum_ia_i\zeta_n^i\in\Z[\zeta_n]$,
$\sigma_p(x)=\sum_ia_i\zeta_n^{pi}\equiv(\sum_ia_i\zeta_n^i)^p\pmod{p}$ (the binomial theorem modulo $p$ and $a_i^p\equiv a_i$). A general $x\in\mathcal O_K$ has $kx\in\Z[\zeta_n]$
for the index $k$, which is prime to $p$, and then $k\sigma_p(x)\equiv(kx)^p\equiv kx^p\pmod{\mathfrak P}$; so $\sigma_p(x)\equiv x^p\pmod{\mathfrak P}$ for every
$\mathfrak P\mid p$, and $\sigma_p=\Frob_p$ (\Cref{prop:ant:frobenius}). It has order $f$; it is trivial iff $p\equiv1\pmod n$, and generates the group iff $p$ is
inert. Now let $p\mid n$. In $\Q(\zeta_m)$, $p$ has $(e,f,g)=(1,f,\varphi(m)/f)$; in $\Q(\zeta_{p^a})$ it has $(\varphi(p^a),1,1)$ (\Cref{thm:ant:cyclotomicintegers}(ii)).
Write $(e',f',g')$ for the invariants of $p$ in $\Q(\zeta_n)$, which contains both fields. By multiplicativity in towers (\Cref{thm:ant:efg}), $e'$ is a
multiple of $\varphi(p^a)$ (pass through $\Q(\zeta_{p^a})$) and $f'$ a multiple of $f$ (pass through $\Q(\zeta_m)$), and $g'\ge\varphi(m)/f$ because each of the
$\varphi(m)/f$ primes of $\Q(\zeta_m)$ above $p$ has at least one prime of $\Q(\zeta_n)$ above it, and these are distinct. As
$e'f'g'=\varphi(n)=\varphi(p^a)\varphi(m)=\varphi(p^a)\cdot f\cdot\varphi(m)/f$, all three inequalities are equalities: $e'=\varphi(p^a)$, $f'=f$, $g'=\varphi(m)/f$. The inertia group has order $\varphi(p^a)$, and its image
in $\Gal(\Q(\zeta_m)/\Q)$ is the inertia group there, which is trivial (\Cref{thm:ant:residuegalois}(ii)); so it is $\Gal(\Q(\zeta_n)/\Q(\zeta_m))$, of order
$\varphi(p^a)$. The decomposition group contains it and maps onto the decomposition group $\langle p\rangle$ of $\Q(\zeta_m)$, so it is the full preimage
$(\Z/p^a)^\times\times\langle p\rangle$.
\end{proof}

\noindent\emph{Proof of \Cref{thm:ant:cyclotomicintegers}(iv).} By \Cref{thm:ant:cyclotomicsplitting}, $p$ ramifies iff $p\mid n$ and $\varphi(p^a)>1$, and
$\varphi(p^a)=1$ only for $p^a=2$, excluded when $n\not\equiv2\pmod4$.\hfill$\square$

\begin{example}\label{ex:ant:cyclotomicexamples}
(Checked with PARI.) $n=3$: $p\equiv1\pmod3$ splits, $p\equiv2\pmod3$ is inert, $3$ is ramified. $n=5$: $11$ splits completely; $19$ and $29$ (order $2$
modulo $5$) give $f=g=2$; $2$ and $3$ (order $4$) are inert. $n=7$: $2$ has order $3$ modulo $7$, so $f=3$, $g=2$; $13$ has order $2$, so $f=2$, $g=3$; $29\equiv1$
splits completely. $n=8$: $2$ is totally ramified; $3$, $5$, $7$ have $f=g=2$; $17$ splits completely; $(\Z/8)^\times$ is not cyclic, so no prime is inert.
\end{example}

\begin{proposition}[$\Q(\sqrt{\ell^*})$]\label{prop:ant:sqrtlstar}
Let $\ell$ be an odd prime and $\ell^*=(-1)^{(\ell-1)/2}\ell$.
\begin{enumerate}[label=(\roman*)]
\item $\Q(\sqrt{\ell^*})$ is the unique quadratic subfield of $\Q(\zeta_\ell)$ (\Cref{ch:galois}); it is ramified only at $\ell$, with discriminant $\ell^*$.
\item For every prime $p\ne\ell$, including $p=2$, $\chi_{\ell^*}(\Frob_p)=\bigl(\frac p\ell\bigr)$: $p$ splits in $\Q(\sqrt{\ell^*})$ iff $p$ is a square modulo $\ell$, and
$p$ is \emph{inert} iff $p$ is a non-square modulo $\ell$. The inert primes have density $\frac12$ (\Cref{cor:ant:chebotarev}).
\item For $\ell=3$, $\Q(\sqrt{-3})=\Q(\zeta_3)$, and the inert primes are the $p\equiv2\pmod3$.
\item $\Q(\sqrt{\ell^*})$ is imaginary iff $\ell\equiv3\pmod4$.
\end{enumerate}
\end{proposition}

\begin{proof}
(i) $\Gal(\Q(\zeta_\ell)/\Q)\simeq\F_\ell^\times$ is cyclic of even order $\ell-1$, so it has a unique subgroup of index $2$ and $\Q(\zeta_\ell)$ has a unique quadratic
subfield $M$ (\Cref{ch:galois}). A prime $p\ne\ell$ is unramified in $\Q(\zeta_\ell)$ (\Cref{thm:ant:cyclotomicintegers}(iv)), hence in $M$ (ramification indices
multiply in towers, \Cref{thm:ant:efg}); and some prime ramifies in $M\ne\Q$ (\Cref{thm:ant:discram}). So $M$ is ramified exactly at $\ell$, and
\Cref{prop:ant:ramonlyat}(i) gives $M=\Q(\sqrt{\ell^*})$, of discriminant $\ell^*$. For example $\Q(\zeta_5)\supset\Q(\sqrt5)$, as $\zeta_5+\zeta_5^{-1}=\frac{-1+\sqrt5}2$, and
$\Q(\zeta_7)\supset\Q(\sqrt{-7})$. (ii) In $\Gal(\Q(\zeta_\ell)/\Q)\simeq\F_\ell^\times$, $\Frob_p=\sigma_p$ corresponds to $p\bmod\ell$
(\Cref{thm:ant:cyclotomicsplitting}). Its restriction to $\Q(\sqrt{\ell^*})$ is $\Frob_p$ there (\Cref{prop:ant:frobenius}(v)), which is trivial iff $\sigma_p$ lies in
the unique subgroup of index $2$ of the cyclic group $\F_\ell^\times$, the squares. (iii) $\zeta_3=(-1+\sqrt{-3})/2$. (iv) $\ell^*<0$ iff $(\ell-1)/2$ is odd.
\end{proof}

\begin{corollary}[quadratic reciprocity, second proof]\label{cor:ant:qr}
For distinct odd primes $p,q$: $\bigl(\frac{q^*}p\bigr)=\bigl(\frac pq\bigr)$, equivalently $\bigl(\frac pq\bigr)\bigl(\frac qp\bigr)=(-1)^{\frac{p-1}2\frac{q-1}2}$. Moreover
$\bigl(\frac{-1}p\bigr)=(-1)^{(p-1)/2}$ and $\bigl(\frac2p\bigr)=(-1)^{(p^2-1)/8}$.
\end{corollary}

\begin{proof}
$\chi_{q^*}(\Frob_p)$ equals $\bigl(\frac{q^*}p\bigr)$ by \Cref{prop:ant:quadchar}(i) together with \Cref{thm:ant:quadsplit} (whose proof uses only the meaning of the
Legendre symbol, not reciprocity), and $\bigl(\frac pq\bigr)$ by \Cref{prop:ant:sqrtlstar}(ii). Since $\bigl(\frac{q^*}p\bigr)=\bigl(\frac{-1}p\bigr)^{(q-1)/2}\bigl(\frac qp\bigr)$,
the second form follows from the first supplement. For the supplements: $\Q(i)=\Q(\zeta_4)$, and $\sigma_p(i)=i^p=(-1)^{(p-1)/2}i$, so $\bigl(\frac{-4}p\bigr)=
\chi_{-4}(\Frob_p)=(-1)^{(p-1)/2}$; and $\sqrt2=\zeta_8+\zeta_8^{-1}\in\Q(\zeta_8)$, with $\sigma_p(\sqrt2)=\zeta_8^p+\zeta_8^{-p}$, which is $\sqrt2$ if $p\equiv\pm1\pmod8$
and $-\sqrt2$ if $p\equiv\pm3\pmod8$; so $\bigl(\frac8p\bigr)=\bigl(\frac2p\bigr)=\chi_8(\Frob_p)=(-1)^{(p^2-1)/8}$.
\end{proof}

\Cref{ch:rings} proves reciprocity by Gauss sums; the proof above explains why the character of $\Q(\sqrt{q^*})$ is the Legendre symbol modulo $q$, which is
what Lemma lem:twist means by ``$\chi_{q^*}$, the quadratic character of conductor $q$''.

\begin{theorem}[Kronecker-Weber]\label{thm:ant:kroneckerweber}
Every finite abelian extension $K/\Q$ is contained in $\Q(\zeta_n)$ for some $n$. The least such $n$ with $n\not\equiv2\pmod4$ is the \emph{conductor} $f_K$
of $K$, and the primes dividing $f_K$ are exactly the primes ramified in $K$. Equivalently, the maximal abelian extension of $\Q$ is
$\Q^{\mathrm{ab}}=\Q(\mu_\infty)=\bigcup_n\Q(\zeta_n)$, and the cyclotomic characters give $G_\Q^{\mathrm{ab}}=\Gal(\Q^{\mathrm{ab}}/\Q)\simeq\varprojlim_n(\Z/n)^\times=
\hat\Z^\times=\prod_\ell\Z_\ell^\times$. For a quadratic field of discriminant $D$, $f_K=\abs D$.
\end{theorem}

\begin{proof}[Reference]
We do not prove this; see \cite[Chapter 14]{Washington} or \cite[Chapter V, Section 1, and Chapter VI]{Neukirch} (chapter-level pointers). For a
quadratic field the inclusion $K\subseteq\Q(\zeta_{\abs D})$ is elementary (\Cref{ch:galois}), and $f_K=\abs D$ is \Cref{thm:ant:dirichletgalois}(i) applied to
the primitive character $n\mapsto\bigl(\frac Dn\bigr)$ of conductor $\abs D$ (\Cref{ch:rings}), which corresponds to $\chi_D$ by \Cref{prop:ant:quadchar}(i).
\end{proof}

\begin{theorem}[Artin reciprocity over $\Q$, explicit form]\label{thm:ant:artinQ}
Let $K/\Q$ be abelian with $K\subseteq\Q(\zeta_n)$.
\begin{enumerate}[label=(\roman*)]
\item The \emph{Artin map} $\Art_K\colon(\Z/n)^\times\to\Gal(K/\Q)$, $a\mapsto\sigma_a|_K$, is surjective with kernel $H_K$, and $K\mapsto H_K$ is an
inclusion-reversing bijection between the subfields of $\Q(\zeta_n)$ and the subgroups of $(\Z/n)^\times$.
\item For every prime $p\nmid n$, $\Frob_p(K/\Q)=\Art_K(p\bmod n)$: the Frobenius of an unramified prime in an abelian extension depends only on $p\bmod n$.
So $p\nmid n$ splits completely in $K$ iff $p\bmod n\in H_K$, and the residue degree of $p$ in $K$ is the order of $p$ in $(\Z/n)^\times/H_K$.
\item (Stated.) In idelic form, the global Artin map $\Q^\times\backslash\A_\Q^\times\to G_\Q^{\mathrm{ab}}$ is surjective with kernel the connected
component of the identity, and $\A_\Q^\times=\Q^\times\cdot(\R_{>0}\times\hat\Z^\times)$. In the normalisation in which a uniformiser at $p$ maps to an
arithmetic Frobenius, an idele $(1,\dots,1,u_p,1,\dots)$ with $u_p\in\Z_p^\times$ maps to an element of the inertia group at $p$ acting on $\mu_{p^\infty}$ by
$u_p^{-1}$. Deligne's normalisation sends uniformisers to geometric Frobenius elements, and then the sign is reversed.
\end{enumerate}
\end{theorem}

\begin{proof}
(i) is Galois theory (\Cref{ch:galois}) for $\Gal(\Q(\zeta_n)/\Q)\simeq(\Z/n)^\times$. (ii) By \Cref{prop:ant:frobenius}(v) and \Cref{thm:ant:cyclotomicsplitting},
$\Frob_p(K/\Q)=\Frob_p(\Q(\zeta_n)/\Q)|_K=\sigma_p|_K$; the residue degree is the order of $\Frob_p$ (\Cref{prop:ant:frobenius}(i)). (iii) We do not prove
this; see \cite[Chapter VI, Sections 5 and 6]{Neukirch} or \cite[Chapter V]{MilneCFT}; nothing in the paper needs it.
\end{proof}

\begin{theorem}[Dirichlet characters and characters of $G_\Q$]\label{thm:ant:dirichletgalois}
\begin{enumerate}[label=(\roman*)]
\item Let $\chi$ be a primitive Dirichlet character, of conductor $f_\chi$ (\Cref{ch:rings}), and let $\psi_\chi$ be the composite
\[
G_\Q\to\Gal(\Q(\zeta_{f_\chi})/\Q)\simeq(\Z/f_\chi)^\times\xrightarrow{\ \chi\ }\C^\times.
\] Then $\psi_\chi(\Frob_p)=\chi(p)$ for every prime $p\nmid f_\chi$, and
$\chi\mapsto\psi_\chi$ is a bijection from primitive Dirichlet characters to continuous homomorphisms $G_\Q\to\C^\times$ (these have finite image,
\Cref{ch:galois}). Each $\psi$ is determined by its values at the $\Frob_p$, $p\nmid f_\chi$. The field cut out by $\psi_\chi$ has conductor $f_\chi$, and
$\psi_{\chi\chi'}=\psi_\chi\psi_{\chi'}$ when $\chi\chi'$ denotes the primitive character inducing the product.
\item $\psi_\chi$ is unramified at $p$ iff $p\nmid f_\chi$. More precisely, if $f_\chi=p^am$ with $p\nmid m$, then $\psi_\chi(I_p)=\chi\bigl((\Z/p^a)^\times\times\{1\}\bigr)$.
\item The same holds for characters with values in $\Qbar_\ell^\times$ or $\Fbar_\ell^\times$ of finite order (through a fixed embedding of roots of unity),
and every continuous $\psi\colon G_\Q\to\Z_\ell^\times$ factors through $G_\Q^{\mathrm{ab}}\simeq\hat\Z^\times$; for instance $\chi_\ell$ is the projection
$\hat\Z^\times\to\Z_\ell^\times$, and $\bar\chi_\ell=\psi_{\omega_\ell}$ for the character $\omega_\ell\colon a\mapsto a\bmod\ell$ of $(\Z/\ell)^\times$.
\item (Stated; \Cref{ch:hecke,ch:modgalrep}.) Under this dictionary, twisting a newform's coefficients, $a_n\mapsto\chi(n)a_n$, corresponds to tensoring its
Galois representation with $\psi_\chi$.
\end{enumerate}
\end{theorem}

\begin{proof}
(i) $\psi_\chi(\Frob_p)=\chi(p)$ by \Cref{thm:ant:cyclotomicsplitting}. Injectivity: $G_\Q\to(\Z/f_\chi)^\times$ is surjective, so $\psi_\chi$ determines $\chi$
on $(\Z/f_\chi)^\times$, and $\psi_\chi$ determines $f_\chi$ as below. Surjectivity: a continuous $\psi\colon G_\Q\to\C^\times$ has finite abelian image,
and its kernel cuts out a finite abelian extension, contained in some $\Q(\zeta_n)$ (\Cref{thm:ant:kroneckerweber}); so $\psi$ factors as
$G_\Q\to(\Z/n)^\times\xrightarrow{\chi'}\C^\times$, and $\psi=\psi_\chi$ for the primitive $\chi$ inducing $\chi'$ (restriction from $\Q(\zeta_n)$ to $\Q(\zeta_{f_\chi})$
is reduction $(\Z/n)^\times\to(\Z/f_\chi)^\times$). That $\psi$ is determined by its values at Frobenius elements outside a finite set follows from the density
of Frobenius elements (\Cref{cor:ant:chebotarev}), or directly from Dirichlet's theorem on primes in progressions. The field $K$ cut out by $\psi_\chi$ lies
in $\Q(\zeta_{f_\chi})$; if it lay in $\Q(\zeta_{f'})$ with $f'\mid f_\chi$, $f'<f_\chi$, then $\chi$ would be trivial on the kernel of $(\Z/f_\chi)^\times\to(\Z/f')^\times$,
hence induced from a character modulo $f'$, contradicting primitivity; and if $K\subseteq\Q(\zeta_n)$ then $K\subseteq\Q(\zeta_n)\cap\Q(\zeta_{f_\chi})=\Q(\zeta_{\gcd(n,f_\chi)})$
(the intersection contains $\Q(\zeta_{\gcd})$ and has degree $\varphi(n)\varphi(f_\chi)/\varphi(\lcm)=\varphi(\gcd)$, by Galois theory of the compositum
$\Q(\zeta_{\lcm})$). So the least $n$ is $f_\chi$, which is $\not\equiv2\pmod4$ for a primitive character. Multiplicativity is clear on a common modulus.
(ii) By \Cref{thm:ant:cyclotomicsplitting} the image of $I_p$ in $\Gal(\Q(\zeta_{f_\chi})/\Q)$ is $(\Z/p^a)^\times\times\{1\}$; $\chi$ is nontrivial on it when $a\ge1$,
by primitivity (otherwise $\chi$ would factor through $(\Z/m)^\times$). (iii) The same proofs; the factorisation through $\hat\Z^\times$ is
\Cref{thm:ant:kroneckerweber}, since a continuous homomorphism to the abelian group $\Z_\ell^\times$ factors through $G_\Q^{\mathrm{ab}}$. $\bar\chi_\ell(\Frob_p)=p\bmod
\ell=\omega_\ell(p)$. (iv) We do not prove this here; see \Cref{ch:hecke,ch:modgalrep}.
\end{proof}

\begin{example}[the paper's characters in both languages]\label{ex:ant:characters}
For $p$ unramified (arithmetic Frobenius):
\begin{center}\footnotesize\setlength{\tabcolsep}{3pt}
\begin{tabular}{lllll}
\toprule
Galois character & field cut out & Dirichlet character & conductor & value at $\Frob_p$\\
\midrule
$\chi_{-4}$ & $\Q(i)=\Q(\zeta_4)$ & the nontrivial character mod $4$ & $4$ & $(-1)^{(p-1)/2}$\\
$\chi_{-3}=\bar\chi_3^{\mathrm{cyc}}$ & $\Q(\sqrt{-3})=\Q(\zeta_3)$ & $\bigl(\frac\cdot3\bigr)$ & $3$ & $1$ iff $p\equiv1\pmod3$\\
$\chi_{q^*}$, $q$ odd & $\Q(\sqrt{q^*})\subseteq\Q(\zeta_q)$ & $\bigl(\frac\cdot q\bigr)$ & $q$ & $\bigl(\frac pq\bigr)$\\
$\chi_S=\prod_{q\in S}\chi_{q^*}$ & $\Q\bigl(\sqrt{\prod_{q\in S}q^*}\bigr)$ & $\prod_{q\in S}\bigl(\frac\cdot q\bigr)$ & $\prod_{q\in S}q$ & $\prod_{q\in S}\bigl(\frac pq\bigr)$\\
$\chi_{\ell^*}=\bar\chi_\ell^{(\ell-1)/2}$ & $\Q(\sqrt{\ell^*})$ & $\bigl(\frac\cdot\ell\bigr)$ & $\ell$ & $\bigl(\frac p\ell\bigr)$\\
$\bar\chi_\ell$ (values in $\F_\ell^\times$) & $\Q(\zeta_\ell)$ & $\omega_\ell\colon a\mapsto a\bmod\ell$ & $\ell$ & $p\bmod\ell$\\
\bottomrule
\end{tabular}
\end{center}
The identity $\chi_{\ell^*}=\bar\chi_\ell^{(\ell-1)/2}$ means: the quadratic character $\bigl(\frac\cdot\ell\bigr)$ equals $a\mapsto a^{(\ell-1)/2}\bmod\ell$ (Euler's criterion),
read in $\{\pm1\}\subset\F_\ell^\times$. The Dirichlet character $\bigl(\frac\cdot q\bigr)$ equals $\bigl(\frac{q^*}\cdot\bigr)$ by quadratic reciprocity (\Cref{ch:rings}), so the two
descriptions of the column agree.
\end{example}

\begin{warning}\label{warn:ant:chil}
$\chi_{\ell^*}$ and $\bar\chi_\ell$ coincide only for $\ell=3$: for $\ell\ge5$, $\bar\chi_\ell$ has order $\ell-1>2$. The cyclotomic character $\chi_\ell$ itself has
infinite image $\Z_\ell^\times$. So ``$\rho\simeq\rho\otimes\chi_\ell$'' (impossible for every two-dimensional $\rho$: taking determinants, $\chi_\ell^2=1$, while $\chi_\ell(G_\Q)=\Z_\ell^\times$ is
infinite), ``$\bar\rho\simeq\bar\rho\otimes\bar\chi_\ell$'' (possible only for $\ell=3$, by determinants: $\bar\chi_\ell^2=1$ forces $\ell-1\mid2$) and ``$\bar\rho\simeq\bar\rho\otimes\chi_{\ell^*}$'' (a statement about the quadratic field $\Q(\sqrt{\ell^*})$,
\Cref{prop:ant:twistdetect}) are three different conditions. The paper's term $h_\ell(E)$ (Lemma lem:h0) is governed by the second, and the hypothesis of
Theorem thm:main by the third.
\end{warning}

\begin{paperlink}\label{pl:ant:twist}
Section 1.4 and Lemma lem:twist: $E^{(q)}$ is the twist by $\chi_{q^*}$, ``the quadratic character of conductor $q$'', that is by $\Q(\sqrt{q^*})$; it is the
character $\bigl(\frac\cdot q\bigr)$ of \Cref{ex:ant:characters}, ramified only at $q$ (\Cref{prop:ant:ramonlyat}), which is why Lemma lem:twist(i) holds at every
$p\ne q$. Theorems thm:twist and thm:quant: $\chi_S=\prod_{q\in S}\chi_{q^*}$ and $\prod_{q\in S}q$ (the paper's $D$) are the character and the conductor of
$\Q(\sqrt{\prod_{q\in S}q^*})$ (\Cref{prop:ant:ramonlyat}(iii)), and ``the twist is unramified outside $S$'' is \Cref{prop:ant:ramonlyat}(iv). The map
$\iota\colon h\mapsto\sum\chi_S(n)a_n(h)q^n$ uses $\chi_S$ as a Dirichlet character, and $\rho_{E_S}\simeq\rho_E\otimes\chi_S$ uses it as a Galois character; the
two are the same character by \Cref{thm:ant:dirichletgalois}. The table of Section 1 uses $\chi_{-3}(q)$ and $\chi_{-4}(q)$ as values of Dirichlet characters,
equal to the Galois values at $\Frob_q$ (\Cref{prop:ant:quadchar}(iii)), which is how they enter the proof of Lemma lem:euler. Proposition prop:forms uses
$q\equiv\chi_{-3}(q)\pmod3$, which holds for every prime $q\ne3$ by \Cref{prop:ant:quadchar}(iii).
\end{paperlink}

\section[Densities, Chebotarev, complex multiplication]{Densities, the Chebotarev density theorem, complex multiplication}\label{sec:ant:chebotarev}

The paper uses Chebotarev's theorem to compare Galois representations through their Frobenius traces (Lemma lem:old, ``by Chebotarev and
Brauer-Nesbitt''); the original version of the note also used it to count primes inert in a quadratic field (Lemma lem:h0 and Section sec:anomaly in
that version). The opposite direction needs no density theorem: in Section 3 of the paper the optimal-level theorem gives $\bar\rho_g\simeq\bar\rho$, and
$a_p(g)\equiv a_p(f)$ for all $p\nmid3N$ follows by taking traces of Frobenius (\Cref{prop:ant:frobGQ}(iv)). This section states the theorem with the Frobenius of
\Cref{prop:ant:frobenius} and over any number field, proves the corollaries the book needs, gives the criterion for a representation to be isomorphic to its
twist by a quadratic character, and explains, in the form needed to read the original version of Lemma lem:h0 correctly (\Cref{pl:ant:h0}), why $a_p(E)=0$ at the inert primes of a quadratic field
happens only for curves with complex multiplication by that field. The form of Chebotarev over $\Q$ and its Brauer-Nesbitt corollaries are in
\Cref{ch:galois}.

\begin{proposition}[the Dedekind zeta function and densities]\label{prop:ant:dedekindzeta}
Let $K$ be a number field of degree $n$.
\begin{enumerate}[label=(\roman*)]
\item $\zeta_K(s)=\sum_{\mathfrak a}N\mathfrak a^{-s}=\prod_{\mathfrak p}(1-N\mathfrak p^{-s})^{-1}$ (sum over the nonzero ideals, product over the nonzero primes)
converges absolutely for $\mathrm{Re}\,s>1$.
\item (Stated.) $\zeta_K$ extends meromorphically to $\mathrm{Re}\,s>1-1/n$ with a simple pole at $s=1$, whose residue is given by the analytic class number
formula.
\item Given (ii): $\sum_{\mathfrak p}N\mathfrak p^{-s}\sim\log\frac1{s-1}$ as $s\to1^+$. So the \emph{Dirichlet density} of a set $X$ of primes of $K$,
$\delta(X)=\lim_{s\to1^+}\sum_{\mathfrak p\in X}N\mathfrak p^{-s}\big/\log\frac1{s-1}$ when the limit exists, is defined as in \Cref{ch:galois} with $N\mathfrak p$ in
place of $p$. The primes $\mathfrak p$ with $f(\mathfrak p\mid p)\ge2$ have density $0$, so the primes of degree one have density $1$; finite sets have
density $0$.
\end{enumerate}
\end{proposition}

\begin{proof}
(i) By unique factorisation (\Cref{thm:ant:idealfactorisation}) and multiplicativity of the norm (\Cref{prop:ant:idealnorm}), the expansion of the finite
products $\prod_{N\mathfrak p\le x}(1-N\mathfrak p^{-s})^{-1}$ contains each $N\mathfrak a^{-s}$ with $\mathfrak a$ composed of these primes exactly once. For real
$s=\sigma>1$, at most $n$ primes lie above each rational $p$, each with $N\mathfrak p\ge p$ (\Cref{thm:ant:efg}), so
$\prod_{\mathfrak p}(1-N\mathfrak p^{-\sigma})^{-1}\le\prod_p(1-p^{-\sigma})^{-n}=\zeta(\sigma)^n<\infty$; this bounds all partial sums of the series of positive terms.
(ii) We do not prove this; see \cite[Chapter VII, Section 5]{Neukirch} or \cite[Chapter VIII]{LangANT}. (iii) $\log\zeta_K(s)=\sum_{\mathfrak p}\sum_{m\ge1}
N\mathfrak p^{-ms}/m$, and the terms with $m\ge2$ are bounded by $n\sum_p\sum_{m\ge2}p^{-m}\le n\sum_p\frac{2}{p^2}<\infty$ near $s=1$; by (ii)
$\log\zeta_K(s)=\log\frac1{s-1}+O(1)$. The primes with $f\ge2$ have $N\mathfrak p\ge p^2$, and $\sum_{f(\mathfrak p)\ge2}N\mathfrak p^{-s}\le n\sum_pp^{-2}$ stays
bounded.
\end{proof}

\begin{theorem}[Chebotarev density theorem]\label{thm:ant:chebotarev}
Let $L/K$ be a finite Galois extension of number fields with group $G$ and $C\subseteq G$ a conjugacy class. The set of primes $\mathfrak p$ of $K$, unramified in
$L$, with $\Frob_{\mathfrak p}=C$ (\Cref{prop:ant:frobenius}(ii)) has Dirichlet density $\#C/\#G$. It also has natural density $\#C/\#G$ (counting the primes with
$N\mathfrak p\le x$). For $G$ abelian the classes are elements, and the statement contains Dirichlet's theorem on primes in arithmetic progressions.
\end{theorem}

\begin{proof}[Reference]
We do not prove this; see \cite[Chapter VII, Section 13]{Neukirch} (the theorem is numbered 13.4 there; this number and the others in this chapter are
pointers not checked against a printed copy) or \cite[Chapter I, Section 2.2]{SerreAbelian}. The idea: for $G$ cyclic it follows from the nonvanishing
at $s=1$ of the Hecke $L$-functions of $K$ (abelian class field theory), and the general case reduces to the cyclic one by restricting to the cyclic
subgroup generated by an element of $C$. Only the Dirichlet density is used in this book; the natural density is not used in any proof.
\end{proof}

\begin{corollary}\label{cor:ant:chebotarev}
In the setting of \Cref{thm:ant:chebotarev}:
\begin{enumerate}[label=(\roman*)]
\item every $\sigma\in G$ is $\Frob_{\mathfrak P}$ for infinitely many $\mathfrak P$ (lying over a set of primes $\mathfrak p$ of positive density);
\item the primes of $K$ that split completely in $L$ have density $1/[L:K]$; if $M/K$ is a finite extension, not necessarily Galois, the primes of $K$ having
a prime of degree one above them in $M$ have positive density;
\item (Bauer) for $L$, $L'$ Galois over $K$: $L\subseteq L'$ iff every prime of $K$ that splits completely in $L'$, outside a set of density zero, splits
completely in $L$; so a Galois extension of $K$ is determined by its set of split primes;
\item (Dirichlet) for $K=\Q$, $L=\Q(\zeta_n)$: each class $a\in(\Z/n)^\times$ contains a set of primes of density $1/\varphi(n)$;
\item for $L/\Q$ quadratic, the split and the inert primes each have density $\frac12$; in particular the primes inert in $\Q(\sqrt{\ell^*})$ (the $p$ with
$\bigl(\frac p\ell\bigr)=-1$) have density $\frac12$, and for $\ell=3$ these are the $p\equiv2\pmod3$;
\item if $L/\Q$ is Galois, possibly infinite, and unramified outside a finite set, the Frobenius elements at the primes outside any set of density zero are
dense in $\Gal(L/\Q)$ (\Cref{ch:galois}).
\end{enumerate}
\end{corollary}

\begin{proof}
(i) Take the class of $\sigma$; a set of positive density is infinite (finite sets have density $0$), and for each such $\mathfrak p$ some $\mathfrak P\mid\mathfrak p$
has $\Frob_{\mathfrak P}=\sigma$ (\Cref{prop:ant:frobenius}(ii)). (ii) $C=\{1\}$ and \Cref{prop:ant:frobenius}(iii). For $M$, let $L$ be a Galois closure of $M/K$,
$G=\Gal(L/K)$, $H=\Gal(L/M)$; by \Cref{prop:ant:frobsubfield} (whose proof is the same over $K$), $\mathfrak p$ unramified in $L$ has a degree-one prime in $M$ iff
$\langle\Frob_{\mathfrak P}\rangle$ fixes a point of $G/H$, which holds when $\Frob_{\mathfrak P}=1$; so the set contains the split primes, of density $1/[L:K]>0$.
(iii) If $L\subseteq L'$, a prime split completely in $L'$ is split completely in $L$ (\Cref{thm:ant:efg}). Conversely, consider the compositum $LL'$, Galois
over $K$, and the injection $\Gal(LL'/K)\to\Gal(L/K)\times\Gal(L'/K)$ by restriction, which sends $\Frob_{\mathfrak P}$ to the pair of Frobenius elements
(\Cref{prop:ant:frobenius}(v)). So a prime unramified in $LL'$ splits completely in $LL'$ iff it splits completely in $L$ and in $L'$; by hypothesis this is
the set of primes splitting completely in $L'$, up to density zero. By (ii), $1/[LL':K]=1/[L':K]$, so $LL'=L'$ and $L\subseteq L'$. (iv) $\Frob_p=\sigma_p$
corresponds to $p\bmod n$ (\Cref{thm:ant:cyclotomicsplitting}), and $\#G=\varphi(n)$. (v) $G$ has two classes of size $1$; for $\Q(\sqrt{\ell^*})$ use
\Cref{prop:ant:sqrtlstar}(ii). (vi) is proved in \Cref{ch:galois} from \Cref{thm:ant:chebotarev}.
\end{proof}

\begin{proposition}[detecting a quadratic twist]\label{prop:ant:twistdetect}
Let $k$ be a field of characteristic $\ne2$ (a finite field, $\Fbar_\ell$, or a finite extension of $\Q_\ell$), $\rho\colon G_\Q\to\GL_2(k)$ continuous,
semisimple, unramified outside a finite set $S$, and $K$ a quadratic field with character $\chi_K$. The following are equivalent:
(a) $\rho\simeq\rho\otimes\chi_K$; (b) $\tr\rho(\Frob_p)=0$ for every prime $p\notin S$ inert in $K$; (c) $\tr\rho(\Frob_p)=0$ for all primes inert in $K$
outside a set of density zero.
\end{proposition}

\begin{proof}
(a)$\Rightarrow$(b): for $p\notin S$ unramified in $K$ and inert, $\tr\rho(\Frob_p)=\tr(\rho\otimes\chi_K)(\Frob_p)=\chi_K(\Frob_p)\tr\rho(\Frob_p)=-\tr\rho(\Frob_p)$, and
$2$ is invertible in $k$; a prime inert in $K$ is unramified in $K$. (b)$\Rightarrow$(c) is trivial. (c)$\Rightarrow$(a): $\rho\otimes\chi_K$ is semisimple (it has
the same invariant subspaces as $\rho$, $\chi_K$ acting by scalars) and unramified outside $S$ and the primes ramified in $K$. For $p$ outside this finite set,
$\det(X-(\rho\otimes\chi_K)(\Frob_p))=X^2-\chi_K(\Frob_p)\tr\rho(\Frob_p)X+\det\rho(\Frob_p)$, which equals $\det(X-\rho(\Frob_p))$ when $p$ splits, and when $p$ is
inert and $\tr\rho(\Frob_p)=0$. By (c) the characteristic polynomials agree at all primes outside a set of density zero, so $\rho\simeq\rho\otimes\chi_K$ by
the Chebotarev and Brauer-Nesbitt theorem of \Cref{ch:galois} (both being semisimple).
\end{proof}

Two applications. (a) If $\bar\rho_{E,3}$ is irreducible, then $\bar\rho_{E,3}\simeq\bar\rho_{E,3}\otimes\chi_{-3}$ iff $3\mid a_p(E)$ for all primes $p\equiv2\pmod 3$ of good
reduction (outside a density zero set), since $\tr\bar\rho_{E,3}(\Frob_p)=a_p\bmod3$ and the inert primes of $\Q(\sqrt{-3})$ are the $p\equiv2\pmod3$ (\Cref{ch:galois}
gives the example $y^2=x^3-6x-8$). (b) If $\rho_{E,\ell}$ is semisimple (it always is, by Faltings's theorem, stated in \Cref{ch:galrep}; and it is irreducible
when $\bar\rho_{E,\ell}$ is), then $\rho_{E,\ell}\simeq\rho_{E,\ell}\otimes\chi_K$ iff $a_p(E)=0$ for almost all $p$ inert in $K$ (outside a density zero set).
The criterion concerns twisting by a quadratic character; it says nothing about $\rho\otimes\chi_\ell$, which is never isomorphic to $\rho$
(\Cref{warn:ant:chil}).

\begin{theorem}[complex multiplication]\label{thm:ant:cm}
Let $E/\Q$ be an elliptic curve with complex multiplication, that is $\End_{\Qbar}(E)\ne\Z$. Then $\End_{\Qbar}(E)$ is an order $\mathcal O$ in an imaginary
quadratic field $K$, the endomorphisms are defined over $K$, and $\sigma\in G_\Q$ acts on them through $\Gal(K/\Q)$, by complex conjugation of
$\mathcal O\subset K$ when $\chi_K(\sigma)=-1$ (\Cref{ch:ellcurves}).
\begin{enumerate}[label=(\roman*)]
\item (Stated.) $\mathcal O$ has class number one; so $K=\Q(\sqrt{-d})$ with $d\in\{1,2,3,7,11,19,43,67,163\}$, and $j(E)$ is one of $13$ rational values.
\item For every prime $\ell$, $\rho_{E,\ell}\simeq\rho_{E,\ell}\otimes\chi_K$, and $a_p(E)=0$ for every prime $p$ of good reduction inert in $K$.
\item Every prime ramified in $K$ is a prime of bad reduction of $E$.
\item (Stated.) Conversely, if $a_p(E)=0$ for all primes $p$ inert in a quadratic field $K$ outside a set of density zero, then $E$ has complex multiplication by
$K$; in particular $K$ is imaginary.
\item (Stated.) If $E$ has no complex multiplication, the set of primes $p$ with $a_p(E)=0$ has density $0$.
\item Consequently, for an odd prime $\ell$, the condition ``$a_p(E)=0$ for almost all $p$ inert in $\Q(\sqrt{\ell^*})$'' holds if and only if $E$ has
complex multiplication by $\Q(\sqrt{\ell^*})$; this needs $\ell\equiv3\pmod4$ and $\ell\in\{3,7,11,19,43,67,163\}$, and then $\ell$ divides the conductor of $E$ by (iii). It never
happens when $\ell\nmid N$.
\end{enumerate}
\end{theorem}

\begin{proof}[Sketch of proof]
The structure of $\End_{\Qbar}(E)$ and the Galois action, $\rho_{E,\ell}(\sigma)\alpha\rho_{E,\ell}(\sigma)^{-1}=\sigma(\alpha)$ on $V_\ell E$ for $\alpha\in\End_{\Qbar}(E)$,
are quoted from \Cref{ch:ellcurves} and \cite[Chapter II, Section 2]{SilvermanATAEC}. (ii) Choose $\alpha\in\mathcal O$ with $\bar\alpha=-\alpha$, $\alpha\ne0$ (for
instance $\alpha=f\sqrt{-d}$ for a suitable integer $f$); it acts invertibly on $V_\ell E$, and $\rho(\sigma)\alpha\rho(\sigma)^{-1}=\sigma(\alpha)=\chi_K(\sigma)\alpha$. So
$\alpha\rho(\sigma)\alpha^{-1}=\chi_K(\sigma)\rho(\sigma)$: $\alpha$ is an isomorphism $\rho\to\rho\otimes\chi_K$. For $p\ne\ell$ of good reduction and inert in $K$,
$\chi_K(\Frob_p)=-1$ gives $a_p=\tr\rho(\Frob_p)=-a_p$, so $a_p=0$ (choose $\ell\ne p$, and $a_p\in\Z$). (iii) If $p$ is ramified in $K$ and $E$ had good reduction
at $p$, take $\ell\ne p$: $\rho$ is unramified at $p$, and there is $\tau\in I_p$ with $\chi_K(\tau)=-1$ (\Cref{prop:ant:quadchar}(ii)); then
$\alpha=\rho(\tau)\alpha\rho(\tau)^{-1}=\chi_K(\tau)\alpha=-\alpha$, a contradiction. (i) is the theorem of Heegner, Baker and Stark on imaginary quadratic fields of
class number one, together with the fact that $j(E)\in\Q$ generates the ring class field of $\mathcal O$; see \cite[Chapter II and Appendix A]{SilvermanATAEC} and
\cite[Sections 7, 11 and 12]{Cox}. (iv) We only indicate the argument. Let $L$ be the compositum of $K$ and the field cut out by $\rho_{E,\ell}$, unramified outside $\ell Nd_K$. By
\Cref{cor:ant:chebotarev}(vi) the Frobenius elements of $\Gal(L/\Q)$ at the primes outside the exceptional set of density zero are dense, so those at the inert
such primes are dense in the open coset $\Gal(L/\Q)\setminus\Gal(L/K)$; the trace of $\rho_{E,\ell}$ is continuous and vanishes at them, so it vanishes on
$G_\Q\setminus G_K$. If $E$ had no complex multiplication, the image of $\rho_{E,\ell}$ would be open in $\GL_2(\Z_\ell)$ (Serre's open image theorem
\cite{Serre72}), and the trace cannot vanish on a coset of an open subgroup of index $2$ of an open subgroup of $\GL_2(\Z_\ell)$ (such a coset contains elements
$g$ and $(1+\ell^kx)g$ for all $x\in\Mat_2(\Z_\ell)$ and $k$ large, and $\tr(xg)$ is not identically zero). So $E$ has complex multiplication by some $K'$. If
$K'\ne K$, a set of primes of positive density splits in $K'$ and is inert in $K$ (\Cref{cor:ant:chebotarev}(i) in $KK'$); at such a prime of good reduction
the reduction is ordinary (Deuring's criterion, \cite[Section 14]{Cox}), so $a_p\not\equiv0\pmod p$ and $a_p\ne0$, a contradiction. (v) is due to Serre \cite[Section 7]{Serre81}. (vi) The forward implication is
(iv); the converse is (ii). The class number one list and $\ell^*<0$ (needed for $K$ imaginary, \Cref{prop:ant:sqrtlstar}(iv)) give the list of $\ell$, and
$\ell$ is ramified in $\Q(\sqrt{\ell^*})$.
\end{proof}

\begin{example}[curves with complex multiplication]\label{ex:ant:cmcurves}
(Checked with PARI: $a_p=0$ at every good prime $p$ inert in the field among the first $70$ primes, and $a_p\ne0$ at the split primes listed.)
$y^2=x^3-x$, conductor $32$, CM by $\Z[i]$: $a_p=0$ for $p\equiv3\pmod 4$, and $a_5=-2$, $a_{13}=6$, $a_{17}=2$, $a_{29}=-10$, $a_{37}=-2$.
$y^2=x^3+1$ (36a1) and $y^2+y=x^3-7$ (27a1), $j=0$, CM by $\Z[\zeta_3]$: $a_p=0$ for $p\equiv2\pmod3$; for 36a1, $a_{13}=2$, $a_{37}=-10$.
121b1 $=[0,-1,1,-7,10]$, $j=-2^{15}$, CM by the maximal order of $\Q(\sqrt{-11})=\Q(\sqrt{11^*})$: $a_p=0$ when $\bigl(\frac p{11}\bigr)=-1$, for example $a_2=a_7=0$,
while $a_5=-3$, $a_{37}=7$. In each case the primes ramified in $K$ ($2$, $3$, $3$, $11$) divide the conductor, as \Cref{thm:ant:cm}(iii) requires.
\end{example}

\begin{paperlink}\label{pl:ant:h0}
Lemma lem:h0 (revised version): $H^0(\Q,W^D)$ is cyclic of length $h_\ell(E)$, $h_\ell(E)=0$ for $\ell\ge5$ and $h_3(E)\le1$, ``for every $E$, with or without complex
multiplication''. The proof is the determinant argument (\Cref{ch:reps,ch:galois}): an isomorphism $\rho\bmod\ell^k\simeq(\rho\otimes\chi_\ell)\bmod\ell^k$ forces
$\chi_\ell^2\equiv1\pmod{\ell^k}$; it does not use complex multiplication or the distribution of inert primes. The original version of the note derived the finiteness
from ``$\rho\simeq\rho\otimes\chi_\ell$ would force $a_p=0$ at all primes inert in $\Q(\sqrt{\ell^*})$'', which confuses the cyclotomic character $\chi_\ell$ with the
quadratic character $\chi_{\ell^*}$ (\Cref{warn:ant:chil}); the correct statements about $a_p=0$ at inert primes are \Cref{prop:ant:twistdetect} and
\Cref{thm:ant:cm}, and by \Cref{thm:ant:cm}(vi) the complex multiplication alternative cannot even occur when $\ell\nmid N$. The equivalence ``$h_3(E)=1$ iff
$\bar\rho_{E,3}\simeq\bar\rho_{E,3}\otimes\chi_{-3}$'' uses $\bar\chi_3^{\mathrm{cyc}}=\chi_{-3}$ (\Cref{ex:ant:characters}), and its link with the reducibility of $\bar\rho|_{G_{\Q(\sqrt{-3})}}$, like
Proposition prop:ns, obtains $\bar\rho\simeq\bar\rho\otimes\chi_{-3}$ from $\bar\rho\otimes\Fbar_3\simeq\Ind\psi$ for a character $\psi$ of
$\Gal(\Qbar/\Q(\sqrt{-3}))$, by Frobenius reciprocity and $\Ind\psi\otimes\chi_{-3}\simeq\Ind\psi$ (\Cref{ch:reps}), not from Frobenius traces. Application (a) of
\Cref{prop:ant:twistdetect} gives the equivalent trace criterion ($3\mid a_p$ for almost all $p\equiv2\pmod3$), used in the example of \Cref{ch:galois}, not in the
paper. Section sec:anomaly no longer needs the original numerical argument for $h_3(E)\le1$ (that some prime $p<1223$
inert in $\Q(\sqrt{-3})$ has $9\nmid a_p$); $h_3(E)\le1$ now follows from Lemma lem:h0.
\end{paperlink}

\section{Coefficient fields and primes above \texorpdfstring{$\ell$}{l}}\label{sec:ant:coefficients}

Lemmas lem:crit and lem:old of the paper work with eigenforms whose coefficients lie in ``the ring of integers $\mathcal O$ of a number field'', with ``a prime
$\lambda\mid\ell$ of $\mathcal O$'', with congruences modulo $\lambda$, with ``a maximal ideal of residue characteristic $\ell$'' of the Hecke algebra, and with the
operation ``after enlarging $\mathcal O$ and $\lambda$''. This section makes each of these precise. The Hecke algebra $\T$ itself, and the facts that its
eigencharacters correspond to eigenforms, are in \Cref{ch:modforms,ch:hecke}; here only the algebra of a commutative ring that is finitely generated as a
$\Z$-module is used.

\begin{proposition}[coefficient fields]\label{prop:ant:coefffield}
Let $R$ be a commutative ring, finitely generated as a $\Z$-module, and $\phi\colon R\to\C$ a ring homomorphism. Then every $\phi(t)$ is an algebraic integer,
$K_\phi=\Q(\phi(R))$ is a number field, and $\phi(R)$ is a subring of $\mathcal O_{K_\phi}$ of finite index (an \emph{order}). For
$\sigma\in G_\Q$, $\sigma\circ\phi$ is again such a homomorphism (a \emph{Galois conjugate}). For the eigencharacter $\phi_f\colon\T\to\C$, $T_n\mapsto a_n(f)$, of a
newform $f$ of weight $2$ on $\Gamma_0(N)$: the \emph{coefficient field} $K_f=\Q(a_n(f):n\ge1)$ is a number field, every $a_n(f)$ lies in $\mathcal O_{K_f}$, and
(stated; \Cref{ch:hecke}) the conjugates $f^\sigma=\sum\sigma(a_n)q^n$ are newforms and $K_f$ is totally real. For $f=f_E$, $K_f=\Q$, $\mathcal O=\Z$, and the only
prime above $\ell$ is $\lambda=\ell$.
\end{proposition}

\begin{proof}
$\phi(R)$ is a subring of $\C$ that is finitely generated as a $\Z$-module, so each of its elements is integral over $\Z$ (\Cref{prop:ant:integral}(iii),
whose proof works in any field). The field $K_\phi$ is generated by the images of finitely many $\Z$-module generators, each algebraic, so it is a number
field, and $\phi(R)\subseteq\mathcal O_{K_\phi}$. The $\Q$-span of $\phi(R)$ is a subring of $\C$ of finite dimension over $\Q$, a domain, hence a field, so it is $K_\phi$; thus
$\phi(R)$ and $\mathcal O_{K_\phi}$ are free abelian groups of the same rank $[K_\phi:\Q]$ (\Cref{thm:ant:integralbasis}), and the index is finite. For $\sigma\in G_\Q$, $\sigma\circ\phi$ is a ring homomorphism. For $\phi_f$, $\T$ is finitely generated as a $\Z$-module
(\Cref{ch:hecke}), $\phi_f(\T)=\Z[a_n(f):n\ge1]$, which spans $K_f$, and $a_1(f)=1$.
\end{proof}

\begin{proposition}[primes above $\ell$ and reduction modulo $\lambda$]\label{prop:ant:primesabove}
Let $K$ be a number field, $\mathcal O=\mathcal O_K$ and $\ell$ a prime.
\begin{enumerate}[label=(\roman*)]
\item The primes $\lambda\mid\ell$ of $\mathcal O$ are the maximal ideals containing $\ell$; there are at most $[K:\Q]$ of them; $k_\lambda=\mathcal O/\lambda\simeq\F_{\ell^f}$
with $f=f(\lambda\mid\ell)$; reduction modulo $\lambda$ is the ring homomorphism $\mathcal O\to k_\lambda$, $a\mapsto\bar a$.
\item For $a,b\in\Z$: $a\equiv b\pmod\lambda$ iff $a\equiv b\pmod\ell$. For $a\in\mathcal O$: $\ell\mid N_{K/\Q}(a)$ iff some $\lambda\mid\ell$ contains $a$; and $a$ lies in
every $\lambda\mid\ell$ iff $a$ lies in the radical of $\ell\mathcal O$, iff $a^m\in\ell\mathcal O$ for some $m\ge1$.
\item (Enlarging.) If $K\subseteq K'$ then $\lambda\mathcal O_{K'}\ne\mathcal O_{K'}$, every prime $\lambda'$ of $\mathcal O_{K'}$ dividing $\lambda\mathcal O_{K'}$ has
$\lambda'\cap\mathcal O=\lambda$, and reduction modulo $\lambda'$ restricts on $\mathcal O$ to reduction modulo $\lambda$ followed by $k_\lambda\subseteq k_{\lambda'}$.
\item (Lying over, for orders.) Let $\alpha$ be integral over $\mathcal O$ and $\mathcal O'=\mathcal O[\alpha]$. For every prime $\lambda$ of $\mathcal O$ there is a maximal
ideal $\mathfrak q$ of $\mathcal O'$ with $\mathfrak q\cap\mathcal O=\lambda$, namely $\mathfrak q=\lambda'\cap\mathcal O'$ for a prime $\lambda'$ above $\lambda$ in the ring of integers of
$K(\alpha)$; its residue field is a finite field containing $k_\lambda$.
\item (The step of Lemma lem:old.) Let $p\ne\ell$ be a prime, $a\in\mathcal O$, and $\alpha_p,\beta_p$ the roots of $X^2-aX+p$; let $\lambda'$ be a prime of the ring of
integers of $K(\alpha_p)$ above $\lambda$. If $X^2-aX+p\equiv(X-u)(X-w)\pmod\lambda$ with $u,w\in\mathcal O$, then $\alpha_p\equiv u$ or $\alpha_p\equiv w\pmod{\lambda'}$,
and $\beta_p$ is congruent to the other one. In particular, if $u=sp$ and $w=s$ with $s=\pm1$, one of $\alpha_p,\beta_p$ is congruent to $s$ modulo $\lambda'$.
\end{enumerate}
\end{proposition}

\begin{proof}
(i) A prime containing $\ell$ is nonzero, hence maximal; it lies above $\ell$ (\Cref{thm:ant:efg}), and there are at most $\sum e_if_i=[K:\Q]$ of them.
$k_\lambda$ is a finite field of characteristic $\ell$ (\Cref{prop:ant:idealnorm}(iii)). (ii) $\lambda\cap\Z=\ell\Z$. $N_{K/\Q}(a\mathcal O)=\prod_{\mathfrak p}N\mathfrak p^{v_{\mathfrak p}(a)}$
(\Cref{prop:ant:idealnorm}), divisible by $\ell$ iff some $\mathfrak p\mid\ell$ has $v_{\mathfrak p}(a)>0$. The intersection of the $\lambda\mid\ell$ is
$\prod_{\lambda\mid\ell}\lambda$ (\Cref{prop:ant:crt}(i)), and $a$ lies in it iff $v_\lambda(a)\ge1$ for all $\lambda\mid\ell$ iff some power of $a$ lies in
$\ell\mathcal O=\prod\lambda^{e_\lambda}$. (iii) \Cref{thm:ant:efg} for $K'/K$ and $\mathfrak p=\lambda$. (iv) Let $K''=K(\alpha)$ and $\lambda'$ a prime of $\mathcal O_{K''}$
above $\lambda$, which exists by (iii). Then $\mathfrak q=\lambda'\cap\mathcal O'$ is a prime ideal of $\mathcal O'\subseteq\mathcal O_{K''}$ with $\mathfrak q\cap\mathcal O=
\lambda$, and $\mathcal O'/\mathfrak q$ embeds in the finite field $\mathcal O_{K''}/\lambda'$, so it is a finite domain, hence a field. (v) $\alpha_p,\beta_p$ are
integral over $\mathcal O$, so they lie in the ring of integers of $K(\alpha_p)$, whose quotient by $\lambda'$ is a field $k$. In $k[X]$,
$(X-\bar\alpha_p)(X-\bar\beta_p)=X^2-\bar aX+\bar p=(X-\bar u)(X-\bar w)$, and unique factorisation in $k[X]$ gives $\{\bar\alpha_p,\bar\beta_p\}=\{\bar u,\bar w\}$ as
multisets.
\end{proof}

\begin{proposition}[maximal ideals from congruences]\label{prop:ant:congruencemaximal}
Let $R$ be a commutative ring, finitely generated as a $\Z$-module, $\mathcal O=\mathcal O_K$, $\lambda\mid\ell$ a prime of $\mathcal O$, and $\phi\colon R\to\mathcal O$ a
ring homomorphism.
\begin{enumerate}[label=(\roman*)]
\item $\mathfrak m=\ker(R\xrightarrow{\phi}\mathcal O\to k_\lambda)$ is a maximal ideal of $R$, and $R/\mathfrak m$ is a subfield of $k_\lambda$, of characteristic $\ell$.
\item If $\phi'\colon R\to\mathcal O$ is another ring homomorphism with $\phi'(t)\equiv\phi(t)\pmod\lambda$ for all $t$ in a set of ring generators of $R$ (for $R=\T$:
$a_n(g)\equiv a_n(f)$ for all $n$), then $\phi$ and $\phi'$ define the same $\mathfrak m$.
\item If $J\subseteq R$ is an ideal contained in $\mathfrak m$ such that $R/J$ is a cyclic group of order $r$, then $\ell\mid r$ (with $\ell\mid0$).
\item (Stated; \Cref{ch:modules}.) $R\otimes\Z_\ell$ is the product of its localisations at its maximal ideals, which are the maximal ideals of $R$ containing $\ell$;
and $\phi$ extends to $R\otimes\Z_\ell\to\mathcal O_\lambda$ (\Cref{prop:ant:otimesZl}), which factors through the factor $(R\otimes\Z_\ell)_{\mathfrak m}$.
\end{enumerate}
\end{proposition}

\begin{proof}
(i) $R/\mathfrak m$ embeds in the finite field $k_\lambda$; a subring of a finite field is a finite domain, hence a field. (ii) The two compositions
$R\to k_\lambda$ are ring homomorphisms agreeing on generators, hence equal, so they have the same kernel. (iii) $R/\mathfrak m$ is a nonzero quotient of the
cyclic group $R/J\simeq\Z/r$, so it is $\Z/s$ with $s\mid r$; it has characteristic $\ell$, so $s=\ell$. (iv) We do not prove this here; the decomposition
is the idempotent-lifting theorem of \Cref{ch:modules}, and the extension of $\phi$ is $\phi\otimes\Z_\ell$ followed by the projection of
$\mathcal O\otimes\Z_\ell$ onto its factor $\mathcal O_\lambda$ (\Cref{prop:ant:otimesZl}), which kills the idempotents of the other factors of $R\otimes\Z_\ell$.
\end{proof}

\begin{proposition}[$\mathcal O_K\otimes\Z_\ell$]\label{prop:ant:otimesZl}
For a number field $K$ and a prime $\ell$,
\[
\mathcal O_K\otimes_\Z\Z_\ell\simeq\prod_{\lambda\mid\ell}\mathcal O_{K,\lambda},\qquad K\otimes_\Q\Q_\ell\simeq\prod_{\lambda\mid\ell}K_\lambda,
\]
where $\mathcal O_{K,\lambda}=\varprojlim_k\mathcal O_K/\lambda^k$ is the $\lambda$-adic completion, a complete discrete valuation ring, free over $\Z_\ell$ of rank $e_\lambda f_\lambda$
($\sum_\lambda e_\lambda f_\lambda=[K:\Q]$), with uniformiser any $\varpi\in\lambda\setminus\lambda^2$ and residue field $k_\lambda$; and $K_\lambda$, its fraction field, is
an extension of $\Q_\ell$ of degree $e_\lambda f_\lambda$. The cases $K=\Q(\zeta_3)$ and $\Q(i)$ are \Cref{prop:ant:zlzeta3}.
\end{proposition}

\begin{proof}[Reference]
We do not prove this here; it is proved in \Cref{ch:local}, and see \cite[Chapter II, Section 8]{Neukirch} or \cite[Chapter II]{SerreLocalFields}. The
algebraic core is the Chinese remainder theorem: $\mathcal O_K/\ell^k\mathcal O_K\simeq\prod_{\lambda\mid\ell}\mathcal O_K/\lambda^{ke_\lambda}$ for every $k$ (\Cref{prop:ant:crt}), and
$\mathcal O_K\otimes\Z_\ell=\varprojlim_k\mathcal O_K/\ell^k\mathcal O_K$ because $\mathcal O_K$ is free of finite rank over $\Z$.
\end{proof}

\begin{example}[level $23$]\label{ex:ant:level23}
The space $S_2(\Gamma_0(23))$ is two-dimensional and spanned by a newform $f$ and its Galois conjugate (stated; PARI's \texttt{mfinit([23,2],0)}, and
LMFDB 23.2.a.a). Its coefficient field is $K_f=\Q(\sqrt5)$, and $a_2(f)$ is a root $\theta$ of $X^2+X-1$, the characteristic polynomial of $T_2$ (computed with
PARI's \texttt{mfheckemat}); the first coefficients are
\[
f=q+\theta q^2+(-2\theta-1)q^3+(-\theta-1)q^4+2\theta q^5+(\theta-2)q^6+(2\theta+2)q^7+\cdots,\qquad\theta^2+\theta-1=0,
\]
and $a_{11}=-2\theta-4$, $a_{13}=3$ (PARI's \texttt{mfeigenbasis}). The ring of integers is $\mathcal O_{K_f}=\Z[\frac{1+\sqrt5}2]=\Z[\theta]$, as $\disc(X^2+X-1)=5$ is
squarefree. Primes: $5=(\sqrt5)^2$ is ramified, with $\theta\equiv2\pmod{\sqrt5}$ since $X^2+X-1\equiv(X-2)^2\pmod5$; $11$ splits ($\bigl(\frac5{11}\bigr)=1$) into
$\lambda=(11,\theta-3)$ and $\lambda'=(11,\theta-7)$, as $X^2+X-1\equiv(X-3)(X-7)\pmod{11}$, both with residue field $\F_{11}$; $2$ is inert ($5\equiv5\pmod8$), with
residue field $\F_4$, in which $\theta$ reduces to a root of $X^2+X+1$, a primitive cube root of unity. So the reductions of $f$ modulo $\lambda$ and
$\lambda'$ are different: modulo $\lambda$, $(a_2,a_3,a_5,a_7)\equiv(3,4,6,8)$, and modulo $\lambda'$, $(7,7,3,5)$; they are exchanged by $\Gal(K_f/\Q)$. A congruence
``$a_n(g)\equiv a_n(f)\pmod\lambda$'' as in Lemma lem:crit depends on $\lambda$, not only on $\ell$. The Hecke algebra of level $23$ is $\T=\Z[T_2]\simeq\Z[\theta]$.
Indeed, let $\T$ be the subring of $\End(S_2(\Gamma_0(23)))$ generated by all the $T_n$, and $\phi_f\colon\T\to\mathcal O_{K_f}$ the eigencharacter of $f$
(\Cref{prop:ant:coefffield}). It is injective: if $\phi_f(t)=0$, then also $\phi_{f^\sigma}(t)=\sigma(\phi_f(t))=0$ for the conjugate $f^\sigma$, so $t$ kills $f$ and
$f^\sigma$, which span $S_2(\Gamma_0(23))$, and $t=0$. Its image is a subring of $\mathcal O_{K_f}=\Z[\theta]$ containing $\phi_f(T_2)=a_2(f)=\theta$, so it is
$\Z[\theta]$, and $\T=\Z[T_2]$ since $\phi_f(T_2)$ generates the image. (Checked with PARI: \texttt{mfdim([23,2],1)} is $2$ and \texttt{nfbasis(x\^{}2+x-1)} is
$[1,x]$, so $\Z[\theta]$ is the full ring of integers of $\Q(\sqrt5)$.)
\end{example}

\begin{paperlink}\label{pl:ant:crit}
Lemma lem:crit: ``an eigenform $g$ ... with coefficients in the ring of integers $\mathcal O$ of a number field, and a prime $\lambda\mid\ell$ of $\mathcal O$ such that
$a_n(g)\equiv a_n(f)\pmod\lambda$ for every $n\ge1$''; then ``$I+J$ lies in the kernel of $\T\to\mathcal O/\lambda$, a maximal ideal of residue characteristic $\ell$, and
$\ell$ divides $\#\T/(I+J)=r$''. The ring $\mathcal O$ exists by \Cref{prop:ant:coefffield} (the eigencharacter $\phi_g$ of $g$ takes values in a number field and in its
ring of integers); $\phi_g\equiv\phi_f\pmod\lambda$ on the generators $T_n$ gives the same kernel $\mathfrak m$ (\Cref{prop:ant:congruencemaximal}(ii)), so
$I=\ker\phi_f\subseteq\mathfrak m$ and $J\subseteq\ker\phi_g\subseteq\mathfrak m$; and $\ell\mid r$ is \Cref{prop:ant:congruencemaximal}(iii) applied to $I+J$. ``The same
statements hold with $\Z$ replaced by $\Z_\ell$'' uses \Cref{prop:ant:congruencemaximal}(iv). Lemma lem:old: ``after enlarging $\mathcal O$ and $\lambda$'' is
\Cref{prop:ant:primesabove}(iii): the eigenform $g$ built from $h$ has coefficients $\mu_p\in\{\alpha_p,\beta_p\}$ in a larger field, and a prime above $\lambda$ there
restricts to $\lambda$; ``one of $\alpha_p,\beta_p$ is congruent to $a_p(f)$ modulo a prime above $\lambda$ in $\mathcal O[\alpha_p]$'' is
\Cref{prop:ant:primesabove}(iv) and (v), with $X^2-a_p(h)X+p\equiv(X-a_p(f)p)(X-a_p(f))\pmod\lambda$ and $a_p(f)=\pm1$. In Section 3 of the paper all forms have
rational coefficients and $\lambda=\ell$; Lemma lem:old is needed for general $h$, whose coefficient field can be larger, as in \Cref{ex:ant:level23}.
\end{paperlink}

\section{Exercises}\label{sec:ant:exercises}

\begin{exercise}\label{ex:ant:ex-sqrt13}
Let $K=\Q(\sqrt{13})$. Show that $\mathcal O_K=\Z[\theta_{13}]$ with $\theta_{13}=\frac{1+\sqrt{13}}2$, that $d_K=13$, and that $\frac{3+\sqrt{13}}2$ is a unit of norm $-1$. Factor
$2$, $3$ and $13$ in $\mathcal O_K$ (answer: $2$ inert, $3$ split, $13$ ramified). Show that Dedekind-Kummer applied to the generator $\sqrt{13}$ (that is, to $\Z[\sqrt{13}]$) at $p=2$ would predict
that $2$ ramifies, and explain why the theorem does not apply there.
\end{exercise}

\begin{exercise}\label{ex:ant:ex-sqrtminus3}
In the non-integrally closed ring $A=\Z[\sqrt{-3}]$ let $\mathfrak a=(2,1+\sqrt{-3})$. Show that $\mathfrak a^2=2\mathfrak a$ but $\mathfrak a\ne(2)$, so ideals cannot be
cancelled in $A$; deduce that $A$ is not a Dedekind domain and identify the hypothesis of \Cref{def:ant:dedekind} that fails. What is $\mathfrak a\Z[\zeta_3]$?
\end{exercise}

\begin{exercise}\label{ex:ant:ex-factor}
Factor $65$, $91$ and $1729=7\cdot13\cdot19$ into primes in $\Z[i]$ and in $\Z[\zeta_3]$. Deduce the number of ways of writing $65$ as $a^2+b^2$ with $0<a<b$, and
explain the answer by counting the elements of norm $65$ in $\Z[i]$.
\end{exercise}

\begin{exercise}\label{ex:ant:ex-sqrtminus5}
Show that a prime $p\ne2,5$ splits in $\Q(\sqrt{-5})$ iff $p\equiv1,3,7,9\pmod{20}$. Show that the primes above $3$ and $7$ are not principal, that the
primes above $3$ lie in the class of $\mathfrak p_2$ of \Cref{ex:ant:sqrtminus5}, and, using \Cref{thm:ant:classgroup}, that $h(\Q(\sqrt{-5}))=2$. Which primes $p$
are of the form $a^2+5b^2$?
\end{exercise}

\begin{exercise}\label{ex:ant:ex-cubic}
Let $L=\Q(2^{1/3},\zeta_3)$, with $\Gal(L/\Q)\simeq S_3$ acting on the three cube roots of $2$. For $p\ge5$, show using \Cref{cor:ant:cycletype}: if $p\equiv2\pmod3$,
$\Frob_p$ is a transposition and $X^3-2$ has exactly one root modulo $p$; if $p\equiv1\pmod3$, $\Frob_p$ is trivial or a $3$-cycle according as $2$ is or is not a
cube modulo $p$. Determine $\Frob_p$ for $p=5,7,31$, and deduce from \Cref{thm:ant:chebotarev} that $2$ is a cube modulo $p$ for a set of primes of density $\frac23$.
\end{exercise}

\begin{exercise}\label{ex:ant:ex-ramify}
Without using \Cref{thm:ant:discram}, show that an odd prime $p$ ramifies in $\Q(\sqrt d)$ ($d$ squarefree) iff $p\mid d$: compute $(p,\sqrt d)^2$ when $p\mid d$,
and when $p\nmid d$ show that $\mathcal O_K/p\mathcal O_K$ is reduced.
\end{exercise}

\begin{exercise}\label{ex:ant:ex-zeta7}
Factor $2$, $13$ and $29$ in $\Z[\zeta_7]$ and check $efg=6$ against \Cref{thm:ant:cyclotomicsplitting}. The quadratic subfield of $\Q(\zeta_7)$ is $\Q(\sqrt{-7})$; check that
the splitting of each of these primes in it agrees with $\bigl(\frac p7\bigr)$ and with the image of $\Frob_p=\sigma_p$ in $(\Z/7)^\times/(\Z/7)^{\times2}$.
\end{exercise}

\begin{exercise}\label{ex:ant:ex-zeta8}
Determine the decomposition and inertia groups in $\Gal(\Q(\zeta_8)/\Q)\simeq(\Z/8)^\times$ of $2$, $3$, $5$, $7$ and $17$. Identify the three quadratic subfields
$\Q(i)$, $\Q(\sqrt2)$, $\Q(\sqrt{-2})$ with the subgroups $\{1,5\}$, $\{1,7\}$, $\{1,3\}$, and deduce that $\bigl(\frac{-2}p\bigr)=1$ iff $p\equiv1,3\pmod8$.
\end{exercise}

\begin{exercise}\label{ex:ant:ex-unramified}
(a) List the quadratic fields unramified outside $\{3,7\}$ (answer: $\Q(\sqrt{-3})$, $\Q(\sqrt{-7})$, $\Q(\sqrt{21})$). Which of them is cut out by $\chi_S$ for $S=\{3,7\}$,
and what is its conductor? (b) List the seven quadratic fields unramified outside $\{2,3\}$, with their discriminants.
\end{exercise}

\begin{exercise}\label{ex:ant:ex-valuations}
In $\Z_3[\zeta_3]$ compute $v_\pi$ of $1+\zeta_3$, $2+\zeta_3$, $3$, $\sqrt{-3}$, $1-\zeta_3^2$ and $4+\zeta_3$ (answers $0,1,2,1,1,0$). Show that the unit group of
$\Z_3[\zeta_3]/(\pi^2)$ is cyclic of order $6$, generated by the class of $-\zeta_3$, and that every unit of $\Z_3[\zeta_3]$ is $\pm\zeta_3^k$ times a unit congruent to $1$
modulo $\pi^2$.
\end{exercise}

\begin{exercise}\label{ex:ant:ex-sigma}
Let $\sigma=\sm1{-3}1{-2}\in\GL_2(\Z_3)$, acting on $T=\Z_3^2$. Check that $\sigma$ has order $3$ and determinant $1$, that its reduction modulo $3$ is a nontrivial
unipotent, and that $e_1$ generates $T$ over $\Z_3[\sigma]$. Show more generally that $v\in T$ generates $T$ over $\Z_3[\sigma]$ iff $v\notin(1-\sigma)T$, and find
$(1-\sigma)T$ and the fixed line of $\sigma$ on $T/3T$ explicitly.
\end{exercise}

\begin{exercise}\label{ex:ant:ex-z7}
Describe $\Z_\ell[\zeta_3]$ for $\ell=2,5,7$. For $\ell=7$, find the roots of $X^2+X+1$ modulo $7$ and modulo $49$ (answers $2,4$ and $30,18$), and show that the
idempotents of $\Z_7[\zeta_3]\simeq\Z_7\times\Z_7$ are congruent modulo $49$ to $4\zeta_3+27$ and $45\zeta_3+23$.
\end{exercise}

\begin{exercise}\label{ex:ant:ex-euler}
(a) Let $\sigma$ have characteristic polynomial $X^2+X+1$ and let $g=-\sigma$ (order $6$, types $\kod{II}$, $\kod{II}^*$). Show that the commutant of $g$ and its
trace-zero line are those of $\sigma$, and that $AgA^{-1}=g^q$ for an odd prime $q\ne3$ gives the same sign $\chi_{-3}(q)$ as in \Cref{prop:ant:commutant}(iii).
(b) Using the table of Section 1 of the paper, compute $e_\ell(E,q)$ for $q\in\{5,7,11,13\}$, every odd $\ell\ne q$, and the types $\kod{IV}$ and $\kod{III}$. Show that
$e_3(E,q)\ge1$ for every $q\ge5$ of type $\kod{II}$, $\kod{II}^*$, $\kod{IV}$, $\kod{IV}^*$, and decide for which $q$ the same holds for type $\kod{III}$.
\end{exercise}

\begin{exercise}\label{ex:ant:ex-signs}
Let $E$ have nonsplit multiplicative reduction at $q\ne\ell$, with $\rho(\Frob)=\sm{-q}{c'}0{-1}$ for an arithmetic Frobenius. Show that the arithmetic Frobenius
acts on $(\ad V)^{I_q}$ by $q$ and on $(\ad V(1))^{I_q}$ by $q^2$, and the geometric one by $q^{-1}$ and $q^{-2}$. Deduce that
$\det(1-\Frob_q^{-1}\mid(\ad V(1))^{I_q})=1-q^{-2}=L_q(\ad V,1)^{-1}$, whereas the arithmetic Frobenius gives $1-q^2$: a different number with the same $\ell$-adic
valuation.
\end{exercise}

\begin{exercise}\label{ex:ant:ex-cm}
Let $E\colon y^2=x^3-x$. Show that $a_p=0$ for every prime $p\equiv3\pmod4$, using the substitution $x\mapsto-x$ and $\bigl(\frac{-1}p\bigr)=-1$. Compute $a_5=-2$,
$a_{13}=6$ and $a_{17}=2$ by counting points, and check that $p=(a_p/2)^2+b^2$ with $b\in\Z$ for $p=5,13,17$ and $29$ ($a_{29}=-10$). Relate the first statement to
\Cref{thm:ant:cm}(ii), and the second to the factorisation $p=\varpi\bar\varpi$ in $\Z[i]$ with $a_p=\varpi+\bar\varpi$ for a suitable prime $\varpi$ above $p$ (stated;
\Cref{ch:ellcurves}).
\end{exercise}

\begin{exercise}\label{ex:ant:ex-twist}
(a) For 11a1, $a_2=-2$. Using \Cref{prop:ant:twistdetect} ((a)$\Rightarrow$(b) needs no semisimplicity), show that $\bar\rho_{E,3}\not\simeq\bar\rho_{E,3}\otimes\chi_{-3}$.
(b) For $E\colon y^2=x^3-6x-8$, compute $a_5$, $a_{11}$, $a_{17}$ (answers $3$, $3$, $-6$) and explain, using \Cref{prop:ant:twistdetect} and \Cref{ch:galois}, why $3\mid a_p$
for every prime $p\equiv2\pmod3$ of good reduction.
\end{exercise}

\begin{exercise}\label{ex:ant:ex-level23}
Let $f$ be the newform of level $23$ of \Cref{ex:ant:level23}, with coefficients $a_n=x_n+y_n\theta$, $x_n,y_n\in\Z$, and $f^\sigma$ its conjugate. Show that
$a_n(f)\equiv a_n(f^\sigma)\pmod{\sqrt5}$ for all $n$, and that $(\sqrt5)$ is the only prime of $\mathcal O_{K_f}$ modulo which all coefficients of $f$ and $f^\sigma$
agree. Compute the reduction of $f$ modulo the prime above $2$ up to $q^7$, in $\F_4$.
\end{exercise}

\begin{exercise}\label{ex:ant:ex-enlarge}
Let $K=\Q(\sqrt5)$ and $K'=\Q(\sqrt5,i)$. Determine $e$, $f$, $g$ for $2$, $3$ and $5$ in $K$ and in $K'$ (answers in $K'$: $(2,2,1)$, $(1,2,2)$, $(2,1,2)$), and describe,
for each prime $\lambda$ of $K$ above these, the primes $\lambda'$ of $K'$ above it and the inclusions $k_\lambda\subseteq k_{\lambda'}$ of \Cref{prop:ant:primesabove}(iii).
\end{exercise}

\begin{exercise}\label{ex:ant:ex-unitsquad}
Show that $\Z[\frac{1+\sqrt5}2]^\times=\{\pm u^k:k\in\Z\}$ with $u=\frac{1+\sqrt5}2$, by the method of \Cref{thm:ant:units}(iii). Show that the reduction of $u$
modulo the prime above $2$ has order $3$ in $\F_4^\times$ and modulo each prime above $11$ has order dividing $10$; which orders occur?
\end{exercise}

\section{Notes and references}\label{sec:ant:notes}

The algebraic number theory of \Cref{sec:ant:integers,sec:ant:dedekind,sec:ant:splitting} follows Neukirch \cite{Neukirch}, Chapter I (Sections 2 to 3 for integrality,
integral bases and Dedekind domains, Sections 8 and 9 for the splitting of primes, decomposition and inertia groups, Section 10 for cyclotomic fields), and
Marcus \cite{Marcus77}, Chapters 2 to 4, which is the gentlest complete treatment and has many exercises of the kind given here. The proof of unique
factorisation of ideals (\Cref{lem:ant:primeproducts}, \Cref{thm:ant:idealfactorisation}) is Neukirch's; the proof of transitivity in
\Cref{prop:ant:galoissplitting} and of the surjectivity onto the residue Galois group in \Cref{thm:ant:residuegalois} are the standard ones (Marcus, Chapter 4).
Atiyah and Macdonald \cite{AtiyahMacdonald}, Chapters 5 and 9, treat integral extensions, lying over and Dedekind domains in the language of commutative
algebra. Lang \cite{LangANT} is a further reference for the Dedekind zeta function and densities (Chapter VIII). Minkowski's theorem, Dirichlet's unit theorem,
the analytic continuation of $\zeta_K$ and Chebotarev's theorem are stated, with references to Neukirch, Chapter I, Sections 5 to 7, and Chapter VII, Sections 5 and
13.

The splitting of primes in quadratic fields and the Gaussian and Eisenstein integers are in Ireland and Rosen \cite{IrelandRosen}, Chapters 1 and 9 and
Chapter 13, and in Cox \cite{Cox}, Sections 5 and 7, where prime discriminants (\Cref{prop:ant:ramonlyat}) are also discussed. Cyclotomic fields, the
discriminant formula, Kronecker-Weber and the dictionary between Dirichlet characters and characters of $G_\Q$ follow Washington \cite{Washington}, Chapters 2, 3
and 14. The tame relation (\Cref{prop:ant:tameconj}) is in Serre \cite{SerreLocalFields}, Chapter IV, Sections 1 and 2, where the higher ramification groups are
defined; \Cref{ch:local} gives the local version. Frobenius elements in $G_\Q$ and their use for $\ell$-adic representations are in Serre \cite{SerreAbelian},
Chapter I. Class field theory in idelic form, with its two sign conventions, is in Neukirch, Chapter VI, and in Milne's notes \cite{MilneCFT}.

For complex multiplication, \Cref{thm:ant:cm} quotes Silverman \cite{SilvermanATAEC}, Chapter II and Appendix A (the thirteen rational $j$-invariants), and
Cox \cite{Cox}, Sections 7 to 14 (class number one and Deuring's reduction criterion); the open image theorem is Serre \cite{Serre72}, and the density zero of
$\{p:a_p(E)=0\}$ for $E$ without complex multiplication is Serre \cite{Serre81}. The argument that a prime ramified in the CM field divides the conductor
(\Cref{thm:ant:cm}(iii)) is the elementary one given here; it can also be read off the conductor of an induced representation.

All theorem, section and chapter numbers of books cited in this chapter (Neukirch, Marcus, Washington, Ireland-Rosen, Cox, Serre, Lang, Silverman) are recalled
and were not checked against printed copies; they should be read as pointers, and no proof in the book depends on them. Every numerical example (class
numbers, discriminants, factorisations of primes, Frobenius classes, the coefficients of the level $23$ newform, the traces of Frobenius of the curves of
\Cref{ex:ant:cmcurves} and \Cref{ex:ant:ex-twist}) was checked with PARI/GP \cite{PARI}, for instance with the following commands.
\begin{lstlisting}[language=GP]
quadclassunit(-20)[1]                     \\ class number 2
K = nfinit(x^3 - 2); idealprimedec(K, 5)  \\ residue degrees 1 and 2
nfdisc(polcyclo(7))                       \\ -16807 = -7^5
mf = mfinit([23,2],0); mfcoefs(mfeigenbasis(mf)[1], 13)
E = ellinit([0,-1,1,-7,10]); [ellap(E,p) | p <- [2,7,5,37]]   \\ 0,0,-3,7
\end{lstlisting}
